% Compile with LuaLaTeX. DejaVu Serif and DejaVu Sans Mono supply the Unicode
% characters retained from the Markdown notes. Keep build products out of Git.
\documentclass[11pt,a4paper,oneside]{report}
\usepackage[margin=25mm]{geometry}
\usepackage{fontspec}
\setmainfont{DejaVu Serif}
\setmonofont{DejaVu Sans Mono}[Scale=0.83]
\usepackage{amsmath,amsthm}
\usepackage{unicode-math}
\setmathfont{Latin Modern Math}
\usepackage{longtable,booktabs,array,calc}
\usepackage{microtype}
\usepackage{xcolor}
\usepackage{url}
\usepackage[colorlinks=true,linkcolor=blue!50!black,urlcolor=blue!60!black]{hyperref}
\usepackage{bookmark}
\usepackage{etoolbox}
\AtBeginEnvironment{longtable}{\small\setlength{\tabcolsep}{4pt}}
\providecommand{\tightlist}{%
  \setlength{\itemsep}{0pt}\setlength{\parskip}{0pt}}
\newtheorem{definition}{Definition}[chapter]
\setcounter{secnumdepth}{2}
\setcounter{tocdepth}{1}
\setlength{\parindent}{0pt}
\setlength{\parskip}{0.45em}
\setlength{\emergencystretch}{3em}
\raggedbottom

\title{Weight systems and lattice polytopes\\[0.5em]
  \large Consolidated literature review}
\author{process-polytopes}
\date{14 October 2026}

\begin{document}
\maketitle
\tableofcontents

\chapter{Basic definitions and conventions}
\label{basic-definitions}

This document collects all seven literature-review notes in
\texttt{docs/learning/}. The following definitions fix the notation needed
to read them. Each review chapter retains its original citations, examples,
provenance, and qualifications about the supporting formal results and native
implementation. The concrete Lean classification is described in
\texttt{proofs/README.md} and \texttt{proofs/STATEMENTS.md}; it does not verify
PALP or certify a completed census. The separate single-weight paper is not a
literature-review chapter.

\section{Lattices, polytopes, and interior}

\begin{definition}[Lattice and dual lattice]
A lattice \(M\) of rank \(d\) is a free abelian group isomorphic to
\(\mathbb{Z}^{d}\). Its real extension is
\(M_{\mathbb{R}}=M\otimes_{\mathbb{Z}}\mathbb{R}\), and its dual lattice is
\(N=\operatorname{Hom}(M,\mathbb{Z})\).
The evaluation pairing extends to
\(\langle\,\cdot,\cdot\,\rangle:N_{\mathbb{R}}\times M_{\mathbb{R}}
\longrightarrow\mathbb{R}\).
Choosing a lattice basis identifies \(M\) with \(\mathbb{Z}^{d}\), but a real
change of basis need not preserve that lattice.
\end{definition}

\begin{definition}[Polytope and lattice polytope]
A convex hull consists of convex combinations:
\[
  \operatorname{conv}\{v_1,\ldots,v_m\}
   =\left\{\sum_{i=1}^{m}\lambda_i v_i:
       \lambda_i\geq0,\ \sum_{i=1}^{m}\lambda_i=1\right\}.
\]
A polytope is the convex hull of a finite set of points. It is a
\emph{lattice polytope in \(M\)} if its vertices lie in \(M\).
It is full dimensional if its affine span is \(M_{\mathbb{R}}\).
Its dimension is the dimension of its affine span, not the number of
coordinates used to present it.
\end{definition}

A face is where a supporting linear functional attains its extreme value
on the polytope. Vertices are zero-dimensional faces; facets of a
\(d\)-dimensional polytope are faces of dimension \(d-1\).

The ambient interior \(\operatorname{int}P\) is computed in the stated real
vector space. The relative interior \(\operatorname{relint}P\) is computed
in the affine span of \(P\). A lower-dimensional polytope has empty ambient
interior, but can have nonempty relative interior. When working in a specified
equation subspace, interior must be taken in that entire subspace; using only
the hull's possibly smaller affine span could hide a dimension defect.

\begin{definition}[IP property and one-point polytope]
An origin-marked, full-dimensional polytope has the \emph{IP property} if
\(0\in\operatorname{int}P\). A \emph{one-point lattice polytope} has exactly
one interior lattice point:
\[
  k(P)=\#\bigl(\operatorname{int}P\cap M\bigr)=1.
\]
These are different predicates. For example, \([-2,2]\subset\mathbb{R}\)
has the IP property in \(\mathbb{Z}\), but has three interior lattice points.
Translate the unique interior lattice point of a one-point polytope to zero
when an origin-fixed convention is needed.
\end{definition}

\section{Polarity and reflexivity}

\begin{definition}[Polar and reflexive polytope]
For a full-dimensional polytope \(P\subset M_{\mathbb{R}}\) with the IP
property, use the sign convention
\[
  P^{*}=\{y\in N_{\mathbb{R}}:
    \langle y,x\rangle\geq -1\text{ for every }x\in P\}.
\]
A lattice polytope \(P\) is \emph{reflexive} if it has the IP property and
\(P^{*}\) is a lattice polytope in the dual lattice \(N\).
Equivalently, every facet of \(P\) is at primitive lattice distance one from
zero. If a facet is \(\langle u,x\rangle\geq-h\), where \(u\in N\) is
primitive and \(h\) is a positive integer, its polar vertex is \(u/h\).
\end{definition}

A reflexive polytope has exactly one interior lattice point, but the converse
is false. The polar of a general IP lattice polytope is rational and need
not be integral. Polarity reverses inclusion and is an involution on the
reflexive class in the fixed dual pair of lattices; this does not give the
same class-preserving argument for arbitrary nonreflexive one-point objects.

\section{Simplex circuits and minimality}

\begin{definition}[IP simplex and positive circuit]
An IP simplex has affinely independent vertices \(v_1,\ldots,v_n\) and
zero in its relative interior. Affine independence means that
\(v_2-v_1,\ldots,v_n-v_1\) are linearly independent.
Its unique normalized barycentric relation is
\[
  q_i>0,\qquad \sum_{i=1}^{n}q_i=1,\qquad
  \sum_{i=1}^{n}q_i v_i=0.
\]
This is a positive circuit: the linear-relation space on these vertices is
one dimensional. Affine independence must not be confused with linear
independence of the vertex vectors, which have the displayed nonzero relation.
\end{definition}

\begin{definition}[Different minimality predicates]
A vertex-minimal IP polytope loses the IP property when any vertex is
deleted. Inclusion-minimality in a class forbids any proper subpolytope in
that class; for example, an \(r\)-minimal reflexive polytope has no proper
reflexive lattice subpolytope. An inclusion-maximal member has no proper
superset within the stated class. The class, lattice, embedding, and interior
convention must be fixed. Vertex-minimal IP seeds and \(r\)-minimal reflexive
objects are not interchangeable.
\end{definition}

\section{Weights and weighted projective spaces}

\begin{definition}[Normalized and primitive weights]
A normalized rational weight system is a positive vector
\(q=(q_1,\ldots,q_n)\in\mathbb{Q}_{>0}^{n}\) with \(\sum_i q_i=1\);
take \(n\geq2\) here unless a source explicitly treats dimension zero.
Its primitive integral representation is
\[
  w=(w_1,\ldots,w_n)\in\mathbb{Z}_{>0}^{n},\qquad
  \gcd(w_1,\ldots,w_n)=1,\qquad
  D=\sum_i w_i,\qquad q_i=w_i/D.
\]
The integer \(D\) is the degree. For a lattice simplex, rational barycentric
coordinates give this integral relation after clearing denominators and
dividing the common gcd.
\end{definition}

The weighted projective space associated with \(w\) is
\[
  \mathbb{P}(w)=\operatorname{Proj}\mathbb{C}[X_1,\ldots,X_n],
  \qquad \deg X_i=w_i.
\]
The weights are well formed when
\(\gcd(w_1,\ldots,\widehat{w_i},\ldots,w_n)=1\) for every \(i\).
Global gcd one alone does not imply this condition.
Conrads uses ``reduced'' for global gcd one and ``normalized'' for
well-formed integral weights; those terms differ from sum-one rational
normalization. Weights do not by themselves specify every ambient-lattice
realization of a simplex or its associated toric variety.

\section{Integer Newton hulls and IP weight systems}

\begin{definition}[Newton hull and IP weight system]
For normalized positive rational \(q\) of length \(n\), define
\[
  A(q)=\{a\in\mathbb{R}_{\geq0}^{n}:q\cdot a=1\},
  \qquad
  \Delta(q)=\operatorname{conv}\bigl(A(q)\cap\mathbb{Z}^{n}\bigr).
\]
The all-ones vector \(\mathbf{1}\) satisfies the equation.
The weight system is IP if \(\mathbf{1}\) is interior to the integer hull
in the entire affine hyperplane \(q\cdot a=1\).
After translation,
\[
  P(q)=\Delta(q)-\mathbf{1}\subset H_q=\ker(q\cdot-),\qquad
  L_q=H_q\cap\mathbb{Z}^{n}.
\]
Equivalently, \(0\) is interior to \(P(q)\) in \(H_q\).
The expected dimension is \(n-1\); the IP test includes full dimension
within this hyperplane.
\end{definition}

Positive simplex weights alone do not imply integer Newton-hull IP.
For example, \((1,1,3)/5\) supplies a positive simplex circuit but fails
the Newton-hull IP condition. Conversely, IP Newton hulls need not be
reflexive in dimension five. Depth-\(c\) shifted hulls use the inequalities
\(x_i\geq-c\); a half-weight reduction can produce depth two rather than
the ordinary depth-one problem.

\section{Combined weights, rank, and containers}

\begin{definition}[Combined weight equations]
For rational nonnegative rows, choose an integral matrix
\(W\in\mathbb{Z}_{\geq0}^{r\times n}\), with every coordinate covered by
a positive coefficient in some row. Put \(\mathbf{D}=W\mathbf{1}\).
The equation data underlying a combined weight system give
\[
  A_W=\{a\in\mathbb{R}_{\geq0}^{n}:Wa=\mathbf{D}\},\qquad
  H_W=\ker W,\qquad L_W=H_W\cap\mathbb{Z}^{n},
\]
\[
  P_W=\operatorname{conv}(A_W\cap\mathbb{Z}^{n})-\mathbf{1}.
\]
Its IP condition is \(0\in\operatorname{int}_{H_W}P_W\).
The support of row \(i\) is \(\{j:W_{ij}>0\}\); zeros pad its relation
outside that block.
Rank-nullity gives \(\dim H_W=n-\operatorname{rank}W\), which is \(n-r\)
only when the rows are independent. Geometric CWS constructions require
appropriate positive simplex supports and a complete relation basis;
nonnegative coordinate coverage alone does not supply those properties.
\end{definition}

A container covers a target when a map with the stated lattice properties
sends that target into the container. Containment is not equality or
equivalence. Real surjectivity can transport interior, but does not imply
integer surjectivity. In particular, projection from a CWS space to a
component weight space requires a lifting hypothesis for component IP.

\section{Lattice equivalence, index, and saturation}

\begin{definition}[Linear and affine lattice equivalence]
Origin-fixed lattice equivalence means \(Q=U(P)\), where \(U\) is a
lattice isomorphism; in lattice coordinates \(U\in\operatorname{GL}(d,\mathbb{Z})\).
Affine lattice equivalence allows \(Q=U(P)+t\), with \(t\in M\).
A real-linear bijection preserving lattice points forwards is weaker:
its inverse must also preserve the lattice to be a lattice isomorphism.
Centering the unique interior lattice point turns affine equivalence of
one-point polytopes into origin-fixed equivalence.
\end{definition}

For full-rank sublattices \(L'\subseteq L\), the index
\([L:L']=\#(L/L')\) is finite. The saturation of \(L'\) in \(L\) is
\[
  \operatorname{sat}_{L}(L')
   =\{x\in L:mx\in L'\text{ for some positive integer }m\}.
\]
A rational equation space has the saturated coordinate lattice
\(H_W\cap\mathbb{Z}^{n}\), but the image of an original lattice in that
space can have a proper finite index. The image lattice and its index
must not be discarded when classifying lattice realizations.

\section{Volume and finiteness conventions}

Lattice-coordinate volume \(\operatorname{vol}_{M}\) assigns volume one
to a fundamental parallelepiped of \(M\). Integer normalized volume is
\(\operatorname{Vol}_{M}(P)=d!\operatorname{vol}_{M}(P)\).
These differ by a factorial. Also distinguish the interior count
\(k(P)\) from the total count \(J(P)=\#(P\cap M)\).
The fixed-\((d,k)\) finiteness results discussed here require \(k\geq1\);
unrestricted origin-IP does not fix \(k\), and its union over positive
interior counts need not be finite. An existential volume or box bound
is not automatically an explicit, computable enumeration bound.

\section{Normal forms and algorithmic claims}

\begin{definition}[Canonical normal form and digest]
For a fixed lattice and equivalence relation, a canonical normal form
\(\operatorname{NF}\) satisfies
\[
  \operatorname{NF}(P)=\operatorname{NF}(Q)
   \quad\Longleftrightarrow\quad P\sim Q.
\]
The integer normal-form matrix, its serialized bytes, and a finite-width
hash are distinct objects. Equal hashes do not prove equal normal forms.
Exact class merging must compare the full canonical values and preserve
the multiplicities of the contributing inputs.
\end{definition}

Soundness means every output belongs to the target class. Completeness
means every target class has a represented member. Termination means the
procedure finishes for each input in its stated domain. Producing one
representative per class is an additional requirement. An existential
reachable-node theorem or noncomputable finite catalogue is not itself
a terminating executable enumeration. Abstract mathematical coverage,
native implementation correctness, and a completed reproducible census
therefore remain separate claims throughout this review.

% BEGIN GENERATED LITERATURE REVIEW
\chapter{IP, reflexivity, polarity, and the pyramid
boundary}\label{ip-reflexivity-and-pyramids--ip-reflexivity-polarity-and-the-pyramid-boundary}

\noindent\textit{Markdown source: }\texttt{docs/\allowbreak{}learning\allowbreak{}/\allowbreak{}ip-\allowbreak{}reflexiv\allowbreak{}ity-\allowbreak{}and-\allowbreak{}pyramids\allowbreak{}.\allowbreak{}md}

This note fixes the distinctions between origin-IP, one-point and
reflexive lattice polytopes, and records the scope of the pyramid
results. The \href{../../proofs/STATEMENTS.md}{formal classification}
targets reflexive polytopes only; the examples here explain why
nonreflexive IP objects still occur as intermediates.

\section{Fix the lattice, interior, and equivalence
first}\label{ip-reflexivity-and-pyramids--fix-the-lattice-interior-and-equivalence-first}

Let M be a rank-d lattice, M\_R its real extension, and N = Hom(M,Z) the
dual lattice. A lattice polytope is the convex hull of finitely many
points of M. Throughout, an IP polytope is full dimensional and contains
the chosen origin in its ambient interior. A one-point polytope has
exactly one interior lattice point. These predicates differ: {[}-2,2{]}
in Z has IP and three interior lattice points. For a one-point polytope,
translate its unique interior lattice point to zero before comparing by
GL(d,Z); without this normalization, the natural unmarked equivalence is
affine unimodular equivalence.

KS {[}K00{]}, §2, printed pp.~5--6, defines the lattice, IP property,
and polar separately. The origin in a lower-dimensional simplex belongs
to its intrinsic interior in its span, not the ambient interior. The
lower-dimensional simplices in the seed decomposition therefore require
a relative-interior predicate. Positive simplex barycentric coordinates
alone do not say that the \emph{integer Newton hull} constructed from a
weight system has IP.

For a positive normalized rational weight vector q of length l, the
exponent-space region is

\texttt{Q =\allowbreak{} \{a ∈ R\textasciicircum{}l :\allowbreak{} a\_\allowbreak{}i ≥ 0,\allowbreak{} Σ q\_\allowbreak{}i a\_\allowbreak{}i =\allowbreak{} 1\}},

with integer Newton hull
\texttt{Δ(\allowbreak{}q)\allowbreak{} =\allowbreak{} conv(\allowbreak{}Q ∩ Z\textasciicircum{}l)}
and distinguished point
\texttt{1 =\allowbreak{} (\allowbreak{}1,\allowbreak{}…,\allowbreak{}1)}.
Interior here is intrinsic to the equation hyperplane and relative to
its affine lattice. After translation a=x+1, the region lies in
\texttt{Σ q\_\allowbreak{}i x\_\allowbreak{}i=\allowbreak{}0},
\texttt{x\_\allowbreak{}i≥−1}. For a full-dimensional Newton hull with 1
in its intrinsic interior, any interior lattice point has positive
integer coordinates. Since each a\_i≥1 and Σq\_i=1 with all
q\_i\textgreater0, the equation forces a=1. Thus IP Newton hulls in this
positive-weight construction are one-point polytopes; this does not make
arbitrary IP polytopes one-point. For CWS, the analogous argument needs
the equations to cover every coordinate by a positive coefficient;
uncovered coordinates must not be silently admitted. Skarke {[}S96{]},
§2, p.~3, gives this exponent-coordinate argument and explicitly
acknowledges lattice ambiguity.

\section{Polar duality: what remains
integral}\label{ip-reflexivity-and-pyramids--polar-duality-what-remains-integral}

Use the sign convention

\texttt{P* =\allowbreak{} \{y ∈ N\_\allowbreak{}R :\allowbreak{} <y,\allowbreak{}x> ≥ −1 for every x ∈ P\}}.

For compact convex full-dimensional P with 0 in its interior, \(P^{*}\)
is compact, convex, and IP, and \(P^{**}\)=P. For arbitrary sets P⊆Q,
\(Q^{*}\)⊆\(P^{*}\). The reverse implication holds when Q is closed,
convex, and contains 0: dualizing again and using P⊆\(P^{**}\) yields
P⊆Q. These are real convex-geometric facts; integrality is additional.

A full-dimensional IP lattice polytope is reflexive precisely when
\(P^{*}\) is also a lattice polytope in the \emph{dual} lattice.
Equivalently, every facet has primitive lattice distance one from 0. If
the facet is \texttt{<u,\allowbreak{}x>≥−h} with u primitive in N and h
a positive integer, its polar vertex is u/h. For h\textgreater1 this is
rational and nonintegral. A rational polar is therefore well defined
without being an object of the same lattice-polytope class.

Batyrev {[}B94{]}, Definition 4.1.1, printed p.~17, defines polarity;
Definitions 4.1.4--4.1.5, p.~18, define primitive integral distance and
reflexive pairs. Theorem 4.1.6, p.~18, states that
\texttt{(\allowbreak{}P,\allowbreak{}M)} reflexive implies
\texttt{(\allowbreak{}P*,\allowbreak{}N)} reflexive. Proposition 4.1.7,
p.~19, reverses face incidence and sends an l-face to a (d−l−1)-face.
Its face equality uses pairing −1. The v1 proof of Theorem 4.1.6 has a
sign inconsistency: it writes outward inequalities ≤1 and then
identifies their normals, rather than their negatives, with vertices of
the ≥−1 polar. Correct the normal orientation when reconstructing this
proof; the theorem and defining convention agree after this correction.

Reflexivity implies a unique interior lattice point. An interior lattice
point pairs strictly above −1 with every integral polar vertex, hence at
least 0. If it were nonzero, polarity and bipolarity would put an
unbounded positive ray in the compact primal polytope. The converse
fails; Section 5 supplies an explicit one-point Newton hull with a
distance-two facet.

\section{Toric interpretation and its
hypotheses}\label{ip-reflexivity-and-pyramids--toric-interpretation-and-its-hypotheses}

Batyrev {[}B94{]}, Theorem 4.1.9, printed p.~19, assumes an
n-dimensional integral polyhedron Δ, its projective toric variety P\_Δ,
and the family of projective \textbf{Δ-regular} hypersurfaces. It
equates the Calabi--Yau property (with the canonical-singularity
convention of Definition 4.1.8) with the associated ample sheaf being
anticanonical, and with Δ having a unique interior lattice point m0 such
that \texttt{(\allowbreak{}Δ−m0,\allowbreak{}M)} is reflexive. This is
not a theorem about every polynomial supported in Δ. The regularity
hypothesis cannot be omitted, and a generic toric Calabi--Yau
hypersurface is not asserted smooth in every dimension.

Proposition 4.1.10 and Corollary 4.1.11, printed p.~20, identify the
anticanonical polyhedra and recover the toric fan from cones over the
faces of \(\Delta^{*}\). The dual pair is the basis of Batyrev's mirror
construction; no mirror-symmetry theorem for arbitrary nonreflexive
one-point polytopes follows. None of these toric objects is represented
in the inspected Lean definitions.

\section{Minimality and the class being
dualized}\label{ip-reflexivity-and-pyramids--minimality-and-the-class-being-dualized}

KS {[}K00{]}, §2.1, printed p.~7, calls an IP polytope vertex-minimal if
removing any vertex destroys IP. Its §2.3, printed p.~15, distinguishes
lattice-point minimality, very minimality, and r-maximal reflexive
polytopes with their r-minimal duals. A vertex-minimal IP polytope need
not be reflexive. An r-minimal reflexive polytope can have a proper IP
subpolytope that is nonreflexive. Neither predicate can replace the
other.

On a fixed dual pair of lattices, reflexive polarity is an involution
that reverses inclusion. Consequently P is inclusion-maximal
\textbf{among reflexive lattice polytopes} iff \(P^{*}\) is
inclusion-minimal among reflexive lattice polytopes in the dual lattice.
Real antitonicity alone proves no such assertion for a class not closed
under polarity. The comparison also fixes embeddings and origins; it is
not an inclusion relation between bare GL-equivalence classes.

KS {[}K00{]}, §2.3, printed pp.~13--15, starts with reflexive targets,
constructs CWS containers on the finest lattice, and requires suitable
sublattices and subpolytopes. Lemma 7 assumes an integral vertex pairing
matrix and describes the intermediate-lattice recovery; a rational
pairing matrix does not meet its hypothesis. A nonreflexive IP container
can still contain reflexive targets. This is different from a claim that
all nonreflexive one-point targets are covered by the same completeness
argument: their rational polars need not admit integral minimal seeds in
the fixed dual lattice. Clearing denominators changes lattice data and
primitive distances and requires a new correctness/completeness
argument. Rational polarity does not supply the missing argument,
lattice enumeration, termination, or output guarantees.

\section{Skarke's theorem and two distinct
counterexamples}\label{ip-reflexivity-and-pyramids--skarkes-theorem-and-two-distinct-counterexamples}

\subsection{The actual low-dimensional
theorem}\label{ip-reflexivity-and-pyramids--the-actual-low-dimensional-theorem}

Skarke {[}S96{]}, §3, \textbf{Lemma 1}, printed pp.~6--7, assumes an
integer pyramid of lattice height h≥2 in a rank-four lattice, and its
double-height pyramid with the same apex and shape. The double contains
a lattice point outside the original pyramid and outside the doubled
base. The remarks on p.~7 extend the proof to dimensions below four and
give a failure in dimension five.

The \textbf{unnumbered Theorem} on p.~7 states that maximal Newton
polyhedra of dimension at most four with 1 in their interior are
reflexive. ``Maximal Newton'' means the entire integer hull of the
prescribed nonnegative exponent region; it does not mean
inclusion-maximal among all reflexive polytopes. The proof uses this
full integer hull: the doubling lies in the region with exponent
coordinates ≥−1, and a lattice point away from the doubled base has
nonnegative coordinates and hence belongs to the original Newton hull.
The following remark includes constructions defined by several weight
equations. KS {[}K00{]}, Lemma 8, printed p.~15, restates the conclusion
as \texttt{n≤4} and \texttt{Δ(\allowbreak{}q)} IP implies reflexive.
Neither statement is ``every IP lattice polytope of dimension ≤4 is
reflexive.''

There is a normalization obligation in the source-to-Lean bridge. Skarke
triangulates the base and writes its vertices as 0,2e1,2e2,2e3. An
arbitrary lattice tetrahedron is not necessarily unimodular: a real
change making its edges coordinate vectors need not preserve the ambient
lattice or keep the apex coordinates integral. The frozen theorem checks
only the displayed normalized family. A proof of the general reduction
must track the transformed lattice or supply additional lattice
geometry; the comment asserting a general GL(4,Z) reduction is not
itself that proof. This note does not conclude that Skarke's stated
lemma is false; it identifies an obligation not discharged by the
inspected theorem type.

\subsection{The five-dimensional height
example}\label{ip-reflexivity-and-pyramids--the-five-dimensional-height-example}

The doubled pyramid has base
\texttt{conv(\allowbreak{}0,\allowbreak{}2e1,\allowbreak{}2e2,\allowbreak{}2e3,\allowbreak{}2e4)}
at height zero and apex
\texttt{(\allowbreak{}2,\allowbreak{}2,\allowbreak{}2,\allowbreak{}2,\allowbreak{}4)}.
Its doubled height is 4; the associated original height parameter is
h=2. At height 1 its first four coordinates are each at least 1/2. A
lattice point must therefore have each at least 1. The cap inequality

\texttt{(\allowbreak{}x1+x2+x3+x4)\allowbreak{}/\allowbreak{}2 − 3x5/\allowbreak{}4 ≤ 1}

would then give 5/4≤1, a contradiction. This is exactly a failure of the
intermediate-height lemma. The pyramid has 0 on its boundary and is not
an IP weight system. The failure of one sufficient proof mechanism does
not logically imply that a nonreflexive IP weight system exists.

\subsection{An actual nonreflexive IP weight system: an independent
derivation}\label{ip-reflexivity-and-pyramids--an-actual-nonreflexive-ip-weight-system-an-independent-derivation}

Take weights
\texttt{(\allowbreak{}1,\allowbreak{}1,\allowbreak{}1,\allowbreak{}1,\allowbreak{}1,\allowbreak{}2)}
and total degree 7. Write t=a6. Its complete integer Newton hull is

\texttt{H =\allowbreak{} \{a\_\allowbreak{}i≥0 :\allowbreak{} a1+…+a5+2t=\allowbreak{}7,\allowbreak{} 0≤t≤3\}}.

To verify equality, all integer exponent solutions lie in H. Conversely
H is the convex hull of the two endpoint simplices: their vertices are
\texttt{7e\_\allowbreak{}i} at t=0 and \texttt{e\_\allowbreak{}i+3e6} at
t=3, for i=1,\ldots,5. These ten vertices are integer solutions, so the
complete hull is H. Its dimension is five, and 1 has all coordinates
positive and 0\textless t\textless3, hence belongs to its intrinsic
interior. Every interior integer point has all six coordinates ≥1; the
degree equation then forces every coordinate to equal 1. Thus H is
one-point.

The face t=3 is a four-dimensional facet. On the translated rank-five
lattice

\texttt{L =\allowbreak{} \{x∈Z\textasciicircum{}6 :\allowbreak{} x1+…+x5+2x6=\allowbreak{}0\}},

the coordinate functional x6 is primitive: \texttt{e6−2e1∈L} has value
1. That facet has lattice distance \texttt{3−1=\allowbreak{}2} from the
interior point. The corresponding polar vertex is
\texttt{−x6/\allowbreak{}2} as a functional on L and is nonintegral
because it takes value −1/2 on \texttt{e6−2e1}. Hence this IP Newton
hull is not reflexive.

This is an explicit original exact-arithmetic derivation, independently
inspectable from the displayed vertices and primitive lattice vector. It
is not attributed to Skarke, is not a checked Lean theorem, and is not a
claim about census counts. It separately establishes the existence which
the height example alone cannot establish.

\section{Formal support: actual types, not
comments}\label{ip-reflexivity-and-pyramids--formal-support-actual-types-not-comments}

Paths are relative to \texttt{proofs/\allowbreak{}Refpoly/};
\href{../../proofs/SOURCES.md}{SOURCES.md} identifies the original
source versions. The table records the supporting declarations'
mathematical scope.

\begin{longtable}[]{@{}
  >{\raggedright\arraybackslash}p{(\linewidth - 4\tabcolsep) * \real{0.2751}}
  >{\raggedright\arraybackslash}p{(\linewidth - 4\tabcolsep) * \real{0.3916}}
  >{\raggedright\arraybackslash}p{(\linewidth - 4\tabcolsep) * \real{0.3333}}@{}}
\toprule\noalign{}
\begin{minipage}[b]{\linewidth}\raggedright
File / declaration
\end{minipage} & \begin{minipage}[b]{\linewidth}\raggedright
Actual scope
\end{minipage} & \begin{minipage}[b]{\linewidth}\raggedright
Boundary
\end{minipage} \\
\midrule\noalign{}
\endhead
\bottomrule\noalign{}
\endlastfoot
\texttt{Basic.\allowbreak{}lean}: \texttt{V},
\texttt{IsLattic\allowbreak{}ePoint},
\texttt{IsLattic\allowbreak{}ePolytop\allowbreak{}e} &
\texttt{V n =\allowbreak{} Fin n → R}; coordinate integrality; finite
lattice generating hull & A chosen Z\^{}n model, not an arbitrary
embedded lattice with a supplied basis \\
\texttt{Basic.\allowbreak{}lean}: \texttt{IsIP} &
\texttt{0 ∈ interior\allowbreak{} P} & Ambient interior; no one-point
count \\
\texttt{Basic.\allowbreak{}lean}: \texttt{IsReflex\allowbreak{}ive} &
Lattice polytope, IP, and polar lattice polytope & Definition uses dual
integrality; primitive facet-distance equivalence not supplied by this
declaration \\
\texttt{Basic.\allowbreak{}lean}:
\texttt{polarDua\allowbreak{}l\_\allowbreak{}antitone} &
\texttt{P⊆Q → polarDua\allowbreak{}l Q⊆polarDua\allowbreak{}l P} for
arbitrary sets & No lattice premise or class closure \\
\texttt{Basic.\allowbreak{}lean}: \texttt{bipolar} & Convex P, closed P,
\texttt{0∈P} imply \(P^{**}\)=P & Origin need not be interior; no
polytope premise required \\
\texttt{Basic.\allowbreak{}lean}:
\texttt{interior\allowbreak{}\_\allowbreak{}latticeP\allowbreak{}oint\_\allowbreak{}unique}
&
\texttt{IsReflex\allowbreak{}ive P → interior\allowbreak{} P ∩ latticeP\allowbreak{}oints n =\allowbreak{} \{0\}}
& No converse \\
\texttt{Basic.\allowbreak{}lean}: \texttt{IsUnimod\allowbreak{}ular},
\texttt{GLEquiv} & Real linear equivalence preserving integers forwards
\textbf{and backwards}; exact image equality & Forward integrality alone
is insufficient; no translation parameter \\
\texttt{Reflexiv\allowbreak{}eDuality\allowbreak{}.\allowbreak{}lean}:
\texttt{IsReflex\allowbreak{}ive.\allowbreak{}bipolar\_\allowbreak{}eq},
\texttt{IsReflex\allowbreak{}ive.\allowbreak{}polarDua\allowbreak{}l\_\allowbreak{}reflexiv\allowbreak{}e}
& Reflexive P has \(P^{**}\)=P and reflexive \(P^{*}\) & Assumes the
primal reflexivity structure; does not prove all IP objects reflexive \\
\texttt{MaxMin.\allowbreak{}lean}:
\texttt{subset\_\allowbreak{}iff\_\allowbreak{}polarDua\allowbreak{}l\_\allowbreak{}subset}
& Closed convex Q containing 0 gives \texttt{P⊆Q ↔ Q*⊆P*} & A real order
statement \\
\texttt{MaxMin.\allowbreak{}lean}:
\texttt{IsMaxima\allowbreak{}lReflexi\allowbreak{}ve},
\texttt{IsMinima\allowbreak{}lReflexi\allowbreak{}ve},
\texttt{maximal\_\allowbreak{}iff\_\allowbreak{}dual\_\allowbreak{}minimal}
& Quantifies over reflexive sets on the same Z\^{}n model; equivalence
assumes reflexive P & Does not apply to all one-point polytopes \\
\texttt{MaxMin.\allowbreak{}lean}: \texttt{IsMinima\allowbreak{}lIPGen}
& A finite generating set whose hull is ambient IP and every proper
subset fails IP & Distinct from r-minimality; no lattice condition in
the predicate \\
\texttt{Minimal.\allowbreak{}lean}: \texttt{IsGoodSi\allowbreak{}mplex},
\texttt{IsMinima\allowbreak{}lIPRel} & Affine independence and intrinsic
interior; relative minimality tested in the fixed parent span & A
correct distinction from ambient IP \\
\texttt{PyramidL\allowbreak{}emma.\allowbreak{}lean}:
\texttt{pyramid\_\allowbreak{}lemma\_\allowbreak{}dim\_\allowbreak{}le\_\allowbreak{}four}
& For integers x,y,z,h with h≥2, produces an integral p in \textbf{V 4},
in
\texttt{conv(\allowbreak{}0,\allowbreak{}2e1,\allowbreak{}2e2,\allowbreak{}2e3,\allowbreak{}(\allowbreak{}2x,\allowbreak{}2y,\allowbreak{}2z,\allowbreak{}2h)\allowbreak{})},
with \texttt{1≤p4≤h−1} & Not dimension-parametric and no
arbitrar\allowbreak{}y-\allowbreak{}base/\allowbreak{}lattice
normalization theorem in this type \\
\texttt{PyramidL\allowbreak{}emma.\allowbreak{}lean}:
\texttt{pyramid\_\allowbreak{}countere\allowbreak{}xample\_\allowbreak{}dim\_\allowbreak{}five}
& Integral m in \texttt{conv(\allowbreak{}pyramid5\allowbreak{})}
implies \texttt{m5≠1} & Height obstruction only; no WS, IP, or
nonreflexivity conclusion \\
\end{longtable}

The normalized pyramid existence and counterexample have proof bodies in
the snapshot. The actual low-dimensional integer Newton-hull reflexivity
theorem is deliberately not declared in
\texttt{PyramidL\allowbreak{}emma.\allowbreak{}lean}. Its facet-distance
and relative-lattice infrastructure remains a separate obligation.
Existing comments saying that the height example makes automatic WS
reflexivity fail are broader than the counterexample's type.

\section{Computational consequences and
gaps}\label{ip-reflexivity-and-pyramids--computational-consequences-and-gaps}

\begin{longtable}[]{@{}
  >{\raggedright\arraybackslash}p{(\linewidth - 4\tabcolsep) * \real{0.2395}}
  >{\raggedright\arraybackslash}p{(\linewidth - 4\tabcolsep) * \real{0.2966}}
  >{\raggedright\arraybackslash}p{(\linewidth - 4\tabcolsep) * \real{0.4639}}@{}}
\toprule\noalign{}
\begin{minipage}[b]{\linewidth}\raggedright
Obligation
\end{minipage} & \begin{minipage}[b]{\linewidth}\raggedright
Evidence
\end{minipage} & \begin{minipage}[b]{\linewidth}\raggedright
Remaining boundary
\end{minipage} \\
\midrule\noalign{}
\endhead
\bottomrule\noalign{}
\endlastfoot
Separate IP from reflexivity in dimension five & Explicit WS derivation
above; low-dimensional source theorem has n≤4 & Native tests must use
the actual lattice and separate accepted
IP/\allowbreak{}reflexiv\allowbreak{}e counts \\
Recover all reflexive targets from CWS containers & {[}K00{]} §2.3;
concrete \texttt{KS} reconstruction and composition & The Lean endpoint
includes lattice realizations; the native descriptor
catalogu\allowbreak{}e/\allowbreak{}referenc\allowbreak{}e lattice is a
separate guarantee \\
Cover all nonreflexive one-point targets & Real polarity exists but may
be nonintegral & Need a separate theorem or rational-dual extension
preserving lattice data \\
Transfer Skarke's pyramid theorem to the normalized Lean family &
{[}S96{]} §3 and inspected exact type & General base
normaliz\allowbreak{}ation/\allowbreak{}lattice tracking and
lower-dimensional specialization need proof \\
Automatic reflexivity of n≤4 IP Newton hulls & {[}S96{]} unnumbered
theorem; {[}K00{]} Lemma 8 & Not present as a theorem in the inspected
pyramid file \\
Toric interpretation & {[}B94{]} Theorem 4.1.9 with Δ-regularity & No
toric formalization in these declarations \\
Minimal/\allowbreak{}maximal equivalence & Frozen
\texttt{maximal\_\allowbreak{}iff\_\allowbreak{}dual\_\allowbreak{}minimal}
& Class and lattice fixed; vertex-minimal seeds are a different
predicate \\
Formal trust boundary & Exact source versions, reviewed interfaces and
compiler-derived axiom closures & See \href{../proof-audit.md}{proof
guarantees}; primary-source prose is not a substitute for the formal
statement \\
\end{longtable}

PALP's accepted CWS rows and exact normal forms are consistency
regression evidence. They do not establish native refinement of the
generic Lean coverage and lattice-realization theorems, or a completed
census.

\section{Bibliography and cached source
evidence}\label{ip-reflexivity-and-pyramids--bibliography-and-cached-source-evidence}

All page references above are \textbf{printed preprint page numbers},
not journal pagination. Cached bytes are local ignored artifacts; the
citations record consulted editions and original-file SHA-256 values.

\begin{itemize}
\tightlist
\item
  \textbf{{[}B94{]}} Victor V. Batyrev, \emph{Dual Polyhedra and Mirror
  Symmetry for Calabi-Yau Hypersurfaces in Toric Varieties}, Journal of
  Algebraic Geometry \textbf{3} (1994), 493--535.
  \href{https://arxiv.org/abs/alg-geom/9310003}{Primary arXiv record},
  \href{https://arxiv.org/pdf/alg-geom/9310003v1}{v1 PDF},
  \href{https://arxiv.org/src/alg-geom/9310003v1}{v1 author source}.
  Version v1, 5 October 1993; arXiv DOI
  \texttt{10.\allowbreak{}48550/\allowbreak{}arXiv.\allowbreak{}alg-\allowbreak{}geom/\allowbreak{}9310003}
  (not asserted as a journal DOI). Cache
  \texttt{sources/\allowbreak{}batyrev-\allowbreak{}1994/\allowbreak{}v1/\allowbreak{}paper.\allowbreak{}pdf},
  SHA-256
  \texttt{ebcd9bc3\allowbreak{}2414d1a9\allowbreak{}f4b52f02\allowbreak{}56d78144\allowbreak{}e371ee7b\allowbreak{}edc47738\allowbreak{}7dcad017\allowbreak{}077cad10};
  original gzip author source \texttt{source.\allowbreak{}download},
  SHA-256
  \texttt{3e724a0c\allowbreak{}2b93b3fb\allowbreak{}69c2b9cc\allowbreak{}944b0b18\allowbreak{}7b5ab558\allowbreak{}bd721281\allowbreak{}b4776fc8\allowbreak{}d6fe2e38}.
  \texttt{tex/\allowbreak{}paper.\allowbreak{}tex} is supplementary gzip
  extraction. Retrieved 2026-10-14; provenance in the same directory.
  The license grant beyond arXiv distribution was not independently
  verified.
\item
  \textbf{{[}S96{]}} Harald Skarke, \emph{Weight systems for toric
  Calabi--Yau varieties and reflexivity of Newton polyhedra}, Modern
  Physics Letters A \textbf{11} (1996), 1637--1652, DOI
  \href{https://doi.org/10.1142/S0217732396001636}{10.\allowbreak{}1142/\allowbreak{}S0217732\allowbreak{}39600163\allowbreak{}6}.
  \href{https://arxiv.org/abs/alg-geom/9603007v1}{Primary v1}, 7 March
  1996. Cache
  \texttt{sources/\allowbreak{}skarke-\allowbreak{}alg-\allowbreak{}geom-\allowbreak{}9603007/\allowbreak{}v1/\allowbreak{}paper.\allowbreak{}pdf},
  SHA-256
  \texttt{ee4357a5\allowbreak{}eadd242e\allowbreak{}9689546a\allowbreak{}7991a9e5\allowbreak{}c291bb7b\allowbreak{}4627674e\allowbreak{}59865a7e\allowbreak{}80698800};
  retrieval 2026-10-14. §2 pp.~3--4 for exponent lattice and hull IP; §3
  pp.~6--7 for Lemma 1 and the unnumbered reflexivity theorem.
\item
  \textbf{{[}K00{]}} Maximilian Kreuzer and Harald Skarke,
  \emph{Reflexive polyhedra, weights and toric Calabi--Yau fibrations},
  Reviews in Mathematical Physics \textbf{14(4)} (April 2002), 343--374,
  DOI
  \href{https://doi.org/10.1142/S0129055X0200120X}{10.\allowbreak{}1142/\allowbreak{}S0129055\allowbreak{}X0200120\allowbreak{}X}.
  Journal volume/pages agree with the primary arXiv metadata;
  DOI/issue/date agree with the publisher-deposited Crossref record.
  \href{https://arxiv.org/abs/math/0001106}{Primary arXiv record},
  \href{https://arxiv.org/pdf/math/0001106v1}{v1 PDF}, v1, 19 January
  2000. Cache
  \texttt{sources/\allowbreak{}kreuzer-\allowbreak{}skarke-\allowbreak{}math-\allowbreak{}0001106/\allowbreak{}v1/\allowbreak{}paper.\allowbreak{}pdf},
  SHA-256
  \texttt{7d400220\allowbreak{}d9bb4f0e\allowbreak{}0e2a121d\allowbreak{}46d75a7a\allowbreak{}ad0b0a98\allowbreak{}ebbc1d6a\allowbreak{}d45e3340\allowbreak{}7d0b01c1};
  retrieval 2026-10-14. §2 pp.~5--6 definitions; §2.1 p.~7 minimality;
  §2.3 pp.~13--15 lattice recovery and refined minimality. The preprint
  is the consulted edition; the journal DOI does not identify a
  separately consulted edition.
\end{itemize}

The original weight-system example is an exact derivation, separate from
the cited source statements and the Lean theorem inventory.

\chapter{Weight systems, simplex circuits, and weighted projective
spaces}\label{weight-systems-and-wps--weight-systems-simplex-circuits-and-weighted-projective-spaces}

\noindent\textit{Markdown source: }\texttt{docs/\allowbreak{}learning\allowbreak{}/\allowbreak{}weight-\allowbreak{}systems-\allowbreak{}and-\allowbreak{}wps.\allowbreak{}md}

This note distinguishes barycentric weights, integer Newton-hull IP and
weighted projective spaces (WPS). Equal weights need not determine the
same lattice polytope. The supporting formal declarations below retain
their precise hypotheses; the complete executable reflexive
classification is described in \href{../../proofs/STATEMENTS.md}{the
statement guide}.

\section{Three normalizations that must remain
distinct}\label{weight-systems-and-wps--three-normalizations-that-must-remain-distinct}

For an IP d-simplex with lattice vertices v0,\ldots,vd, the unique
barycentric coordinates of 0 satisfy

\texttt{q\_\allowbreak{}i>0,\allowbreak{} Σq\_\allowbreak{}i=\allowbreak{}1,\allowbreak{} Σq\_\allowbreak{}i v\_\allowbreak{}i=\allowbreak{}0}.

They are rational because they solve an invertible rational affine
system. Clear denominators and divide the common gcd to obtain a
\textbf{primitive integral relation}
\texttt{Σw\_\allowbreak{}i v\_\allowbreak{}i=\allowbreak{}0} with
positive integers w\_i and gcd(w\_i)=1. Conversely q\_i=w\_i/D, D=Σw\_i.
Thus normalized rational weights mean sum one in KS notation, whereas
primitive integral weights mean global gcd one. Neither phrase by itself
means well formed.

Conrads {[}C02{]}, Definition 2.2, pp.~217--218, uses different
terminology:

\begin{longtable}[]{@{}
  >{\raggedright\arraybackslash}p{(\linewidth - 4\tabcolsep) * \real{0.3274}}
  >{\raggedright\arraybackslash}p{(\linewidth - 4\tabcolsep) * \real{0.2389}}
  >{\raggedright\arraybackslash}p{(\linewidth - 4\tabcolsep) * \real{0.4336}}@{}}
\toprule\noalign{}
\begin{minipage}[b]{\linewidth}\raggedright
Condition
\end{minipage} & \begin{minipage}[b]{\linewidth}\raggedright
Meaning here
\end{minipage} & \begin{minipage}[b]{\linewidth}\raggedright
Conrads' terminology
\end{minipage} \\
\midrule\noalign{}
\endhead
\bottomrule\noalign{}
\endlastfoot
Σq\_i=1 for positive rational q & Barycentric normalization & Not his
definition of normalized integral weights \\
gcd(w0,\ldots,wd)=1 & Primitive integral relation & Reduced \\
gcd(w0,\ldots,omit w\_i,\ldots,wd)=1 for every i & Well-formed WPS
weights & Normalized \\
\end{longtable}

Example: (2,3,4) is globally primitive, but omitting 3 leaves
gcd(2,4)=2. Conrads' Example 3.9, p.~223, replaces it by the well-formed
tuple (1,3,2), with the same WPS fan. Dividing a \emph{common} gcd
changes only scale; well-forming a tuple may change individual entries
and therefore the barycentric relation of a particular chosen simplex.
The coarse WPS can remain isomorphic even when the original integral
simplex changes. One must preserve the original relation and lattice
data when using weights as polytope invariants.

Conrads' Definition 2.3, p.~218, assigns the \textbf{determinant weight}
\texttt{Q\_\allowbreak{}i=\allowbreak{}|det(\allowbreak{}v0,\allowbreak{}…,\allowbreak{}omit v\_\allowbreak{}i,\allowbreak{}…,\allowbreak{}vd)\allowbreak{}|}.
Its factor λ=gcd(Q\_i) is the index of the vertex-generated lattice in
the ambient lattice (Lemma 2.4). The primitive relation is
\texttt{Q\_\allowbreak{}red=\allowbreak{}Q/\allowbreak{}λ}, and the
simplex type is
\texttt{(\allowbreak{}Q\_\allowbreak{}red,\allowbreak{}λ)}. This Q is
not automatically the primitive relation. The barycentric coordinates
are
\texttt{Q\_\allowbreak{}i/\allowbreak{}ΣQ\_\allowbreak{}j=\allowbreak{}Q\_\allowbreak{}red,\allowbreak{}i/\allowbreak{}ΣQ\_\allowbreak{}red,\allowbreak{}j}.

For an affinely independent family of m lattice vectors in a larger
ambient space with a positive relation, the same rational positive
circuit exists in the family span; intrinsic interior replaces ambient
interior. The circuit specifies one relation. A CWS requires an
independent family of such relations plus support coverage, not
arbitrary repeated rows.

\section{The WPS and the lattice
realization}\label{weight-systems-and-wps--the-wps-and-the-lattice-realization}

Conrads {[}C02{]}, §3, p.~218, defines
\texttt{P(\allowbreak{}w)\allowbreak{}=\allowbreak{}Proj C[\allowbreak{}X0,\allowbreak{}…,\allowbreak{}Xd]},
graded by deg Xi=w\_i. Theorem 3.6, p.~221, constructs a fan with d+1
rays and an integral simplex of type
\texttt{(\allowbreak{}Q\_\allowbreak{}red,\allowbreak{}1)} spanning that
fan. Corollary 4.6, p.~224, makes this factor-one simplex unique up to
unimodular equivalence. Corollary 3.7, p.~222, says its vertices are
primitive precisely when the reduced integral weight tuple is well
formed; passing to primitive ray generators produces the well-formed
tuple and preserves the fan. Corollary 3.8 identifies WPS isomorphism
classes by the well-formed weights, with entry permutations understood.

A complete simplicial fan with d+1 \textbf{primitive} rays whose
positive relation has weights w defines the WPS P(w) when the rays
generate the ambient lattice. Buczyńska {[}Bu08{]}, Notation 5.1 and
Proposition 5.4, §5, explicitly requires those primitive rays and
lattice generation. A larger ambient lattice can carry the same rational
fan and same positive weights while producing a finite quotient of that
WPS. This is a \textbf{fake WPS} in her Definition 6.1. Lemma 6.3 and
Theorem 6.4, §6, realize the quotient via the inclusion of the
vertex/ray-generated rank-d lattice N' into N; the WPS on N' is the
unique universal cover in codimension one. The primitive-ray hypothesis
ensures no branching in codimension one. This conclusion is stronger
than a generic finite toric quotient and must not be asserted when a
nonprimitive vertex is being used as a ray generator.

Conrads' Proposition 4.7, p.~224, also shows that the toric variety
associated with any integral simplex is a finite-group quotient of P(Q),
with quotient group abstractly isomorphic to
\texttt{Z\textasciicircum{}d/\allowbreak{}N\_\allowbreak{}Δ}. The
distinction is that his simplex definition permits nonprimitive
vertices: its generated sublattice is not always the lattice generated
by \textbf{primitive ray generators}. That proposition alone does not
assert a cover unramified in codimension one. When the vertices are
primitive, the two lattice constructions coincide. For a reflexive
lattice simplex the vertices are primitive: a nonprimitive multiple kv,
k\textgreater1, cannot pair to ±1 with an integral polar vertex on a
facet through kv.

For a well-formed primitive relation w, one can express the WPS lattice
as the free quotient
\texttt{Z\textasciicircum{}(\allowbreak{}d+1)\allowbreak{}/\allowbreak{}Z w}.
Rational circuit data recover the real relation space, but an integral
simplex can embed its vertex-generated lattice with nontrivial finite
index in another ambient lattice. Real-linear equivalence and a forward
integral map are consequently weaker than a lattice isomorphism. The
primitive weights alone do not enumerate all reflexive simplex lattice
realizations, let alone all CWS realizations.

\section{What Conrads proves about reflexive
simplices}\label{weight-systems-and-wps--what-conrads-proves-about-reflexive-simplices}

Conrads uses the polar convention \texttt{<y,\allowbreak{}x>≤1} in
equation (2.1), p.~216. The frozen package and KS use
\texttt{<y,\allowbreak{}x>≥−1}. The polars differ by negation, so dual
integrality and reflexivity agree; signs of displayed facet normals must
be translated.

Theorem 4.4, pp.~223--224, states that every integral d-simplex of type
\texttt{(\allowbreak{}Q\_\allowbreak{}red,\allowbreak{}λ)} is
unimodularly equivalent to \texttt{H Δ\_\allowbreak{}Q}, for H in the
finite set of determinant-λ Hermite normal forms. Definition 4.1 and
Theorem 4.2 fix that matrix set. Remark 4.5 warns that different H may
give isomorphic simplices: generating these matrices is coverage, not
canonical deduplication.

Proposition 5.1, pp.~224--225, has two directions with different
hypotheses. If an integral simplex is reflexive, each entry of its
reduced weight divides the total reduced weight. Conversely, when an
integral weight tuple satisfies this divisibility, the unique
\textbf{factor-one} simplex associated to its reduced tuple is
reflexive. For a well-formed tuple w, this gives the familiar WPS
Gorenstein criterion \texttt{w\_\allowbreak{}i | D} for every i,
D=Σw\_i. The relation argument is explicit: an integral polar vertex for
the facet opposite v\_i pairs with v\_i by
\texttt{1−D/\allowbreak{}w\_\allowbreak{}i} in Conrads' sign convention,
so D/w\_i must be integral. For factor one this divisibility is
sufficient; for a larger ambient lattice, the same rational polar may
cease to be integral.

There is a source hypothesis qualification that should be retained.
Corollary 5.2, p.~225, prints the WPS criterion without restating
well-formedness. Read as a criterion for \emph{arbitrary} presentations
of the same WPS it fails: Example 3.9 identifies P(2,3,4) with P(1,3,2).
The second has total 6 divisible by each entry, but the first has total
9 not divisible by 2 or 4. Apply the criterion to the well-formed
presentation. Similarly, the proof of Proposition 5.1 says the divisible
tuple is ``necessarily normalized''; that needs prior gcd reduction (the
scaled tuple (2,2,2) satisfies divisibility but is not well formed as
written). The simplex construction is scale invariant, so the factor-one
conclusion survives after reduction. These qualifications come from
checking the source's own definitions and examples, not from silently
changing its terminology.

Lemma 5.3, p.~225, shows that dual reflexive simplices have the same
reduced weights, possibly different factors. For the factor-one simplex
Δ\_Q let m\_Q be the factor of its dual (Definition 5.4). Proposition
5.5 proves that the factor λ of any reflexive simplex with these reduced
weights divides m\_Q. This is a necessary restriction. The algorithm on
p.~226 still tests whether
\texttt{(\allowbreak{}H\textasciicircum{}T)\allowbreak{}\textasciicircum{}(\allowbreak{}-\allowbreak{}1)\allowbreak{} Δ\_\allowbreak{}Q*}
is integral, then compares isomorphism classes. Divisibility of λ alone
does not assert that every determinant-λ H gives a reflexive simplex.
The classification algorithm is for simplices and lattice realizations;
it does not prove WS/CWS subpolytope coverage of all reflexive polytopes
or native-pipeline correctness.

\section{Simplex IP and integer Newton-hull IP are different
tests}\label{weight-systems-and-wps--simplex-ip-and-integer-newton-hull-ip-are-different-tests}

For positive rational q with sum one, write

\texttt{Q\_\allowbreak{}exp(\allowbreak{}q)\allowbreak{}=\allowbreak{}\{a\_\allowbreak{}i≥0 :\allowbreak{} Σq\_\allowbreak{}i a\_\allowbreak{}i=\allowbreak{}1\}},
\texttt{Δ\_\allowbreak{}Newton(\allowbreak{}q)\allowbreak{}=\allowbreak{}conv(\allowbreak{}Q\_\allowbreak{}exp(\allowbreak{}q)\allowbreak{}∩Z\textasciicircum{}(\allowbreak{}d+1)\allowbreak{})}.

The all-ones vector belongs to the exponent equation. The \textbf{IPWS
condition} asks whether it lies in the intrinsic interior of the full
integer Newton hull. It is stronger than positivity of q in an abstract
simplex circuit. For well-formed weights with all w\_i dividing D, the
exponent simplex has integer vertices
\texttt{(\allowbreak{}D/\allowbreak{}w\_\allowbreak{}i)\allowbreak{}e\_\allowbreak{}i};
its integer hull is the entire exponent simplex, and 1 is in its
intrinsic interior. When divisibility fails, the integer hull may still
have IP, but it requires an independent test.

An original exact example separates the predicates: take weights
(1,1,3), D=5. The lattice triangle with vertices
\texttt{(\allowbreak{}-\allowbreak{}1,\allowbreak{}-\allowbreak{}3)\allowbreak{},\allowbreak{}(\allowbreak{}1,\allowbreak{}0)\allowbreak{},\allowbreak{}(\allowbreak{}0,\allowbreak{}1)}
has positive relation with those weights, so 0 is interior. But integer
exponent solutions of \texttt{a0+a1+3a2=\allowbreak{}5}, ai≥0 have a2≤1;
1 has a2=1 on a supporting boundary, so the Newton hull is not IP.
Positive normalized weights can therefore describe an IP simplex while
failing IPWS.

Conversely, an IPWS need not yield a reflexive Newton hull in dimension
five: the exact example
\texttt{(\allowbreak{}1,\allowbreak{}1,\allowbreak{}1,\allowbreak{}1,\allowbreak{}1,\allowbreak{}2)\allowbreak{}/\allowbreak{}7}
and its primitive distance-two facet are derived in
\hyperref[ip-reflexivity-and-pyramids--an-actual-nonreflexive-ip-weight-system-an-independent-derivation]{the
polarity\allowbreak{}/\allowbreak{}pyramid note}. It is a one-point
integer Newton hull, without a Gorenstein WPS weight tuple. Skarke's
low-dimensional theorem concerns the integer hull in dimension at most
four, not all simplex circuits or arbitrary ambient-lattice
realizations.

The converse connection relevant to KS requires a geometric witness. If
a positive simplex circuit is formed from \textbf{dual lattice points of
a reflexive polytope P}, the pairing map sends integral generators of P
into the depth-one weight lattice, and surjectivity over the reals sends
an open neighborhood of 0 onto an open neighborhood in the weight
hyperplane. Its image lies in the Newton hull, proving IPWS. No claim
that the pairing map is onto the \textbf{integer} weight lattice is
needed for that interior argument.

\section{Two examples about lattice and projection
hypotheses}\label{weight-systems-and-wps--two-examples-about-lattice-and-projection-hypotheses}

\subsection{Same weights, distinct lattice
polytopes}\label{weight-systems-and-wps--same-weights-distinct-lattice-polytopes}

The triangle
\texttt{T=\allowbreak{}conv(\allowbreak{}(\allowbreak{}1,\allowbreak{}0)\allowbreak{},\allowbreak{}(\allowbreak{}0,\allowbreak{}1)\allowbreak{},\allowbreak{}(\allowbreak{}-\allowbreak{}1,\allowbreak{}-\allowbreak{}1)\allowbreak{})}
has reduced weights (1,1,1), factor one, and area 3/2. Apply the integer
matrix
\texttt{H=\allowbreak{}(\allowbreak{}(\allowbreak{}1,\allowbreak{}1)\allowbreak{},\allowbreak{}(\allowbreak{}0,\allowbreak{}3)\allowbreak{})},
of determinant three. The result has vertices
\texttt{(\allowbreak{}1,\allowbreak{}0)\allowbreak{},\allowbreak{}(\allowbreak{}1,\allowbreak{}3)\allowbreak{},\allowbreak{}(\allowbreak{}-\allowbreak{}2,\allowbreak{}-\allowbreak{}3)},
reduced weights still (1,1,1), factor three, and area 9/2. All three
vertices are primitive. Its facet inequalities are

\texttt{x≤1}, \texttt{x−y≤1}, \texttt{−2x+y≤1}.

Their normals are integral, so the triangle is reflexive (negate normals
for the package convention). Its fan in the generated index-three
sublattice gives P²; in the ambient Z² it gives a fake WPS finite
quotient of P². The two triangles are not unimodularly equivalent
because their lattice areas differ. This is an independently derived
example; no census or formal-certificate claim is made. It explains why
simplex weights alone cannot distinguish even reflexive lattice
polytopes.

\subsection{An abstract CWS can lose block projection
surjectivity}\label{weight-systems-and-wps--an-abstract-cws-can-lose-block-projection-surjectivity}

Take two rows \texttt{(\allowbreak{}1,\allowbreak{}1)} and
\texttt{(\allowbreak{}1,\allowbreak{}2)} in R². They are nonnegative,
linearly independent, and cover both coordinates, hence meet the frozen
abstract CWS structure. Their joint kernel is \{0\}; the hull of its
allowed lattice points is \{0\}, with 0 interior in this
zero-dimensional subspace, so abstract CWS IP holds. The second row's
kernel is the line \texttt{y1+2y2=\allowbreak{}0}; its depth-one integer
generators are \texttt{(\allowbreak{}0,\allowbreak{}0)} and
\texttt{(\allowbreak{}2,\allowbreak{}−1)}. Their hull has 0 as an
endpoint, so that row fails IPWS. Projection from the joint zero space
is not surjective onto the row's line. This elementary example explains
the explicit lifting hypothesis in
\texttt{cws\_\allowbreak{}ip\_\allowbreak{}componen\allowbreak{}t}. The
geometric simplex-block theorem below proves its own surjectivity under
stronger geometric hypotheses.

\section{Supporting formal declarations and exact
boundaries}\label{weight-systems-and-wps--supporting-formal-declarations-and-exact-boundaries}

Paths are relative to \texttt{proofs/\allowbreak{}Refpoly/}.
\href{../../proofs/SOURCES.md}{SOURCES.md} identifies the original
source versions. The table describes supporting-library statements, not
the stronger concrete \texttt{KS} endpoint.

\begin{longtable}[]{@{}
  >{\raggedright\arraybackslash}p{(\linewidth - 4\tabcolsep) * \real{0.2065}}
  >{\raggedright\arraybackslash}p{(\linewidth - 4\tabcolsep) * \real{0.4783}}
  >{\raggedright\arraybackslash}p{(\linewidth - 4\tabcolsep) * \real{0.3152}}@{}}
\toprule\noalign{}
\begin{minipage}[b]{\linewidth}\raggedright
Declaration
\end{minipage} & \begin{minipage}[b]{\linewidth}\raggedright
Exact hypothes\allowbreak{}is/\allowbreak{}conclusi\allowbreak{}on
support
\end{minipage} & \begin{minipage}[b]{\linewidth}\raggedright
Boundary
\end{minipage} \\
\midrule\noalign{}
\endhead
\bottomrule\noalign{}
\endlastfoot
\texttt{Minimal.\allowbreak{}lean:\allowbreak{} Refpoly.\allowbreak{}simplex\_\allowbreak{}weights}
& Affine independence of finite S and
\texttt{0∈interior\allowbreak{}(\allowbreak{}conv S)} produce v indexed
by \texttt{Fin(\allowbreak{}n+1)} and real q, with exact range S, affine
independence, positive q, sum one, relation zero & No lattice premise,
rationality, gcd normalization, uniqueness clause, or WPS object in this
theorem type \\
\texttt{WeightSy\allowbreak{}stem.\allowbreak{}lean:\allowbreak{} Refpoly.\allowbreak{}ip\_\allowbreak{}simplex\_\allowbreak{}iff\_\allowbreak{}weights\_\allowbreak{}pos}
& An affine basis b, sum-one real w, and zero affine combination give
interior iff all w positive & Full ambient simplex; does not assert
Newton-hull IP \\
\texttt{WeightSy\allowbreak{}stem.\allowbreak{}lean:\allowbreak{} Refpoly.\allowbreak{}wsHyperp\allowbreak{}lane}
& Kernel of the real dot functional q & Coordinate kernel alone does not
identify all lattice realizations \\
\texttt{WeightSy\allowbreak{}stem.\allowbreak{}lean:\allowbreak{} Refpoly.\allowbreak{}wsLattic\allowbreak{}eGensAt},
\texttt{Refpoly.\allowbreak{}wsPolyto\allowbreak{}peAt} &
Coordinate-integral points in the kernel with every coordinate ≥−c, then
their convex hull \textbf{inside} that subspace & Uses the inherited
coordinate lattice and subspace topology \\
\texttt{WeightSy\allowbreak{}stem.\allowbreak{}lean:\allowbreak{} Refpoly.\allowbreak{}IsIPWeig\allowbreak{}htSystem\allowbreak{}At},
\texttt{Refpoly.\allowbreak{}IsIPWeig\allowbreak{}htSystem} & Positive q
and origin interior to that hull; depth one in the standard predicate &
No rationality or sum-one premise in the definition \\
\texttt{WeightSy\allowbreak{}stem.\allowbreak{}lean:\allowbreak{} Refpoly.\allowbreak{}wsPoints},
\texttt{Refpoly.\allowbreak{}wsPoints\allowbreak{}\_\allowbreak{}shift}
& Nonnegative integral x with \texttt{x·q=\allowbreak{}Σq}; exact
correspondence after subtracting ones to depth-one generators & Works
without sum-one normalization; q is real \\
\texttt{WeightSy\allowbreak{}stem.\allowbreak{}lean:\allowbreak{} Refpoly.\allowbreak{}not\_\allowbreak{}ip\_\allowbreak{}of\_\allowbreak{}weight\_\allowbreak{}gt\_\allowbreak{}half}
& Length n≥2 and \texttt{Σq<2q\_\allowbreak{}j} imply not IPWS &
Necessary restriction, no sufficiency \\
\texttt{WeightSy\allowbreak{}stem.\allowbreak{}lean:\allowbreak{} Refpoly.\allowbreak{}CWS},
\texttt{Refpoly.\allowbreak{}CWS.\allowbreak{}subspace},
\texttt{Refpoly.\allowbreak{}CWS.\allowbreak{}IsIP} & Real nonnegative
independent rows, coordinate coverage; joint kernel; IP hull in joint
subspace & No rationality, row sum normalization, or automatic
block-lifting property \\
\texttt{WeightSy\allowbreak{}stem.\allowbreak{}lean:\allowbreak{} Refpoly.\allowbreak{}cws\_\allowbreak{}ip\_\allowbreak{}componen\allowbreak{}t}
& CWS IP, injective support listing e, matching row weights, and
\textbf{hlift} for every vector in the row kernel imply row IPWS & hlift
is surjectivity over R; abstract CWS conditions alone do not discharge
it \\
\texttt{Geometri\allowbreak{}cCWS.\allowbreak{}lean:\allowbreak{} Refpoly.\allowbreak{}simplex\_\allowbreak{}relation\allowbreak{}\_\allowbreak{}kernel}
& Affine-independent v, sum-one q, zero relation imply coefficient-map
kernel is \texttt{R∙q} & No positivity or lattice premise; real relation
space \\
\texttt{Geometri\allowbreak{}cCWS.\allowbreak{}lean:\allowbreak{} Refpoly.\allowbreak{}simplex\_\allowbreak{}pairing\_\allowbreak{}surjecti\allowbreak{}ve}
& Same affine-independence, normalization and relation hypotheses imply
every z with q·z=0 lifts to real x with pairings v\_i·x=z\_i & Includes
lower-dimensional circuits; no integer-lift or lattice-isomorphism
conclusion \\
\texttt{Geometri\allowbreak{}cCWS.\allowbreak{}lean:\allowbreak{} Refpoly.\allowbreak{}reflexiv\allowbreak{}e\_\allowbreak{}simplex\_\allowbreak{}block\_\allowbreak{}ip}
& Reflexive P, affine-independent positive normalized circuit v,q, every
v\_i integral and in \(P^{*}\), imply IPWS q & Crucial extra dual
geometry; surjectivity proved internally, no separate hlift premise \\
\texttt{Geometri\allowbreak{}cCWS.\allowbreak{}lean:\allowbreak{} Refpoly.\allowbreak{}algorith\allowbreak{}m\_\allowbreak{}reaches\_\allowbreak{}simplex\_\allowbreak{}block}
& Same geometry and a \texttt{Selectio\allowbreak{}nRule m} produce
finite S, \texttt{Reachabl\allowbreak{}e sel S},
\texttt{sel.\allowbreak{}sel S=\allowbreak{}q}, and
\texttt{S.\allowbreak{}card≤m+1} & Abstract reachability for the
selection rule; no native executable output list, deduplication, or WPS
classification \\
\end{longtable}

The pairing-surjectivity proof compares real dimensions using the single
relation kernel and matrix rank under transpose. The block-IP proof
restricts the pairing map to the weight hyperplane, uses it as an open
surjective real-linear map, and uses integral pairings only in the
forward direction. Thus it avoids confusing real surjectivity with an
integral isomorphism.

\section{Obligations and computational
implications}\label{weight-systems-and-wps--obligations-and-computational-implications}

\begin{longtable}[]{@{}
  >{\raggedright\arraybackslash}p{(\linewidth - 4\tabcolsep) * \real{0.2222}}
  >{\raggedright\arraybackslash}p{(\linewidth - 4\tabcolsep) * \real{0.3142}}
  >{\raggedright\arraybackslash}p{(\linewidth - 4\tabcolsep) * \real{0.4636}}@{}}
\toprule\noalign{}
\begin{minipage}[b]{\linewidth}\raggedright
Required connection
\end{minipage} & \begin{minipage}[b]{\linewidth}\raggedright
Evidence available
\end{minipage} & \begin{minipage}[b]{\linewidth}\raggedright
Gap or consequence
\end{minipage} \\
\midrule\noalign{}
\endhead
\bottomrule\noalign{}
\endlastfoot
Real simplex weights → rational primitive integral weights & Lattice
vertices and rational affine-system argument; Conrads determinant
weights & \texttt{simplex\_\allowbreak{}weights} alone is real;
\texttt{KS/\allowbreak{}Rational\allowbreak{}Weights} and
\texttt{KS/\allowbreak{}CircuitW\allowbreak{}eights} supply the
rational\allowbreak{}/\allowbreak{}primitiv\allowbreak{}e bridge \\
Primitive relation → well-formed WPS presentation & {[}C02{]}
Definitions 2.2--2.3, Corollary 3.7 & Preserve simplex data when passing
to primitive fan rays; gcd one alone is insufficient \\
WPS weights → all reflexive simplex realizations & {[}C02{]} Theorem
4.4, Proposition 5.5 and dual-integrality test &
\texttt{KS/\allowbreak{}Realizat\allowbreak{}ions} reconstructs
compatible lattices; WPS-specific toric interpretation and exact normal
forms are separate \\
Positive circuit → Newton-hull IP & Fails for
(\allowbreak{}1,\allowbreak{}1,\allowbreak{}3)\allowbreak{}/\allowbreak{}5;
geometric block theorem supplies sufficient witness & Native weight
acceptance must test the integer hull, not positivity alone \\
IPWS → reflexive Newton hull & Low-dimensional Skarke theorem; explicit
5D failure & Separate IP/\allowbreak{}reflexiv\allowbreak{}e predicates
and counts in dimension five \\
CWS IP → component IP & Frozen
\texttt{cws\_\allowbreak{}ip\_\allowbreak{}componen\allowbreak{}t} with
hlift; geometric circuits supply a different bridge & Arbitrary abstract
rows do not guarantee projection surjectivity \\
Abstract reachability → executable WS/\allowbreak{}CWS search &
\texttt{KS/\allowbreak{}WeightSe\allowbreak{}arch} and
\texttt{KS/\allowbreak{}Assembly} prove terminating
sound/\allowbreak{}complete enumeration & Native input pools,
descriptors and pruning need separate refinement; this is not a
production census \\
Source's WPS Gorenstein criterion & Divisibility on a well-formed
presentation & Conrads Corollary 5.2 needs that presentation
qualification; do not copy its shorthand to arbitrary reduced tuples \\
\end{longtable}

Equality of weights is not equality of lattice polytopes. The concrete
Lean classifier includes compatible lattice realizations and subpolytope
search, while native refinement and exact normal-form deduplication
remain separate guarantees.

\section{Verified bibliography and
provenance}\label{weight-systems-and-wps--verified-bibliography-and-provenance}

\begin{itemize}
\tightlist
\item
  \textbf{{[}C02{]}} Heinke Conrads, \emph{Weighted projective spaces
  and reflexive simplices}, \textbf{manuscripta mathematica 107(2)}
  (February 2002), \textbf{215--227}, DOI
  \href{https://doi.org/10.1007/s002290100235}{10.\allowbreak{}1007/\allowbreak{}s0022901\allowbreak{}00235}.
  Author, title, journal, volume, pages and DOI verified on the
  \href{https://link.springer.com/article/10.1007/s002290100235}{publisher
  record}; full 13-page
  \href{https://link.springer.com/content/pdf/10.1007/s002290100235.pdf}{publisher
  PDF} consulted. Received 30 March 1999; revised 18 October 2001. Cache
  \texttt{sources/\allowbreak{}conrads-\allowbreak{}2002/\allowbreak{}publishe\allowbreak{}r/\allowbreak{}paper.\allowbreak{}pdf},
  SHA-256
  \texttt{5108867a\allowbreak{}34b4f449\allowbreak{}44e1f2e6\allowbreak{}5a34c642\allowbreak{}8a5132f8\allowbreak{}9124c8ac\allowbreak{}b949f7fe\allowbreak{}054a7f1e};
  retrieval 2026-10-14. Page citations are journal pages. Copyright
  Springer-Verlag 2002; no author TeX source archive is identified here.
\item
  \textbf{{[}Bu08{]}} Weronika Buczyńska, \emph{Fake weighted projective
  spaces},
  \href{https://arxiv.org/abs/0805.1211v1}{arXiv:\allowbreak{}0805.\allowbreak{}1211v1},
  submitted 8 May 2008; title page dates the work June 2002 and the
  primary record identifies it as a Warsaw University master's thesis
  defended then. DOI
  \href{https://doi.org/10.48550/arXiv.0805.1211}{10.\allowbreak{}48550/\allowbreak{}arXiv.\allowbreak{}0805.\allowbreak{}1211}
  is the arXiv DOI; no journal publication is asserted. Read author TeX
  from \href{https://arxiv.org/src/0805.1211v1}{v1 source} and checked
  statement numbering against
  \href{https://arxiv.org/html/0805.1211v1}{primary HTML}: §5
  Proposition 5.4; §6 Definition 6.1, Lemma 6.3, Theorem 6.4. Cache
  \texttt{sources/\allowbreak{}buczynsk\allowbreak{}a-\allowbreak{}0805-\allowbreak{}1211/\allowbreak{}v1/\allowbreak{}paper.\allowbreak{}pdf},
  SHA-256
  \texttt{c6f0e514\allowbreak{}c078f861\allowbreak{}86dcb4f6\allowbreak{}6b97f237\allowbreak{}904fcb8d\allowbreak{}689261db\allowbreak{}3825408f\allowbreak{}7977a284};
  original gzip tar \texttt{source.\allowbreak{}download}, SHA-256
  \texttt{b50f6f0b\allowbreak{}cfcf7936\allowbreak{}7e9eabe0\allowbreak{}4a7e337c\allowbreak{}717ad4e6\allowbreak{}ec107239\allowbreak{}8aa0d337\allowbreak{}d9c3f3f2};
  extraction
  \texttt{tex/\allowbreak{}fakeWeig\allowbreak{}htedProj\allowbreak{}ectiveSp\allowbreak{}aces.\allowbreak{}tex}
  and \texttt{.\allowbreak{}bbl}. Retrieved 2026-10-14; perpetual
  non-exclusive arXiv distribution license, no broader grant inferred.
  PDF text extraction had glyph-encoding problems; author TeX supplied
  the actual readable arguments.
\item
  \textbf{{[}S96{]}} Harald Skarke, \emph{Weight systems for toric
  Calabi--Yau varieties and reflexivity of Newton polyhedra}, Modern
  Physics Letters A \textbf{11} (1996), 1637--1652, DOI
  \href{https://doi.org/10.1142/S0217732396001636}{10.\allowbreak{}1142/\allowbreak{}S0217732\allowbreak{}39600163\allowbreak{}6};
  \href{https://arxiv.org/abs/alg-geom/9603007v1}{primary arXiv v1}, 7
  March 1996. Shared cache
  \texttt{sources/\allowbreak{}skarke-\allowbreak{}alg-\allowbreak{}geom-\allowbreak{}9603007/\allowbreak{}v1/\allowbreak{}paper.\allowbreak{}pdf},
  SHA-256
  \texttt{ee4357a5\allowbreak{}eadd242e\allowbreak{}9689546a\allowbreak{}7991a9e5\allowbreak{}c291bb7b\allowbreak{}4627674e\allowbreak{}59865a7e\allowbreak{}80698800}.
  §2, printed pp.~3--4, distinguishes barycentric weights from integer
  Newton-hull IP; §3, p.~7, gives the low-dimensional theorem.
  Provenance records retrieval 2026-10-14.
\item
  \textbf{{[}K00{]}} Maximilian Kreuzer and Harald Skarke,
  \emph{Reflexive polyhedra, weights and toric Calabi--Yau fibrations},
  Reviews in Mathematical Physics \textbf{14(4)} (April 2002), 343--374,
  DOI
  \href{https://doi.org/10.1142/S0129055X0200120X}{10.\allowbreak{}1142/\allowbreak{}S0129055\allowbreak{}X0200120\allowbreak{}X}.
  DOI/issue/date agree with the publisher-deposited Crossref record;
  journal volume/pages agree with the
  \href{https://arxiv.org/abs/math/0001106}{primary arXiv record}.
  Consulted version \href{https://arxiv.org/pdf/math/0001106v1}{v1}, 19
  January 2000, cache
  \texttt{sources/\allowbreak{}kreuzer-\allowbreak{}skarke-\allowbreak{}math-\allowbreak{}0001106/\allowbreak{}v1/\allowbreak{}paper.\allowbreak{}pdf},
  SHA-256
  \texttt{7d400220\allowbreak{}d9bb4f0e\allowbreak{}0e2a121d\allowbreak{}46d75a7a\allowbreak{}ad0b0a98\allowbreak{}ebbc1d6a\allowbreak{}d45e3340\allowbreak{}7d0b01c1}.
  §2.2, printed pp.~9--12, supplies the weights/container construction;
  §2.3, pp.~13--15, supplies lattice realization. The preprint is the
  consulted edition; the journal DOI does not identify a separately
  consulted edition. Retrieval 2026-10-14.
\end{itemize}

Original examples are derivations rather than formal certificates. The
source versions, theorem hypotheses and original-file hashes above
distinguish primary mathematical evidence from the formal statements.

\chapter{Minimal IP seeds, positive circuits, and finite structure
data}\label{minimal-ip-circuits-and-types--minimal-ip-seeds-positive-circuits-and-finite-structure-data}

\noindent\textit{Markdown source: }\texttt{docs/\allowbreak{}learning\allowbreak{}/\allowbreak{}minimal-\allowbreak{}ip-\allowbreak{}circuits\allowbreak{}-\allowbreak{}and-\allowbreak{}types.\allowbreak{}md}

This note explains the minimal-seed and circuit geometry behind the
dimension-generic Kreuzer--Skarke construction, and distinguishes it
from archived five-dimensional structure lists. The
\href{../../proofs/STATEMENTS.md}{formal endpoint} proves reflexive
coverage with lattice realizations and subpolytope recovery; equality
with a native descriptor catalogue is a separate claim.

\section{Primary references and evidence
scope}\label{minimal-ip-circuits-and-types--primary-references-and-evidence-scope}

\textbf{{[}KS97{]}} Maximilian Kreuzer and Harald Skarke, \emph{On the
Classification of Reflexive Polyhedra}, Communications in Mathematical
Physics \textbf{185} (1997), 495--508,
\href{https://doi.org/10.1007/s002200050100}{journal DOI}. The inspected
primary preprint is
\href{https://arxiv.org/abs/hep-th/9512204v2}{hep-\allowbreak{}th/\allowbreak{}9512204v\allowbreak{}2},
submitted 4 February 1997; its cover records December 1995 and a
November 1996 revision. References below use the printed pages of the
13-page \href{https://arxiv.org/pdf/hep-th/9512204v2}{v2 PDF}, not
guessed journal pagination. The geometric statements are in §2,
especially printed page 4; the example continues on page 5. The
small-dimensional lists are in §4, pages 9--10.

\textbf{{[}KS00{]}} Maximilian Kreuzer and Harald Skarke,
\emph{Reflexive polyhedra, weights and toric Calabi-Yau fibrations},
Reviews in Mathematical Physics \textbf{14} (2002), 343--374. The
inspected primary version is
\href{https://arxiv.org/abs/math/0001106v1}{math/\allowbreak{}0001106v\allowbreak{}1},
submitted 19 January 2000,
\href{https://arxiv.org/pdf/math/0001106v1}{36-page PDF}. Section 2.1,
pages 7--9, reformulates the structural argument using IP simplex
structures; §2.2, page 10, states the CWS pairing-space and lattice
result as Lemma 4. Publication metadata was checked against the arXiv
records. The publication year and preprint year intentionally differ.

The consulted primary PDF editions have the following original-file
hashes. Retrieval date is 2026-10-14 UTC; cache paths are relative to
\texttt{docs/\allowbreak{}learning\allowbreak{}/\allowbreak{}sources/}:

\begin{longtable}[]{@{}
  >{\raggedright\arraybackslash}p{(\linewidth - 4\tabcolsep) * \real{0.1667}}
  >{\raggedright\arraybackslash}p{(\linewidth - 4\tabcolsep) * \real{0.3333}}
  >{\raggedright\arraybackslash}p{(\linewidth - 4\tabcolsep) * \real{0.5000}}@{}}
\toprule\noalign{}
\begin{minipage}[b]{\linewidth}\raggedright
Source/\allowbreak{}version
\end{minipage} & \begin{minipage}[b]{\linewidth}\raggedright
Original cached PDF
\end{minipage} & \begin{minipage}[b]{\linewidth}\raggedright
SHA-256
\end{minipage} \\
\midrule\noalign{}
\endhead
\bottomrule\noalign{}
\endlastfoot
KS97, hep-\allowbreak{}th/\allowbreak{}9512204v\allowbreak{}2 &
\texttt{kreuzer-\allowbreak{}skarke-\allowbreak{}hep-\allowbreak{}th-\allowbreak{}9512204/\allowbreak{}v2/\allowbreak{}paper.\allowbreak{}pdf}
&
\texttt{c703beef\allowbreak{}d69d2f20\allowbreak{}3c32cfaf\allowbreak{}cf46dc97\allowbreak{}1b7721b8\allowbreak{}ee23b831\allowbreak{}738123cd\allowbreak{}cf896075} \\
KS00, math/\allowbreak{}0001106v\allowbreak{}1 &
\texttt{kreuzer-\allowbreak{}skarke-\allowbreak{}math-\allowbreak{}0001106/\allowbreak{}v1/\allowbreak{}paper.\allowbreak{}pdf}
&
\texttt{7d400220\allowbreak{}d9bb4f0e\allowbreak{}0e2a121d\allowbreak{}46d75a7a\allowbreak{}ad0b0a98\allowbreak{}ebbc1d6a\allowbreak{}d45e3340\allowbreak{}7d0b01c1} \\
\end{longtable}

Each directory has the primary URL and retrieval timestamp in
\texttt{provenan\allowbreak{}ce.\allowbreak{}json};
\texttt{paper.\allowbreak{}txt} is a derived reading aid. Source
licenses are unverified, so this is a local reading cache rather than a
redistribution claim.
\texttt{SINGLE\_\allowbreak{}WS\_\allowbreak{}5d.\allowbreak{}tex} is a
separate single-weight paper, not the specification or proof-status
authority for this construction.

Formal declaration paths below are relative to
\texttt{proofs/\allowbreak{}Refpoly/};
\href{../../proofs/SOURCES.md}{SOURCES.md} identifies the original
editions. Archived catalogue evidence is fixed to root-repository ref
\texttt{gpu-\allowbreak{}run-\allowbreak{}43types}, commit
\texttt{452fbe8e\allowbreak{}1b296bf7\allowbreak{}5fa9e8be\allowbreak{}cf6c34d8\allowbreak{}d48bf00a},
at
\texttt{lean/\allowbreak{}MinimalP\allowbreak{}olytopes\allowbreak{}5D/\allowbreak{}Profiles\allowbreak{}.\allowbreak{}lean},
\texttt{lean/\allowbreak{}MinimalP\allowbreak{}olytopes\allowbreak{}5D/\allowbreak{}Structur\allowbreak{}es.\allowbreak{}lean},
\texttt{paper/\allowbreak{}minimal\_\allowbreak{}polytope\allowbreak{}s\_\allowbreak{}5d.\allowbreak{}tex},
and
\texttt{paper/\allowbreak{}minimal\_\allowbreak{}polytope\allowbreak{}s\_\allowbreak{}5d\_\allowbreak{}table.\allowbreak{}tex}.

\section{Definitions that control the
interpretation}\label{minimal-ip-circuits-and-types--definitions-that-control-the-interpretation}

Fix a real vector space of dimension \texttt{n}, a finite generating set
\texttt{S}, and origin \texttt{0}. The IP property means that \texttt{0}
lies in the interior of \texttt{conv(\allowbreak{}S)} in the stated
ambient space. It does \textbf{not} assert that there is exactly one
interior lattice point. A lower-dimensional component has empty ambient
interior, so its appropriate property is interior relative to its affine
span. If the hull contains zero, that affine span is the corresponding
linear span, which makes quotient arguments legitimate.

A vertex-minimal full-dimensional IP set retains IP itself and loses
full-dimensional IP after removing any generator. For a finite set,
testing all proper subsets and testing each one-point deletion are
equivalent by monotonicity of convex hull and interior. This is distinct
from inclusion-minimal reflexivity, minimality among all lattice
subpolytopes, or a globally smallest vertex count across unrelated
polytopes. No reflexivity or lattice hypothesis is needed for the real
geometric decomposition.

A \textbf{good simplex}, or \textbf{IP simplex}, has affinely
independent vertices and origin in its intrinsic interior. Equivalently,
its unique normalized barycentric relation has strictly positive
coefficients:

\texttt{q\_\allowbreak{}i > 0},
\texttt{sum q\_\allowbreak{}i =\allowbreak{} 1},
\texttt{sum q\_\allowbreak{}i v\_\allowbreak{}i =\allowbreak{} 0}.

The vertex vectors of such a nontrivial simplex are linearly dependent:
the positive relation is exactly that dependence. Affine independence
instead says that any relation whose coefficients also sum to zero is
trivial. Subtracting suitable multiples of the normalized positive
relation shows that its linear-relation space is one-dimensional. In the
homogenized coordinates
\texttt{(\allowbreak{}v\_\allowbreak{}i,\allowbreak{}1)}, the same
vertices are linearly independent. One must specify the representation
before calling simplex vertices ``independent.''

The positive relation guarantees simplex IP in its own span. An IP
\textbf{weight system} in the enumeration sense requires a separate
integer Newton-hull condition. Merely exhibiting a positive circuit does
not establish that condition. Real weights also do not automatically
give primitive integral weights: for lattice/rational vertices,
rationality and normalization require their own bridge.

\section{The structural reduction and what follows from
it}\label{minimal-ip-circuits-and-types--the-structural-reduction-and-what-follows-from-it}

\subsection{A single
step}\label{minimal-ip-circuits-and-types--a-single-step}

{[}KS97{]} §2, Lemma 1, page 4, concerns an \texttt{n}-dimensional
minimal polytope with \texttt{k} vertices. It either is a simplex, or
has a lower-dimensional minimal subpolytope \texttt{M'} with dimension
\texttt{n'} and \texttt{k'} vertices, together with a good simplex
\texttt{T}. The simplex uses all vertices outside \texttt{M'} and
possibly some old vertices; their union is the original vertex set. The
quantitative conclusions are

\texttt{k -\allowbreak{} k' =\allowbreak{} n -\allowbreak{} n' + 1} and
\texttt{dim(\allowbreak{}T)\allowbreak{} <=\allowbreak{} n'}.

The quotient by the span of \texttt{M'} is crucial: the new vertices
form a simplex there. This is stronger than finding an arbitrary proper
subset whose hull contains zero. It is also stronger than covering each
vertex by some positive circuit. Maximality of the proper relative-IP
component, positive lifting from the quotient, and minimality supply the
dimension and counting assertions.

{[}KS00{]} §2.1, Lemma 1, pages 7--8, gives the same reduction in real
spaces. Thus the geometry used here is not restricted to a particular
integral lattice realization.

\subsection{Iteration, vertices, and
supports}\label{minimal-ip-circuits-and-types--iteration-vertices-and-supports}

{[}KS97{]} Corollary 1, page 4, iterates Lemma 1 into an ordered
sequence of good simplices and partial minimal polytopes
\texttt{M\_\allowbreak{}mu}. If \texttt{k\_\allowbreak{}mu} counts the
vertices present through stage \texttt{mu}, then

\texttt{dim(\allowbreak{}M\_\allowbreak{}mu)\allowbreak{} =\allowbreak{} k\_\allowbreak{}mu -\allowbreak{} mu}.

For a final \texttt{n}-dimensional object with \texttt{k} vertices, the
number of stages is therefore
\texttt{lambda =\allowbreak{} k -\allowbreak{} n}. Each stage adds
\texttt{m\_\allowbreak{}mu >=\allowbreak{} 2} vertices, contributes
\texttt{m\_\allowbreak{}mu -\allowbreak{} 1} new dimensions, and the
dimension increments sum to \texttt{n}. {[}KS00{]} §2.1, page 8,
explicitly defines IP simplex structures with partial minimality and at
least two newly introduced vertices, and proves existence in its
unnumbered corollary.

{[}KS97{]} Corollary 2, page 4, gives
\texttt{n + 1 <=\allowbreak{} k <=\allowbreak{} 2n} for full-dimensional
vertex-minimal IP polytopes in positive dimension. Hence
\texttt{1 <=\allowbreak{} lambda <=\allowbreak{} n}. Every good-simplex
support has at most \texttt{n + 1} vertices; for a non-simplex step the
stronger bound
\texttt{dim(\allowbreak{}T)\allowbreak{} <=\allowbreak{} n'} constrains
its overlap with the old component.

In dimension five, the resulting sorted \textbf{new-vertex} profiles are

\texttt{[\allowbreak{}6]}, \texttt{[\allowbreak{}5,\allowbreak{}2]},
\texttt{[\allowbreak{}4,\allowbreak{}3]},
\texttt{[\allowbreak{}4,\allowbreak{}2,\allowbreak{}2]},
\texttt{[\allowbreak{}3,\allowbreak{}3,\allowbreak{}2]},
\texttt{[\allowbreak{}3,\allowbreak{}2,\allowbreak{}2,\allowbreak{}2]},
\texttt{[\allowbreak{}2,\allowbreak{}2,\allowbreak{}2,\allowbreak{}2,\allowbreak{}2]}.

These entries count vertices introduced at each stage. They are not
simplex support sizes: an attached simplex may additionally use old
vertices. Sorting the increments is an arithmetic operation, not a proof
that every reordered support family is still a KS decomposition.

The upper vertex bound gives a finite abstract structural search space
in every fixed dimension: finitely many supports among at most
\texttt{2n} labeled vertices, finitely many finite support families, and
finitely many relevant orderings. This observation is a finiteness route
for combinatorial candidates. It does not imply finitely many arbitrary
real shapes, finitely many unrestricted IP lattice polytopes, effective
weight bounds, or a geometrically sound list returned by a particular
program.

\subsection{No-lonely restriction and decomposition
ambiguity}\label{minimal-ip-circuits-and-types--no-lonely-restriction-and-decomposition-ambiguity}

{[}KS97{]} Lemma 2, page 4, excludes a simplex support with
\textbf{exactly one} vertex occurring in none of the other supports of a
spanning good-simplex family. {[}KS00{]} Lemma 2, page 8, states the
condition for an IP simplex structure. Zero exclusive vertices are
allowed. The five-dimensional example on {[}KS97{]} pages 4--5 and
{[}KS00{]} pages 8--9 shows a simplex whose vertices are all covered by
other blocks and shows nonunique choices of decomposition.

This is a necessary geometric restriction, not an if-and-only-if
characterization of arbitrary support hypergraphs. Passing the no-lonely
filter does not prove affine independence, positive realizability, full
span, partial minimality, independence of equations, or
equivalence-class correctness. A simplex entirely covered by the
\textbf{union} of other blocks is possible. A good-simplex support
properly contained in \textbf{one} other good-simplex support is
impossible: the smaller support gives barycentric coordinates with zeros
on the extra vertices, whereas intrinsic interior in the larger affinely
independent simplex requires every coordinate strictly positive.
Uniqueness of barycentric coordinates gives a contradiction.

\section{Supporting seed and circuit
declarations}\label{minimal-ip-circuits-and-types--supporting-seed-and-circuit-declarations}

The namespace is \texttt{Refpoly}, and \texttt{V n} is the real
coordinate space used by the package. The table describes
supporting-library statements; the concrete \texttt{KS} modules supply
primitive encodings, bounded seeds and executable enumeration.

\begin{longtable}[]{@{}
  >{\raggedright\arraybackslash}p{(\linewidth - 4\tabcolsep) * \real{0.2153}}
  >{\raggedright\arraybackslash}p{(\linewidth - 4\tabcolsep) * \real{0.5088}}
  >{\raggedright\arraybackslash}p{(\linewidth - 4\tabcolsep) * \real{0.2759}}@{}}
\toprule\noalign{}
\begin{minipage}[b]{\linewidth}\raggedright
Source/\allowbreak{}declarat\allowbreak{}ion
\end{minipage} & \begin{minipage}[b]{\linewidth}\raggedright
Actual hypotheses and conclusion
\end{minipage} & \begin{minipage}[b]{\linewidth}\raggedright
Scope consequence
\end{minipage} \\
\midrule\noalign{}
\endhead
\bottomrule\noalign{}
\endlastfoot
\texttt{MaxMin.\allowbreak{}lean}: \texttt{IsMinima\allowbreak{}lIPGen},
\texttt{exists\_\allowbreak{}minimalI\allowbreak{}PGen\_\allowbreak{}subset}
& Finite hull has ambient interior origin; every proper subset loses it.
Every finite IP set contains such a subset. & Finite descent, without
lattice/\allowbreak{}reflexiv\allowbreak{}ity assumptions. \\
\texttt{Minimal.\allowbreak{}lean}: \texttt{IsGoodSi\allowbreak{}mplex}
& Affine independence together with intrinsic-interior origin. & Correct
relative notion for lower-dimensional simplex pieces. \\
\texttt{Minimal.\allowbreak{}lean}: \texttt{IsMinima\allowbreak{}lIPRel}
& Intrinsic-interior origin, and no proper subset with the same affine
span retaining intrinsic interior. & Minimality is measured in the fixed
parent span; deleting to a smaller span is not prohibited. \\
\texttt{CircuitG\allowbreak{}eometry.\allowbreak{}lean}:
\texttt{positive\allowbreak{}\_\allowbreak{}circuit\_\allowbreak{}through}
& A normalized nonnegative zero relation with a prescribed point's
coefficient positive yields an affinely independent supported positive
relation retaining that point. & Prescribed-point Carathéodory
reduction. \\
\texttt{CircuitG\allowbreak{}eometry.\allowbreak{}lean}:
\texttt{ip\_\allowbreak{}positive\allowbreak{}\_\allowbreak{}circuit\_\allowbreak{}through};
\texttt{Minimal.\allowbreak{}lean}:
\texttt{exists\_\allowbreak{}goodSimp\allowbreak{}lex\_\allowbreak{}through}
& Ambient IP and membership of a chosen generator yield a good simplex
through it. & Circuit coverage requires neither lattice points nor
minimality. \\
\texttt{CircuitI\allowbreak{}nterior.\allowbreak{}lean}:
\texttt{positive\allowbreak{}\_\allowbreak{}relation\allowbreak{}\_\allowbreak{}intrinsi\allowbreak{}cInterio\allowbreak{}r},
\texttt{exists\_\allowbreak{}positive\allowbreak{}\_\allowbreak{}relation\allowbreak{}\_\allowbreak{}of\_\allowbreak{}intrinsi\allowbreak{}cInterio\allowbreak{}r}
& Strict positive normalized zero relations imply intrinsic interior;
finite intrinsic-IP sets admit such relations. & Both directions are
present; affine independence is unnecessary here. \\
\texttt{CircuitI\allowbreak{}nterior.\allowbreak{}lean}:
\texttt{finset\_\allowbreak{}positive\allowbreak{}\_\allowbreak{}relation\allowbreak{}\_\allowbreak{}interior}
& Positive normalized relation \textbf{and full linear span} imply
ambient IP. & Positivity alone does not establish ambient interior. \\
\texttt{CircuitG\allowbreak{}eometry.\allowbreak{}lean}:
\texttt{conical\_\allowbreak{}caratheo\allowbreak{}dory} & A nonnegative
conical combination has a linearly independent positive support. & This
concerns nonzero conical representations, not an IP simplex's positive
zero relation. \\
\texttt{CircuitG\allowbreak{}eometry.\allowbreak{}lean}:
\texttt{projecte\allowbreak{}d\_\allowbreak{}positive\allowbreak{}\_\allowbreak{}circuit},
\texttt{projecte\allowbreak{}d\_\allowbreak{}relation\allowbreak{}\_\allowbreak{}coeffici\allowbreak{}ents}
& Circuit representatives in a linear projection, and proportionality of
every linear relation on an affinely independent projected simplex to
its normalized relation. & Supplies quotient geometry and relation
uniqueness. \\
\texttt{Minimal.\allowbreak{}lean}:
\texttt{minimalI\allowbreak{}PGen\_\allowbreak{}card\_\allowbreak{}lower},
\texttt{simplex\_\allowbreak{}card} & Ambient IP gives
\texttt{n+1 <=\allowbreak{} S.\allowbreak{}card}; adding affine
independence gives equality. & Lower bound does not need minimality. \\
\texttt{Minimal.\allowbreak{}lean}:
\texttt{simplex\_\allowbreak{}weights} & Affine independence and ambient
IP give an enumeration by \texttt{Fin(\allowbreak{}n+1)} and positive
normalized real zero relation. & The existential signature does not
itself assert rationality, primitive integral normalization, or
Newton-hull IP. \\
\texttt{Minimal.\allowbreak{}lean}:
\texttt{minimal\_\allowbreak{}structur\allowbreak{}e} & Minimal
ambient-IP set that is not affinely independent yields \texttt{S'}
proper, \texttt{T} contained, union equal to \texttt{S},
\texttt{IsMinima\allowbreak{}lIPRel S'},
\texttt{IsGoodSi\allowbreak{}mplex T},
\texttt{S.\allowbreak{}card =\allowbreak{} S'.\allowbreak{}card + (\allowbreak{}n -\allowbreak{} finrank(\allowbreak{}vectorSp\allowbreak{}an S')\allowbreak{})\allowbreak{} + 1},
and
\texttt{finrank(\allowbreak{}vectorSp\allowbreak{}an T)\allowbreak{} <=\allowbreak{} finrank(\allowbreak{}vectorSp\allowbreak{}an S')}.
& A proof body for the one-step KS statement is present. Relative rank
equals the relevant dimension because origin lies in the component
hull. \\
\end{longtable}

\texttt{Minimal.\allowbreak{}lean} also contains maximal-dimension
proper relative-IP selection and lifting lemmas used in its
decomposition proof. The general upper cardinality bound is proved in
\texttt{KS/\allowbreak{}SeedBoun\allowbreak{}d.\allowbreak{}lean}.

\subsection{A complete relation basis is stronger than circuit
coverage}\label{minimal-ip-circuits-and-types--a-complete-relation-basis-is-stronger-than-circuit-coverage}

Define
\texttt{A :\allowbreak{} R\textasciicircum{}S -\allowbreak{}> R\textasciicircum{}n}
by
\texttt{A(\allowbreak{}q)\allowbreak{} =\allowbreak{} sum q\_\allowbreak{}x x}.
\texttt{CircuitS\allowbreak{}pan.\allowbreak{}lean} defines
\texttt{simplexC\allowbreak{}ircuitWe\allowbreak{}ights S} as normalized
positive affine-simplex relations, extended by zero outside their
support, and \texttt{vertexCo\allowbreak{}mbinatio\allowbreak{}n S} as
\texttt{A}.

\texttt{normaliz\allowbreak{}ed\_\allowbreak{}relation\allowbreak{}\_\allowbreak{}mem\_\allowbreak{}span\_\allowbreak{}circuits}
and
\texttt{nonnegat\allowbreak{}ive\_\allowbreak{}relation\allowbreak{}\_\allowbreak{}mem\_\allowbreak{}span\_\allowbreak{}circuits}
first place nonnegative relations in the real span of those circuits.
\texttt{simplex\_\allowbreak{}circuits\allowbreak{}\_\allowbreak{}span\_\allowbreak{}kernel}
then proves, assuming intrinsic-IP,

\texttt{span(\allowbreak{}simplexC\allowbreak{}ircuitWe\allowbreak{}ights S)\allowbreak{} =\allowbreak{} ker(\allowbreak{}vertexCo\allowbreak{}mbinatio\allowbreak{}n S)}.

The final step uses a strictly positive relation on all coordinates: any
signed relation can be shifted by a sufficiently large multiple of it
into the nonnegative cone, then recovered by subtraction. This is why
coverage of coordinates by arbitrary rows is not a substitute for kernel
spanning.

\texttt{CWSBasis\allowbreak{}.\allowbreak{}lean} packages
\texttt{CircuitB\allowbreak{}asis S} with a finite row index
\texttt{Fin B.\allowbreak{}k}, linearly independent rows, a
positive-simplex witness for every row, and an equality between their
span and the complete kernel.
\texttt{exists\_\allowbreak{}circuitB\allowbreak{}asis} assumes
intrinsic-IP and returns
\texttt{Nonempty\allowbreak{} (\allowbreak{}CircuitB\allowbreak{}asis S)}
using an abstract independent-subset choice. It is not a computed basis
extraction routine.
\texttt{CircuitB\allowbreak{}asis.\allowbreak{}nonneg} gives nonnegative
padded rows; \texttt{CircuitB\allowbreak{}asis.\allowbreak{}cover} shows
every coordinate occurs positively somewhere; the converse is not
asserted.

\texttt{CircuitB\allowbreak{}asis.\allowbreak{}pairing\_\allowbreak{}range}
identifies the range of the vertex-pairing matrix with the joint kernel
of the selected row equations.
\texttt{CircuitB\allowbreak{}asis.\allowbreak{}coordina\allowbreak{}te\_\allowbreak{}count}
assumes ambient IP and proves
\texttt{S.\allowbreak{}card =\allowbreak{} n + B.\allowbreak{}k}. The
latter theorem does not require minimality despite its descriptive
comment. These are real linear-space statements. In {[}KS00{]} Lemma 4,
page 10, the corresponding finest lattice statement additionally uses
rational CWS data and the coarsest lattice generated by the vertices.
Forward integrality or real bijectivity alone does not certify a
two-sided unimodular lattice equivalence.

\texttt{CWSAssem\allowbreak{}bly.\allowbreak{}lean} contains abstract
finite row/CWS candidate sets and a \textbf{noncomputable}
\texttt{cwsCatal\allowbreak{}ogue}, followed by a noncomputable
dimension catalogue.
\texttt{finite\_\allowbreak{}ws\_\allowbreak{}cws\_\allowbreak{}containe\allowbreak{}r\_\allowbreak{}coverage}
states that for each dimension and family of selection rules there is a
bound and a finite catalogue supplying IP containers for every reflexive
polytope. Its displayed map is a real linear bijection with a forward
lattice-point preservation property and container inclusion. It does not
claim a computable bound, the optimized 47-entry structure list, every
lattice realization, exact normal forms, or a terminating native
end-to-end enumeration. Its dimension bound uses a separate
lattice-point-count theorem rather than the unported structural
\texttt{2n} bound.

\section{The archived 47 entries and the concrete 46 CWS
descriptors}\label{minimal-ip-circuits-and-types--the-archived-47-entries-and-the-concrete-46-cws-descriptors}

\texttt{Profiles\allowbreak{}.\allowbreak{}lean} computes lists with
atoms
\texttt{[\allowbreak{}6,\allowbreak{}5,\allowbreak{}4,\allowbreak{}3,\allowbreak{}2]},
lengths one through five, a nonincreasing predicate, and
\texttt{sum(\allowbreak{}m-\allowbreak{}1)\allowbreak{}=\allowbreak{}5}.
Its
\texttt{admissib\allowbreak{}leProfil\allowbreak{}e\_\allowbreak{}complete}
requires membership in that explicit search space. It verifies the
resulting seven arithmetic profiles; it does not derive those hypotheses
from a real minimal polytope.

\texttt{Structur\allowbreak{}es.\allowbreak{}lean} hard-codes four
\texttt{known3} and nine \texttt{known4} bases. The primary
small-dimensional classification is stated in {[}KS97{]} §4 and
{[}KS00{]} §2.1, Lemma 3, page 9 (the four-dimensional structures refer
to its appendix table). Those papers are the provenance for the intended
bases, but the archived declarations contain no geometric interpretation
theorem for them.

\texttt{extendTo\allowbreak{}5} extends the supplied bases with two new
vertices for \texttt{n'=\allowbreak{}4}, or three for
\texttt{n'=\allowbreak{}3}, attaches allowable old labels, and filters
only with \texttt{lemma2Ok}.
\texttt{rawFiveD\allowbreak{}imension\allowbreak{}alStruct\allowbreak{}ures}
adds the simplex. \texttt{canonica\allowbreak{}l} permutes the simplex
list, then relabels vertices in first-occurrence order;
\texttt{insertSt\allowbreak{}ructure} removes duplicate resulting lists.
The native-decide theorems verify that this defined computation has 47
entries, with vertex-count distribution
\texttt{(\allowbreak{}6,\allowbreak{}1)\allowbreak{},\allowbreak{}(\allowbreak{}7,\allowbreak{}6)\allowbreak{},\allowbreak{}(\allowbreak{}8,\allowbreak{}25)\allowbreak{},\allowbreak{}(\allowbreak{}9,\allowbreak{}13)\allowbreak{},\allowbreak{}(\allowbreak{}10,\allowbreak{}2)},
and the stated profile counts. No type mentions real vertices, affine
independence, intrinsic interior, or minimality.

Two concrete distinctions are essential:

\begin{enumerate}
\def\labelenumi{\arabic{enumi}.}
\tightlist
\item
  \textbf{Geometric soundness:} numbered row 8 of
  \texttt{paper/\allowbreak{}minimal\_\allowbreak{}polytope\allowbreak{}s\_\allowbreak{}5d\_\allowbreak{}table.\allowbreak{}tex}
  contains \texttt{\{1,\allowbreak{}2\}} and
  \texttt{\{1,\allowbreak{}2,\allowbreak{}3,\allowbreak{}4\}} as
  purported good-simplex supports, plus
  \texttt{\{5,\allowbreak{}6,\allowbreak{}7,\allowbreak{}8\}}. Proper
  containment of the first support in the second is impossible for two
  good simplices of the same vertex realization, by the barycentric
  uniqueness argument above. No-lonely passes because the first support
  has zero exclusive vertices and the larger block has two. Thus
  satisfying that filter does not supply geometric soundness. Other
  numbered entries likewise have contained blocks; they may be redundant
  equations or overgenerated candidates, but require separate
  interpretation rather than being called minimal-IP simplex
  decompositions.
\item
  \textbf{Unlabeled classification:} numbered rows 16, 17, and 18 have
  \texttt{\{1,\allowbreak{}2\}},
  \texttt{\{3,\allowbreak{}4,\allowbreak{}5\}}, and
  \texttt{\{j,\allowbreak{}6,\allowbreak{}7,\allowbreak{}8\}} for
  \texttt{j=\allowbreak{}3,\allowbreak{}4,\allowbreak{}5}. Swapping
  labels 3 and 4 or 3 and 5 while keeping the other labels fixed
  identifies their support hypergraphs. First-occurrence relabeling
  without all internal vertex permutations does not prove
  canonicalization under arbitrary vertex relabeling. Consequently 47 is
  an exact count of this computation's output lists, not a verified
  count of distinct unlabeled overlap hypergraphs.
\end{enumerate}

These are direct deductions from inspected table entries and
definitions, not results attributed to KS. The archived paper's
assertion that passing Lemma 2 is sufficient for a valid structure, and
its claim that this canonicalization quotients exactly by vertex
relabeling, need separate justification and are contradicted by the
examples above under the stated good-simplex/unlabeled interpretations.
This review does not alter the archival source or propose a replacement
count.

The 47-entry list includes the single-weight simplex. The project's 46
genuine CWS descriptors omit that single-weight case; subtraction alone
does not prove correspondence. Descriptor labels, overlaps, redundant
rows, orderings, row independence, component input pools and accepted
lattice realizations require a separate native-refinement argument.
Neither source arithmetic nor a container theorem proves that concrete
census complete beyond its declared input universe.

\section{Connection to the concrete
classification}\label{minimal-ip-circuits-and-types--connection-to-the-concrete-classification}

\begin{longtable}[]{@{}
  >{\raggedright\arraybackslash}p{(\linewidth - 2\tabcolsep) * \real{0.1885}}
  >{\raggedright\arraybackslash}p{(\linewidth - 2\tabcolsep) * \real{0.8115}}@{}}
\toprule\noalign{}
\begin{minipage}[b]{\linewidth}\raggedright
Connection
\end{minipage} & \begin{minipage}[b]{\linewidth}\raggedright
Formal guarantee
\end{minipage} \\
\midrule\noalign{}
\endhead
\bottomrule\noalign{}
\endlastfoot
Bounded minimal seed & \texttt{KS/\allowbreak{}SeedBoun\allowbreak{}d}
combines minimal decomposition and relative-coordinate transport to
prove \(n+1\le N\le2n\). \\
Complete primitive circuit rows & \texttt{KS/\allowbreak{}Circuits}
preserves supports, weighted-zero relations, independence and full
relation-kernel span. \\
Weight search and row assembly &
\texttt{KS/\allowbreak{}WeightSe\allowbreak{}arch} and
\texttt{KS/\allowbreak{}Assembly} enumerate every required ordered
primitive IP row and bounded independent covering CWS. \\
Compatible lattices and subpolytopes &
\texttt{KS/\allowbreak{}Recovery}, \texttt{KS/\allowbreak{}Lattices} and
\texttt{KS/\allowbreak{}Realizat\allowbreak{}ions} prove exact hull
recovery, finite lattice choices and reconstruction. \\
All reflexive targets &
\texttt{Refpoly.\allowbreak{}KS.\allowbreak{}classifi\allowbreak{}es\_\allowbreak{}reflexiv\allowbreak{}e}
in \texttt{KS/\allowbreak{}Obligati\allowbreak{}ons} composes the
executable routines into terminating coverage up to two-sided lattice
equivalence. \\
Additional promises & Ordered simplex-decomposition catalogues,
no-lonely pruning, native descriptor refinement and exact deduplication
are not supplied by the headline theorem. \\
\end{longtable}

The concrete enumeration conservatively assembles bounded relation rows;
it need not identify every output row family with a unique minimal-IP
simplex decomposition. The archived 47-entry computation is not used as
its completeness certificate. Nonreflexive one-point targets remain
outside the theorem.

\chapter{Weight enumeration, positive selection, and degree
bounds}\label{weight-enumeration-and-degree-bounds--weight-enumeration-positive-selection-and-degree-bounds}

\noindent\textit{Markdown source: }\texttt{docs/\allowbreak{}learning\allowbreak{}/\allowbreak{}weight-\allowbreak{}enumerat\allowbreak{}ion-\allowbreak{}and-\allowbreak{}degree-\allowbreak{}bounds.\allowbreak{}md}

This note explains positive selection, finite branching and degree
bounds in KS weight enumeration. The
\href{../../proofs/STATEMENTS.md}{concrete Lean classifier} uses exact
rational selection and rank-bounded search, together with CWS assembly,
compatible lattices and subpolytope recovery. The papers reviewed here
target the reflexive construction. Their nonreflexive Newton polytopes
do not establish coverage of all nonreflexive one-point lattice
polytopes.

\section{Primary sources and
conventions}\label{weight-enumeration-and-degree-bounds--primary-sources-and-conventions}

The main sources actually read are:

\begin{itemize}
\tightlist
\item
  Harald Skarke, \emph{Weight systems for toric Calabi--Yau varieties
  and reflexivity of Newton polyhedra}, \textbf{Modern Physics Letters A
  11} (1996), 1637--1652,
  \href{https://doi.org/10.1142/S0217732396001636}{DOI
  10.\allowbreak{}1142/\allowbreak{}S0217732\allowbreak{}39600163\allowbreak{}6}.
  Consulted edition:
  \href{https://arxiv.org/abs/alg-geom/9603007v1}{arXiv:\allowbreak{}alg-\allowbreak{}geom/\allowbreak{}9603007v\allowbreak{}1},
  7 March 1996. Locations below use its printed preprint page numbers.
\item
  Friedrich Schöller and Harald Skarke, \emph{All Weight Systems for
  Calabi--Yau Fourfolds from Reflexive Polyhedra},
  \textbf{Communications in Mathematical Physics 372} (2019), 657--678,
  \href{https://link.springer.com/article/10.1007/s00220-019-03331-9}{publisher
  record}, \href{https://doi.org/10.1007/s00220-019-03331-9}{DOI
  10.\allowbreak{}1007/\allowbreak{}s00220-\allowbreak{}019-\allowbreak{}03331-\allowbreak{}9}.
  Consulted edition:
  \href{https://arxiv.org/abs/1808.02422v1}{arXiv:\allowbreak{}1808.\allowbreak{}02422v1},
  7 August 2018, abbreviated SS18. The journal record verifies online
  publication on 12 February 2019 and December 2019 issue date.
  References to SS18 pages use the preprint's printed numbering; PDF
  page is printed page plus one.
\item
  K. Soundararajan, \emph{Approximating 1 from below using n Egyptian
  fractions},
  \href{https://arxiv.org/abs/math/0502247v1}{arXiv:\allowbreak{}math/\allowbreak{}0502247v\allowbreak{}1},
  11 February 2005,
  \href{https://doi.org/10.48550/arXiv.math/0502247}{arXiv DOI
  10.\allowbreak{}48550/\allowbreak{}arXiv.\allowbreak{}math/\allowbreak{}0502247}.
  This is a two-page author note, described by its arXiv record as not
  intended for publication, not a journal article.
\end{itemize}

Let \texttt{n} denote \textbf{number of weights}, so a single weight
system normally describes a polytope of dimension
\texttt{n-\allowbreak{}1}. Skarke 1996 instead uses \texttt{l} for
number of weights and \texttt{n} for polytope dimension. Write positive
normalized weights \texttt{q\_\allowbreak{}i} with
\texttt{∑q\_\allowbreak{}i=\allowbreak{}1}, and primitive positive
integral weights \texttt{w\_\allowbreak{}i},
\texttt{gcd(\allowbreak{}w\_\allowbreak{}i)\allowbreak{}=\allowbreak{}1},
degree \texttt{D=\allowbreak{}∑w\_\allowbreak{}i}, with
\texttt{q\_\allowbreak{}i=\allowbreak{}w\_\allowbreak{}i/\allowbreak{}D}.

In SS18 §3, p.~5, equations (5)--(7), the integer Newton hull is

\[
\Delta_q=\operatorname{conv}\{x\in\mathbb Z_{\ge0}^n:x\cdot q=1\}.
\]

Its IP condition is that
\texttt{ones=\allowbreak{}(\allowbreak{}1,\allowbreak{}…,\allowbreak{}1)}
lies in its \textbf{relative} interior in
\texttt{\{x:\allowbreak{}x·q=\allowbreak{}1\}}. In shifted coordinates
\texttt{y=\allowbreak{}x-\allowbreak{}ones}, the equivalent condition is
origin interior in the lattice hull with \texttt{q·y=\allowbreak{}0} and
\texttt{y\_\allowbreak{}i≥-\allowbreak{}1}, using the topology of
\texttt{W\_\allowbreak{}q=\allowbreak{}ker(\allowbreak{}q·−)}. Ambient
interior in \texttt{R\textasciicircum{}n} would be empty. Positive
weights put \texttt{ones} in the real intercept simplex; they do not
alone put it in the integer Newton hull's relative interior. IP checking
remains an exact separate predicate.

\section{The finite search argument and its
hypotheses}\label{weight-enumeration-and-degree-bounds--the-finite-search-argument-and-its-hypotheses}

Skarke 1996 §2, pp.~4--5, describes the basic separating-functional
search, candidate generation and subsequent IP filtering. SS18 §3,
pp.~5--7, states the improved recursion precisely. In the normalization
\texttt{r=\allowbreak{}1}, a finite set of chosen nonnegative lattice
points \texttt{S}, initially \texttt{\{ones\}}, determines

\[
Q_S=\{p:p_i\ge0,\ \sum p_i=1,\ x\cdot p=1\ (x\in S)\}.
\]

This is a bounded rational polytope because its equations have integer
coefficients and the sum constraint bounds coordinates. Empty
feasibility terminates a branch. For a viable node, the candidate must
be \textbf{strictly positive}, not just feasible with nonnegative
entries, to obtain finite branching. SS18 takes the arithmetic mean of
all vertices of \texttt{Q\_\allowbreak{}S} (equation (8), p.~6). If
\texttt{Q\_\allowbreak{}S} contains a point positive in every
coordinate, then for each coordinate at least one vertex has positive
value, and every vertex is nonnegative. Their mean is consequently
strictly positive. The vertices and this mean are rational, permitting
exact arithmetic.

For positive \texttt{p}, the next-point set

\[
B_p=\{x\in\mathbb Z_{\ge0}^n:x\cdot p<1\}
\]

is finite: \texttt{0≤x\_\allowbreak{}i<1/\allowbreak{}p\_\allowbreak{}i}
for every coordinate. If \texttt{p\_\allowbreak{}j=\allowbreak{}0}, all
multiples of \texttt{e\_\allowbreak{}j} lie in
\texttt{B\_\allowbreak{}p}; nonnegative feasibility alone does not
justify this finiteness claim. A rational positive selection rule gives
computable finite integer coordinate bounds; an arbitrary real-valued
choice does not itself provide an algorithm for comparisons or
enumeration.

The completeness argument is the \textbf{descent lemma} on SS18 p.~6.
Fix an IP target \texttt{q} whose Newton points include \texttt{S}. If
the selected normalized \texttt{p} differs from \texttt{q}, there must
be an integer Newton point \texttt{x} of \texttt{q} with \texttt{x·p<1}.
Otherwise \texttt{p·−} would attain its minimum at the interior point
\texttt{ones} of \texttt{Δ\_\allowbreak{}q}, forcing it to be constant
on the target's affine span; normalization would then give
\texttt{p=\allowbreak{}q}. Retaining all such children therefore retains
a path to each target.

Every feasible step adds an independent equation. All old points satisfy
\texttt{z·p=\allowbreak{}1}, while the new point does not, so it lies
outside their affine span. Because that old affine span is in a
hyperplane not containing zero, affine independence there corresponds to
linear independence of the point vectors. At \texttt{n} linearly
independent points, the equations determine a unique \texttt{q}. SS18
stops the branch at \texttt{k+1=\allowbreak{}n}, with branching levels
\texttt{k=\allowbreak{}0,\allowbreak{}…,\allowbreak{}n-\allowbreak{}2}.
Recording an early IP candidate does not justify discarding all
descendants: other targets can still satisfy the accumulated equations
when the feasible region has positive dimension.

The \textbf{whole-tree termination} conclusion needs bounded depth,
finite branching, and stopping on empty/nonpositive feasibility.
Target-path descent by itself does not prove that an arbitrary tree
exploration terminates. A complete finite implementation also requires
exact IP decisions on its candidates. Soundness is the assertion that
the list records only candidates passing that predicate; it needs a link
to the implemented checker, not merely the comment that PALP checks IP.

A normalized IP vector is rational even if initially defined over the
reals: its integer Newton hull spans the affine hyperplane containing
\texttt{ones}, so one can choose \texttt{n} linearly independent integer
Newton points and solve an invertible integer linear system
\texttt{Xq=\allowbreak{}ones}. This rationality argument is a
mathematical deduction from full relative dimension, formalized in
\texttt{KS/\allowbreak{}Rational\allowbreak{}Weights.\allowbreak{}lean}.
Clearing denominators and making the weights primitive gives a unique
integral encoding of such a rational ray.

\section{Half-weight reduction and the domains of pruning
rules}\label{weight-enumeration-and-degree-bounds--half-weight-reduction-and-the-domains-of-pruning-rules}

SS18 §3, pp.~5--6, gives necessary conditions in its
positive-dimensional setting (\texttt{n≥2}) for normalized IP weights:
\texttt{q\_\allowbreak{}i≤1/\allowbreak{}2}; for \texttt{n>2} at most
one entry is \texttt{1/\allowbreak{}2}. If an entry is larger, Newton
points have that coordinate at most one, so \texttt{ones} lies on a
supporting coordinate hyperplane. The two-equal-halves restriction uses
positivity of the remaining weights and normalization; it excludes no
general \texttt{n=\allowbreak{}2} exception, where
\texttt{(\allowbreak{}1/\allowbreak{}2,\allowbreak{}1/\allowbreak{}2)}
is possible.

For
\texttt{q=\allowbreak{}(\allowbreak{}q',\allowbreak{}1/\allowbreak{}2)}
with \texttt{∑q'\_\allowbreak{}i=\allowbreak{}1/\allowbreak{}2}, the
reduced integer hull uses \texttt{x'·q'=\allowbreak{}1} and interior
point \texttt{(\allowbreak{}2,\allowbreak{}…,\allowbreak{}2)}. Thus the
split is a six-weight \texttt{r=\allowbreak{}1} run plus a five-weight
\texttt{r=\allowbreak{}1/\allowbreak{}2} run, appending the half-weight
at the end. It is not an ordinary depth-one IP problem for the truncated
weights. In shifted coordinates
\texttt{z=\allowbreak{}x'-\allowbreak{}2ones}, it is the
\textbf{depth-two} hull \texttt{q'·z=\allowbreak{}0},
\texttt{z\_\allowbreak{}i≥-\allowbreak{}2}. Positively scaling
\texttt{q'} to sum one preserves this kernel but does not change depth
two into depth one. In the reduced
\texttt{r=\allowbreak{}1/\allowbreak{}2} run, feasibility is encoded by
\texttt{x\textasciicircum{}(\allowbreak{}0)\allowbreak{}=\allowbreak{}2ones},
while the right-hand side of the Newton equation remains one.

SS18 §3, p.~7, lists seven rejection rules, organized by cost:
coordinate and total-sum exclusions; orbit representatives under
permutations stabilizing old points; lexicographic violations of
previous minimum choices; a line extension test; failure of positive
feasibility; and the corresponding minimum-choice test over integer
points of the accumulated affine span. The canonical path chooses the
Newton point minimizing the current candidate pairing, breaking ties by
the lexicographically largest point. This supplies the rationale for
retaining a path while reducing duplicates, but requires separate proofs
for the actual tests and their domains. SS18 §5, p.~11, also says some
redundancy checks were turned off near the bottom of the tree and
remaining duplicates removed by sorting. The algorithm is not claimed to
generate each target once before that final sorting.

Two scope details matter for a faithful port:

\begin{itemize}
\tightlist
\item
  The stabilizer/orbit rule assumes a permutation-compatible search,
  including its candidate choices. It cannot automatically be attached
  to a completely arbitrary selection rule whose behavior breaks
  permutation symmetry.
\item
  The total-sum rule \texttt{∑x\_\allowbreak{}i≤2} uses the split search
  where remaining candidate weights are below \texttt{1/\allowbreak{}2}.
  In the unrestricted normalized feasible region,
  \texttt{(\allowbreak{}2,\allowbreak{}0,\allowbreak{}…,\allowbreak{}0)}
  \textbf{does} admit a positive solution with
  \texttt{q\_\allowbreak{}1=\allowbreak{}1/\allowbreak{}2}. Therefore
  SS18's observation that this rule implies failed positive feasibility
  must retain the split-domain restrictions. It is not a lemma about the
  frozen unrestricted \texttt{qPolytop\allowbreak{}e} alone.
\end{itemize}

\section{Supporting formal weight
theory}\label{weight-enumeration-and-degree-bounds--supporting-formal-weight-theory}

Paths are relative to \texttt{proofs/\allowbreak{}Refpoly/};
\href{../../proofs/SOURCES.md}{SOURCES.md} identifies the original
editions. Names below lie in \texttt{Refpoly}, except the explicit
\texttt{WeightSy\allowbreak{}stemFini\allowbreak{}te} subnamespace. The
abstract statements here are distinct from the executable definitions in
\texttt{KS/\allowbreak{}WeightSe\allowbreak{}arch}.

\begin{longtable}[]{@{}
  >{\raggedright\arraybackslash}p{(\linewidth - 2\tabcolsep) * \real{0.2005}}
  >{\raggedright\arraybackslash}p{(\linewidth - 2\tabcolsep) * \real{0.7995}}@{}}
\toprule\noalign{}
\begin{minipage}[b]{\linewidth}\raggedright
Declaration
\end{minipage} & \begin{minipage}[b]{\linewidth}\raggedright
Exact scope and limit
\end{minipage} \\
\midrule\noalign{}
\endhead
\bottomrule\noalign{}
\endlastfoot
\texttt{WeightSy\allowbreak{}stem.\allowbreak{}lean},
\texttt{IsIPWeig\allowbreak{}htSystem\allowbreak{}At},
\texttt{IsIPWeig\allowbreak{}htSystem} & Strict positivity plus origin
interior of an integer convex hull \textbf{in the subtype}
\texttt{wsHyperp\allowbreak{}lane q}; ordinary IP uses depth one. \\
\texttt{Algorith\allowbreak{}m.\allowbreak{}lean},
\texttt{qPolytop\allowbreak{}e}, \texttt{Selectio\allowbreak{}nRule} &
Feasible vectors are nonnegative with sum one; selection membership is
required only when feasibility is nonempty. Strict positivity and
rationality are not fields of \texttt{Selectio\allowbreak{}nRule}. \\
\texttt{Algorith\allowbreak{}m.\allowbreak{}lean},
\texttt{Reachabl\allowbreak{}e} & Root \texttt{\{ones\}} and insertion
of any nonnegative lattice \texttt{x} with
\texttt{x·sel(\allowbreak{}S)\allowbreak{}<1}. The inductive relation
has no guard requiring positive feasibility or nonempty child
feasibility. \\
\texttt{average\_\allowbreak{}pos} & Positivity of a mean of a supplied
finite family of nonnegative vectors, provided every coordinate is
positive somewhere. It does not construct all feasible-polytope vertices
or prove they are rational. \\
\texttt{branch\_\allowbreak{}finite} & Requires
\texttt{∀i,\allowbreak{}0<p\_\allowbreak{}i}; proves the branch set is
finite. It does not prove finite branching for every
\texttt{Selectio\allowbreak{}nRule}. \\
\texttt{descent\_\allowbreak{}lemma} & For normalized IP \texttt{q} and
normalized \texttt{p≠q}, produces \texttt{x∈wsPoints\allowbreak{} q}
with \texttt{x·p<1}; no positivity hypothesis on \texttt{p} is needed
for this target-path existence. \\
\texttt{reach\_\allowbreak{}step},
\texttt{notMem\_\allowbreak{}affineSp\allowbreak{}an\_\allowbreak{}of\_\allowbreak{}dot\_\allowbreak{}lt},
\texttt{finrank\_\allowbreak{}vectorSp\allowbreak{}an\_\allowbreak{}lt\_\allowbreak{}insert}
& Target-compatible reachable nodes either select the target or have a
target-compatible child of larger affine rank. \\
\texttt{algorith\allowbreak{}m\_\allowbreak{}complete} & For
\textbf{every} \texttt{Selectio\allowbreak{}nRule R} and normalized IP
target \texttt{q}, obtains
\texttt{∃ S,\allowbreak{} Reachabl\allowbreak{}e R S ∧ R.\allowbreak{}sel S=\allowbreak{}q ∧ S.\allowbreak{}card≤n+1}.
This finite witness/\allowbreak{}cardinal\allowbreak{}ity statement does
not return a list or prove the full unrestricted tree finite. The
printed bound is looser than the paper's \texttt{n}-point stopping
argument. \\
\texttt{WeightSy\allowbreak{}stemFini\allowbreak{}te.\allowbreak{}Positive\allowbreak{}Feasible},
\texttt{positive\allowbreak{}Selectio\allowbreak{}n} & Positive
feasibility is existence of a strictly positive member of
\texttt{qPolytop\allowbreak{}e S};
\texttt{positive\allowbreak{}Selectio\allowbreak{}n} chooses one by
classical choice, with fallback feasible choices or zero. It is
\texttt{noncompu\allowbreak{}table}. \\
\texttt{WeightSy\allowbreak{}stemFini\allowbreak{}te.\allowbreak{}nodes},
\texttt{nodes\_\allowbreak{}finite} & A depth-indexed finite set union
branches only at positively feasible parents; the chosen candidate is
then positive. This is a guarded finite catalogue construction in
\texttt{Set}, not executable enumeration. \\
\texttt{normaliz\allowbreak{}ed\_\allowbreak{}ip\_\allowbreak{}weight\_\allowbreak{}systems\_\allowbreak{}finite}
&
\texttt{\{q :\allowbreak{} V n :\allowbreak{} ∑q\_\allowbreak{}i=\allowbreak{}1 ∧ IsIPWeig\allowbreak{}htSystem\allowbreak{} q\}.\allowbreak{}Finite},
obtained by covering targets with the finite image of candidates in
\texttt{nodes n n}. It does not use Hensley or an assumed numerical
degree cutoff. \\
\texttt{WeightSy\allowbreak{}stem.\allowbreak{}lean},
\texttt{ip\_\allowbreak{}half\_\allowbreak{}reductio\allowbreak{}n} &
For positive \texttt{q'} with sum \texttt{1/\allowbreak{}2},
\texttt{IsIPWeig\allowbreak{}htSystem\allowbreak{} (\allowbreak{}snoc q' (\allowbreak{}1/\allowbreak{}2)\allowbreak{})\allowbreak{} ↔ IsIPWeig\allowbreak{}htSystem\allowbreak{}At q' 2}.
This includes the integer-hull and interior geometry, not just weight
positivity. \\
\texttt{DegreeBo\allowbreak{}und.\allowbreak{}lean},
\texttt{normaliz\allowbreak{}edWeight\allowbreak{}s\_\allowbreak{}injectiv\allowbreak{}e\_\allowbreak{}on\_\allowbreak{}primitiv\allowbreak{}e}
& Equal normalized values imply equal natural weight vectors when both
have gcd one and positive total degree. \\
\texttt{primitiv\allowbreak{}e\_\allowbreak{}ip\_\allowbreak{}weight\_\allowbreak{}systems\_\allowbreak{}finite},
\texttt{primitiv\allowbreak{}e\_\allowbreak{}ip\_\allowbreak{}degrees\_\allowbreak{}bounded}
& Fixed-length positive primitive integral IP vectors form a finite set,
hence
\texttt{∃D,\allowbreak{} ∀w,\allowbreak{} … → ∑w\_\allowbreak{}i≤D}.
This is existential boundedness without a supplied numerical value. \\
\end{longtable}

For example, a selection rule may pick a vertex with a zero coordinate
at the root. Its branch relation can then have infinitely many children.
At an infeasible node the rule is unconstrained, so
\texttt{Reachabl\allowbreak{}e} can also include arbitrary further
insertions. Such behavior does not refute
\texttt{algorith\allowbreak{}m\_\allowbreak{}complete}: that theorem
only needs one bounded witness path for each target. It does refute
reading its source comment as a theorem about total executable
termination for all selection rules.

The audited port should preserve the supported abstract theorem while
placing the claims about a runnable, terminating list in a separate
refinement theorem: computable positive rational choice; decidable
feasibility and IP tests; finite enumeration of branches; stopping; safe
pruning; and exact output normalization/permutation handling. The
current PALP/reference-lattice native scope is narrower still than
all-dimensional KS coverage.

\section{Census numbers and the status of
3,263,442}\label{weight-enumeration-and-degree-bounds--census-numbers-and-the-status-of-3263442}

SS18 abstract and §6, pp.~12--13, report \textbf{322,383,760,930}
six-weight IP systems, of which \textbf{185,269,499,015} directly
determine reflexive polytopes and \textbf{137,114,261,915} determine
nonreflexive IP polytopes. The reflexive systems yield
\textbf{532,600,483} distinct Hodge-data sets. These are published
enumeration results, not independently reproduced measurements here.
SS18 p.~13 explicitly says distinct weight systems can define the same
polytope and that their frequency was not checked. None of these weight
counts is the number of distinct five-dimensional reflexive polytopes,
and the Hodge count is neither a polytope nor an equivalence-class
count.

SS18 §5, pp.~11--12, distinguishes the \textbf{337,301,206,339
candidate} systems from the accepted IP census and describes
cross-checks against the degree-300 list and the transverse-weight list.
These support the published census but do not substitute for the
requested native-algorithm refinement proof.

The Sylvester sequence starts \texttt{a\_\allowbreak{}1=\allowbreak{}2}
and satisfies
\texttt{a\_\allowbreak{}(\allowbreak{}j+1)\allowbreak{}=\allowbreak{}1+∏\_\allowbreak{}(\allowbreak{}i≤j)\allowbreak{}a\_\allowbreak{}i=\allowbreak{}a\_\allowbreak{}j(\allowbreak{}a\_\allowbreak{}j-\allowbreak{}1)\allowbreak{}+1}.
SS18 §6, p.~14, equations (23)--(27), constructs the two Fermat
candidates

\[
q_{ct}^{(6)}=(1/2,1/3,1/7,1/43,1/1807,1/3263442),
\qquad
q_{lt}^{(6)}=(1/2,1/3,1/7,1/43,1/3612,1/3612).
\]

The paper gives their normalization and their location at the central
and left tips of its Hodge plot; it identifies Fermat weights as IP.
Indeed, when all
\texttt{q\_\allowbreak{}i=\allowbreak{}1/\allowbreak{}b\_\allowbreak{}i}
and their sum is one, the integral intercept vertices
\texttt{b\_\allowbreak{}i e\_\allowbreak{}i} are present, and
\texttt{ones=\allowbreak{}∑q\_\allowbreak{}i(\allowbreak{}b\_\allowbreak{}i e\_\allowbreak{}i)}
is a strict convex combination inside the full relative-dimensional
intercept simplex. This proves IP for these particular Newton hulls. It
does not prove a universal degree upper bound.

The frozen \texttt{Sylveste\allowbreak{}r.\allowbreak{}lean} proves
\texttt{prod\_\allowbreak{}sylveste\allowbreak{}r},
\texttt{sylveste\allowbreak{}r\_\allowbreak{}sum\_\allowbreak{}inv}, the
positive-entry and sum-one theorems \texttt{qct6\_\allowbreak{}pos},
\texttt{qct6\_\allowbreak{}sum}, \texttt{qlt6\_\allowbreak{}pos},
\texttt{qlt6\_\allowbreak{}sum}, and the arithmetic factorization
\texttt{3263442=\allowbreak{}1806·1807}. Those theorem types do
\textbf{not} assert IP, reflexivity, primitive degree, Hodge extremality
or maximal degree. \texttt{max\_\allowbreak{}degree\_\allowbreak{}dim5}
proves only \texttt{∃D,\allowbreak{} DegreeBo\allowbreak{}undDim5 D}.
\texttt{degree\_\allowbreak{}dim5\_\allowbreak{}le\_\allowbreak{}of\_\allowbreak{}certific\allowbreak{}ate}
assumes an explicit certificate; none for
\texttt{D=\allowbreak{}3263442} is supplied.

There is also a citation discrepancy: the frozen header attributes a
maximum degree of 3,263,442 to SS18 \textbf{§2}, but that section of the
consulted v1 does not state it. The number appears in \textbf{§6} as the
Fermat tip candidate. The reviewed sources therefore support that
candidate and arithmetic, not that particular attribution of a
universally sharp cutoff. If a complete computed catalogue elsewhere
reports 3,263,442 as its observed maximum, that report still needs
catalogue coverage/provenance or a mathematical bound to serve as a
certified universal cutoff. This note did not independently recompute or
verify the global observed maximum.

\subsection{Unit fractions: a verified theorem and a restricted
consequence}\label{weight-enumeration-and-degree-bounds--unit-fractions-a-verified-theorem-and-a-restricted-consequence}

D. R. Curtiss, \emph{On Kellogg's Diophantine Problem}, \textbf{American
Mathematical Monthly 29} (10) (November--December 1922), 380--387,
\href{https://www.jstor.org/stable/2299023}{JSTOR primary record},
\href{https://doi.org/10.2307/2299023}{DOI
10.\allowbreak{}2307/\allowbreak{}2299023}, is a verified bibliographic
lead. The original text was \textbf{not directly consulted}, so no
theorem number, exact statement or proof in Curtiss is asserted as
directly verified.

The actually read substitute is Soundararajan's note, p.~1, opening
statement and equation (1): for positive integral denominators
\texttt{b\_\allowbreak{}1≤…≤b\_\allowbreak{}k} with
\texttt{∑1/\allowbreak{}b\_\allowbreak{}i<1},

\[
\sum_{i=1}^k\frac1{b_i}\le\sum_{i=1}^k\frac1{a_i}
=1-\frac1{\prod_{i=1}^k a_i},
\]

with equality only for the Sylvester denominators. Repeated denominators
are allowed. The proof spans pp.~1--2, using its proposition comparing
products and sums of decreasing positive real sequences; this is an
independently consulted primary proof, not confirmation of the
inaccessible Curtiss proof.

As a deduction, for \texttt{n} genuine unit fractions summing to one,
delete the term with the largest denominator. The remaining
\texttt{n-\allowbreak{}1} terms sum below one, so the deleted fraction
is at least
\texttt{1/\allowbreak{}∏\_\allowbreak{}(\allowbreak{}i<n)\allowbreak{}a\_\allowbreak{}i}.
Every denominator is therefore at most
\texttt{∏\_\allowbreak{}(\allowbreak{}i<n)\allowbreak{}a\_\allowbreak{}i};
for six fractions this is 3,263,442, attained by the displayed Sylvester
completion. This only bounds denominators in the
\textbf{Fermat/unit-fraction subclass}. General primitive IP weights
have
\texttt{q\_\allowbreak{}i=\allowbreak{}w\_\allowbreak{}i/\allowbreak{}D}
with numerators not necessarily one. Moreover their primitive degree in
the Fermat subclass is \texttt{lcm(\allowbreak{}b\_\allowbreak{}i)}, not
in general the largest denominator. The denominator bound alone yields,
for example, the coarse effective degree bound
\texttt{lcm(\allowbreak{}1,\allowbreak{}…,\allowbreak{}3263442)} for
that subclass, not automatically the sharp degree cutoff 3,263,442 for
all IP weights. For the displayed \texttt{q\_\allowbreak{}ct}, the last
denominator is divisible by the others, so its degree does equal
3,263,442 after primitive clearing. That arithmetic fact about one
example must not be promoted to a universal theorem.

\section{Executable search and
scope}\label{weight-enumeration-and-degree-bounds--executable-search-and-scope}

\begin{longtable}[]{@{}
  >{\raggedright\arraybackslash}p{(\linewidth - 2\tabcolsep) * \real{0.1937}}
  >{\raggedright\arraybackslash}p{(\linewidth - 2\tabcolsep) * \real{0.8063}}@{}}
\toprule\noalign{}
\begin{minipage}[b]{\linewidth}\raggedright
Connection
\end{minipage} & \begin{minipage}[b]{\linewidth}\raggedright
Concrete guarantee or boundary
\end{minipage} \\
\midrule\noalign{}
\endhead
\bottomrule\noalign{}
\endlastfoot
Complete normalized target coverage &
\texttt{KS/\allowbreak{}WeightSe\allowbreak{}arch} returns every ordered
primitive IP code, using the descent argument and exact
normalization. \\
Finite branching &
\texttt{KS/\allowbreak{}Rational\allowbreak{}Selectio\allowbreak{}n}
computes a positive rational selector; exact arithmetic bounds every
branch coordinate. \\
Rank/\allowbreak{}depth stopping & Strict affine-rank growth justifies
the structural depth bound; infeasible nodes terminate without
children. \\
Soundness & Exact integer Newton-hull IP acceptance is proved equivalent
to \texttt{Primitiv\allowbreak{}eIP}; positivity alone is
insufficient. \\
Native pruning and half-weight splits & Optimized native rules require
their own domain, normalization and reconstruction guarantees; the
reference search does not depend on those optimizations. \\
Primitive degree bounds & Supporting results are existential. No
numerical universal certificate or sharp 3,263,442 cutoff is used by the
search. \\
Generic reflexive coverage &
\texttt{KS/\allowbreak{}Composit\allowbreak{}ion} and
\texttt{KS/\allowbreak{}Obligati\allowbreak{}ons} include exhaustive
CWS, compatible lattices and subpolytope recovery. \\
Census and other targets & Native refinement, exact distinct-class
counts and nonreflexive one-point coverage are separate claims. \\
\end{longtable}

\section{Cache and source
evidence}\label{weight-enumeration-and-degree-bounds--cache-and-source-evidence}

The following original-file hashes identify the consulted editions.
Local provenance records primary URLs, versions and retrieval date
\textbf{2026-10-14}; downloaded PDFs and author sources are ignored by
Git.

\begin{longtable}[]{@{}
  >{\raggedright\arraybackslash}p{(\linewidth - 4\tabcolsep) * \real{0.1655}}
  >{\raggedright\arraybackslash}p{(\linewidth - 4\tabcolsep) * \real{0.3597}}
  >{\raggedright\arraybackslash}p{(\linewidth - 4\tabcolsep) * \real{0.4748}}@{}}
\toprule\noalign{}
\begin{minipage}[b]{\linewidth}\raggedright
Source/\allowbreak{}version
\end{minipage} & \begin{minipage}[b]{\linewidth}\raggedright
Cache path relative to \texttt{docs/\allowbreak{}learning}
\end{minipage} & \begin{minipage}[b]{\linewidth}\raggedright
Original PDF SHA-256
\end{minipage} \\
\midrule\noalign{}
\endhead
\bottomrule\noalign{}
\endlastfoot
Skarke 1996 v1 &
\texttt{sources/\allowbreak{}skarke-\allowbreak{}alg-\allowbreak{}geom-\allowbreak{}9603007/\allowbreak{}v1/\allowbreak{}paper.\allowbreak{}pdf}
&
\texttt{ee4357a5\allowbreak{}eadd242e\allowbreak{}9689546a\allowbreak{}7991a9e5\allowbreak{}c291bb7b\allowbreak{}4627674e\allowbreak{}59865a7e\allowbreak{}80698800} \\
Schöller--Skarke 2018 v1 &
\texttt{sources/\allowbreak{}schoelle\allowbreak{}r-\allowbreak{}skarke-\allowbreak{}1808-\allowbreak{}02422/\allowbreak{}v1/\allowbreak{}paper.\allowbreak{}pdf}
&
\texttt{1272b653\allowbreak{}5b727824\allowbreak{}744d5bd7\allowbreak{}33b640f2\allowbreak{}46e03a5e\allowbreak{}374a7ea0\allowbreak{}cdd42867\allowbreak{}7001cc57} \\
Soundararajan 2005 v1 &
\texttt{sources/\allowbreak{}soundara\allowbreak{}rajan-\allowbreak{}math-\allowbreak{}0502247/\allowbreak{}v1/\allowbreak{}paper.\allowbreak{}pdf}
&
\texttt{96f1716e\allowbreak{}e0f84ed6\allowbreak{}13eb8118\allowbreak{}fa9199a4\allowbreak{}7f66fe93\allowbreak{}ea70d074\allowbreak{}1038da55\allowbreak{}ca8e1d32} \\
\end{longtable}

The Soundararajan source download is
\texttt{https:\allowbreak{}/\allowbreak{}/\allowbreak{}arxiv.\allowbreak{}org/\allowbreak{}src/\allowbreak{}math/\allowbreak{}0502247v\allowbreak{}1},
saved as \texttt{source-\allowbreak{}original} (gzip compressed
\texttt{Egyptian\allowbreak{}.\allowbreak{}tex}) with a supplementary
\texttt{source-\allowbreak{}extracte\allowbreak{}d.\allowbreak{}tex}.
Its original gzip SHA-256 is
\texttt{4c962562\allowbreak{}5c2f3cd6\allowbreak{}f2619f82\allowbreak{}e98b7650\allowbreak{}f215f860\allowbreak{}7e71209c\allowbreak{}e0f30361\allowbreak{}f25f8c51}.
Curtiss is a bibliographic reference only; no original PDF/text is
identified as directly verified.

Snapshot file hashes for the inspected declarations:

\begin{longtable}[]{@{}
  >{\raggedright\arraybackslash}p{(\linewidth - 2\tabcolsep) * \real{0.2747}}
  >{\raggedright\arraybackslash}p{(\linewidth - 2\tabcolsep) * \real{0.7253}}@{}}
\toprule\noalign{}
\begin{minipage}[b]{\linewidth}\raggedright
File in \texttt{Refpoly/}
\end{minipage} & \begin{minipage}[b]{\linewidth}\raggedright
SHA-256
\end{minipage} \\
\midrule\noalign{}
\endhead
\bottomrule\noalign{}
\endlastfoot
\texttt{Algorith\allowbreak{}m.\allowbreak{}lean} &
\texttt{0071873d\allowbreak{}ae87fa80\allowbreak{}5377c5a2\allowbreak{}4d4fddac\allowbreak{}29182630\allowbreak{}0030bfb0\allowbreak{}f454a157\allowbreak{}c85b568d} \\
\texttt{WeightSy\allowbreak{}stemFini\allowbreak{}te.\allowbreak{}lean}
&
\texttt{ee3474a8\allowbreak{}ede4ecaf\allowbreak{}80963191\allowbreak{}c95f1450\allowbreak{}a504f7cf\allowbreak{}6e62e60e\allowbreak{}1fb131c2\allowbreak{}16f811f1} \\
\texttt{WeightSy\allowbreak{}stem.\allowbreak{}lean} &
\texttt{c874f210\allowbreak{}1fb94463\allowbreak{}7922b022\allowbreak{}b765a357\allowbreak{}f96f38ba\allowbreak{}1a1595c7\allowbreak{}16e811df\allowbreak{}5ca5f85a} \\
\texttt{DegreeBo\allowbreak{}und.\allowbreak{}lean} &
\texttt{46819d08\allowbreak{}4db5bf1f\allowbreak{}f0397977\allowbreak{}3a72b7ae\allowbreak{}7ca9fc90\allowbreak{}c4bb6743\allowbreak{}cc932349\allowbreak{}09583103} \\
\texttt{Sylveste\allowbreak{}r.\allowbreak{}lean} &
\texttt{ae1c4efb\allowbreak{}04ed58b3\allowbreak{}ec54cd6f\allowbreak{}8c7f596e\allowbreak{}e423c73d\allowbreak{}1eaf04f8\allowbreak{}520eeb7d\allowbreak{}a92f2856} \\
\end{longtable}

The displayed Sylvester arithmetic is not a numerical census or a
certified universal maximum degree.

\chapter{CWS equations, containers, and lattice
realizations}\label{cws-and-lattice-realizations--cws-equations-containers-and-lattice-realizations}

\noindent\textit{Markdown source: }\texttt{docs/\allowbreak{}learning\allowbreak{}/\allowbreak{}cws-\allowbreak{}and-\allowbreak{}lattice-\allowbreak{}realizat\allowbreak{}ions.\allowbreak{}md}

This note explains how simplex weights become combined equations, how
complete relations recover a real pairing space, and how compatible
lattices recover every lattice realization. The
\href{../../proofs/STATEMENTS.md}{concrete classification} proves these
connections for its Lean reference algorithms. The
\hyperref[minimal-ip-circuits-and-types--minimal-ip-seeds-positive-circuits-and-finite-structure-data]{minimal-seed
note} treats decomposition, circuits and archived structure lists
separately.

\section{Primary source and
provenance}\label{cws-and-lattice-realizations--primary-source-and-provenance}

\textbf{{[}KS00{]}} Maximilian Kreuzer and Harald Skarke,
\emph{Reflexive polyhedra, weights and toric Calabi-Yau fibrations},
Reviews in Mathematical Physics \textbf{14}, no. 4 (April 2002),
343--374, \href{https://doi.org/10.1142/S0129055X0200120X}{DOI
10.\allowbreak{}1142/\allowbreak{}S0129055\allowbreak{}X0200120\allowbreak{}X}.
The DOI, authors, journal, issue, pages, and publication date were
verified against the publisher-deposited
\href{https://api.crossref.org/works/10.1142/S0129055X0200120X}{Crossref
record}. Publisher HTML could not be fetched; the paper actually read
was the primary \href{https://arxiv.org/abs/math/0001106v1}{arXiv
math/\allowbreak{}0001106v\allowbreak{}1}, submitted 19 January 2000,
\href{https://arxiv.org/pdf/math/0001106v1}{PDF}.

The consulted PDF is
\texttt{docs/\allowbreak{}learning\allowbreak{}/\allowbreak{}sources/\allowbreak{}kreuzer-\allowbreak{}skarke-\allowbreak{}math-\allowbreak{}0001106/\allowbreak{}v1/\allowbreak{}paper.\allowbreak{}pdf},
SHA-256
\texttt{7d400220\allowbreak{}d9bb4f0e\allowbreak{}0e2a121d\allowbreak{}46d75a7a\allowbreak{}ad0b0a98\allowbreak{}ebbc1d6a\allowbreak{}d45e3340\allowbreak{}7d0b01c1}.
Retrieval was 2026-10-14 UTC; \texttt{paper.\allowbreak{}txt} is a
derived reading aid. Source-license status is unverified. Pagination
below is the paper's printed pagination, not its PDF leaf index or
journal pagination.

The central sources are \textbf{§2.2}, pages 9--12, and \textbf{§2.3},
pages 13--17. In particular, the lattice classification algorithm is
Lemma 7 in §2.3, pages 13--15; §3 starts the geometric interpretation of
classification results on page 17. Confusing the section number would
direct a later implementation review to the wrong material.

Lean paths below are relative to \texttt{proofs/\allowbreak{}Refpoly/};
\href{../../proofs/SOURCES.md}{SOURCES.md} identifies the original
versions. All supporting declarations are in namespace \texttt{Refpoly};
filenames identify source locations rather than extra namespaces.

\section{Equations determine a space only when the relation family is
complete}\label{cws-and-lattice-realizations--equations-determine-a-space-only-when-the-relation-family-is-complete}

Let
\texttt{v\_\allowbreak{}1,\allowbreak{}.\allowbreak{}.\allowbreak{}.\allowbreak{},\allowbreak{}v\_\allowbreak{}N}
span an \texttt{n}-dimensional real vector space. Let
\texttt{A :\allowbreak{} R\textasciicircum{}N -\allowbreak{}> R\textasciicircum{}n}
take coefficients to
\texttt{sum\_\allowbreak{}i c\_\allowbreak{}i v\_\allowbreak{}i}. The
pairing map
\texttt{F :\allowbreak{} (\allowbreak{}R\textasciicircum{}n)\allowbreak{}* -\allowbreak{}> R\textasciicircum{}N}
takes \texttt{x} to the coordinates
\texttt{(\allowbreak{}<v\_\allowbreak{}i,\allowbreak{}x>)\allowbreak{}\_\allowbreak{}i}.
Its image is the annihilator of \texttt{ker A}. If rows of a matrix
\texttt{W} are relations among the vertices, then \texttt{im F} is
contained in \texttt{ker W}; equality requires their row span to equal
\textbf{all} of \texttt{ker A}.

For any real \texttt{k}-by-\texttt{N} matrix, rank-nullity gives

\texttt{dim(\allowbreak{}ker W)\allowbreak{} =\allowbreak{} N -\allowbreak{} rank(\allowbreak{}W)}.

If the rows are independent, this is \texttt{N-\allowbreak{}k}. An
arbitrary list of \texttt{k} blocks does not establish independence.
Coordinate coverage does not establish it either, and independent rows
need not constitute every relation of an independently supplied vertex
configuration. Redundant blocks must be removed or handled with rank
rather than raw row count.

{[}KS00{]} §2.2 defines a CWS from the positive normalized relations of
an \textbf{IP simplex structure}, padding each row with zeros off its
support. Its Lemma 4, page 10, assumes a minimal polytope and its
corresponding simplex-structure CWS. It identifies the image of the
pairing map with the simultaneous equation subspace, identifies the
polar with coordinate inequalities
\texttt{x\_\allowbreak{}i >=\allowbreak{} -\allowbreak{}1}, and
identifies the finest lattice when the CWS is rational. These
conclusions do not apply to an arbitrary nonnegative matrix merely
called a CWS. The IP simplex-structure hypotheses carry geometric
information.

The formal \texttt{WeightSy\allowbreak{}stem.\allowbreak{}lean}
structure \texttt{CWS k N} explicitly stores nonnegative rows, a
positive occurrence for every coordinate, and
\texttt{LinearIn\allowbreak{}dependen\allowbreak{}t R W}. It does not
store a vertex realization, affine independence of supports,
rationality, primitive integral row normalization, or a complete
relation-kernel witness. \texttt{CWS.\allowbreak{}subspace} is the
intersection of row kernels; \texttt{CWS.\allowbreak{}IsIP} says that
origin is interior, in that subspace, to the convex hull of its integer
points satisfying the coordinate inequalities. Thus full ambient
interior in \texttt{R\textasciicircum{}N} is neither intended nor
possible for nonzero equations.

\texttt{CWSBasis\allowbreak{}.\allowbreak{}lean} separately supplies
completeness. \texttt{CircuitB\allowbreak{}asis S} has independent
positive simplex rows whose span equals
\texttt{ker(\allowbreak{}vertexCo\allowbreak{}mbinatio\allowbreak{}n S)}.
\texttt{exists\_\allowbreak{}circuitB\allowbreak{}asis} obtains such a
basis abstractly from intrinsic-IP geometry.
\texttt{CircuitB\allowbreak{}asis.\allowbreak{}pairing\_\allowbreak{}range}
then proves equality with the joint row kernel;
\texttt{coordina\allowbreak{}te\_\allowbreak{}count} assumes ambient IP
and proves
\texttt{S.\allowbreak{}card =\allowbreak{} n + B.\allowbreak{}k}.
\texttt{CWSConta\allowbreak{}iner.\allowbreak{}lean} reindexes this
basis with the noncomputable \texttt{coeffEqu\allowbreak{}iv},
constructs \texttt{B.\allowbreak{}toCWS h0}, and proves that reindexing
preserves membership in the equations. These conclusions use the
stronger basis object, not the bare \texttt{CWS} fields.

\section{Component IP and the projection
hypothesis}\label{cws-and-lattice-realizations--component-ip-and-the-projection-hypothesis}

{[}KS00{]} §2.2, page 11, gives the corollary that an IP CWS has IP
constituent weight systems. The argument projects its lattice hull to
one support: the projected hull is contained in the corresponding
one-row hull. For origin to remain interior in the \textbf{entire
one-row solution space}, the restricted projection must be surjective
onto that space. Containment alone establishes interior only relative to
the image.

In \texttt{WeightSy\allowbreak{}stem.\allowbreak{}lean},
\texttt{cws\_\allowbreak{}ip\_\allowbreak{}componen\allowbreak{}t}
assumes an injective support map \texttt{e}, the exact positive support
of a row, agreement with the component weights \texttt{q'}, and the
explicit lifting hypothesis

\texttt{for every z with q' dot z =\allowbreak{} 0,\allowbreak{} there exists y in the joint kernel with y(\allowbreak{}e b)\allowbreak{}=\allowbreak{}z(\allowbreak{}b)}.

It returns \texttt{IsIPWeig\allowbreak{}htSystem\allowbreak{} q'}. That
lifting hypothesis is an essential part of the actual theorem; a
docstring's broad description must not hide it.

An independent elementary counterexample shows why a positive joint
dimension and independent covering rows do not suffice. Take
\texttt{N=\allowbreak{}3}, \texttt{k=\allowbreak{}2}, with normalized
rows

\texttt{q=\allowbreak{}(\allowbreak{}1,\allowbreak{}1,\allowbreak{}3)\allowbreak{}/\allowbreak{}5},
\texttt{r=\allowbreak{}(\allowbreak{}1,\allowbreak{}1,\allowbreak{}0)\allowbreak{}/\allowbreak{}2}.

They are nonnegative, independent, and cover all coordinates. The joint
kernel is
\texttt{\{(\allowbreak{}t,\allowbreak{}-\allowbreak{}t,\allowbreak{}0)\allowbreak{}\}}.
Its permitted integer points are
\texttt{(\allowbreak{}-\allowbreak{}1,\allowbreak{}1,\allowbreak{}0)},
\texttt{(\allowbreak{}0,\allowbreak{}0,\allowbreak{}0)},
\texttt{(\allowbreak{}1,\allowbreak{}-\allowbreak{}1,\allowbreak{}0)},
so its lattice hull has origin in its one-dimensional interior. The
full-support component \texttt{q}, however, is not IP: its integer
generators satisfy
\texttt{x\_\allowbreak{}1+x\_\allowbreak{}2+3x\_\allowbreak{}3=\allowbreak{}0}
and \texttt{x\_\allowbreak{}i>=\allowbreak{}-\allowbreak{}1}; any
positive integer \texttt{x\_\allowbreak{}3} would require
\texttt{x\_\allowbreak{}1+x\_\allowbreak{}2<=\allowbreak{}-\allowbreak{}3},
contradicting
\texttt{x\_\allowbreak{}1+x\_\allowbreak{}2>=\allowbreak{}-\allowbreak{}2}.
Thus every generator has \texttt{x\_\allowbreak{}3<=\allowbreak{}0}, so
origin cannot be interior in the two-dimensional component hyperplane.
The block projection has one-dimensional image and is not onto. This is
a counterexample to omitting the lifting hypothesis for arbitrary
matrices, not a counterexample to the geometric KS corollary.

The frozen geometric development supplies a different route.
\texttt{Geometri\allowbreak{}cCWS.\allowbreak{}lean} proves
\texttt{simplex\_\allowbreak{}relation\allowbreak{}\_\allowbreak{}kernel}
from affine independence, normalization, and the circuit equation.
\texttt{simplex\_\allowbreak{}pairing\_\allowbreak{}surjecti\allowbreak{}ve}
uses those same hypotheses to lift every vector in the component
hyperplane to a pairing with simplex vertices, even for a
lower-dimensional simplex. No ambient full-dimensionality of the simplex
is required.
\texttt{reflexiv\allowbreak{}e\_\allowbreak{}simplex\_\allowbreak{}block\_\allowbreak{}ip}
additionally assumes a reflexive polytope, positive weights, integral
simplex vertices, and their membership in the dual polytope; its proof
maps lattice generators into the component hull and uses the surjective
real map's openness to obtain IP. For a complete global circuit basis,
the global pairing image and this local pairing surjectivity together
provide the geometric lifting interpretation. Surjectivity on
\textbf{real spaces} is enough for this interior proof; it must not be
confused with surjectivity on integer lattices.

\section{Finest, coarsest, image, and saturated
lattices}\label{cws-and-lattice-realizations--finest-coarsest-image-and-saturated-lattices}

Assume rational/lattice vertices spanning \texttt{N\_\allowbreak{}R}.
Let \texttt{N\_\allowbreak{}0} be their integer span, the paper's
\texttt{N\_\allowbreak{}coarsest}, and let
\texttt{M\_\allowbreak{}f=\allowbreak{}N\_\allowbreak{}0*}, the paper's
\texttt{M\_\allowbreak{}finest}. Rationality and full span make these
full-rank discrete lattices. If the vertices are integral in an ambient
lattice \texttt{N}, then

\texttt{N\_\allowbreak{}0 <=\allowbreak{} N}, hence
\texttt{M=\allowbreak{}N* <=\allowbreak{} M\_\allowbreak{}f}.

The pairing embedding identifies \texttt{M\_\allowbreak{}f} with
\texttt{L\_\allowbreak{}f =\allowbreak{} ker(\allowbreak{}W)\allowbreak{} intersec\allowbreak{}t Z\textasciicircum{}N}
under the hypotheses of {[}KS00{]} Lemma 4(c). Integral coordinate
pairings characterize membership in the dual of
\texttt{N\_\allowbreak{}0}, which is why this is the \textbf{finest}
lattice. The image of an originally chosen lattice \texttt{M} can be a
proper finite-index sublattice of \texttt{L\_\allowbreak{}f}.

The distinction can be stated using saturation. In the rational joint
solution space, \texttt{L\_\allowbreak{}f} is saturated in the
coordinate lattice: if \texttt{a z} lies in the real kernel for a
nonzero integer \texttt{a}, then \texttt{z} lies there too. A full-rank
image lattice can nevertheless be unsaturated within
\texttt{L\_\allowbreak{}f}. Rational equations, an integral kernel
basis, and primitive/saturated image data must be tracked separately
from a real basis. The finite quotient
\texttt{L\_\allowbreak{}f/\allowbreak{}F(\allowbreak{}M)} records the
missing index data; it is not encoded just by row support or real row
span.

Here is a concrete example satisfying the geometric container setting.
In \texttt{R\textasciicircum{}2} with lattice
\texttt{Z\textasciicircum{}2}, let

\texttt{P=\allowbreak{}conv(\allowbreak{}±e\_\allowbreak{}1,\allowbreak{}±e\_\allowbreak{}2)}.

This reflexive diamond has dual square
\texttt{[\allowbreak{}-\allowbreak{}1,\allowbreak{}1]\allowbreak{}\textasciicircum{}2}.
Choose the vertex-minimal IP seed in that square consisting of
\texttt{±(\allowbreak{}1,\allowbreak{}1)} and
\texttt{±(\allowbreak{}1,\allowbreak{}-\allowbreak{}1)}. Two positive
segment circuits give rows
\texttt{(\allowbreak{}1,\allowbreak{}1,\allowbreak{}0,\allowbreak{}0)\allowbreak{}/\allowbreak{}2}
and
\texttt{(\allowbreak{}0,\allowbreak{}0,\allowbreak{}1,\allowbreak{}1)\allowbreak{}/\allowbreak{}2}.
The pairing map is

\texttt{F(\allowbreak{}x,\allowbreak{}y)\allowbreak{}=\allowbreak{}(\allowbreak{}x+y,\allowbreak{}-\allowbreak{}x-\allowbreak{}y,\allowbreak{}x-\allowbreak{}y,\allowbreak{}-\allowbreak{}x+y)}.

Its real image is the CWS space
\texttt{\{(\allowbreak{}a,\allowbreak{}-\allowbreak{}a,\allowbreak{}b,\allowbreak{}-\allowbreak{}b)\allowbreak{}\}},
and it is bijective onto that space. Integer
\texttt{(\allowbreak{}x,\allowbreak{}y)} give integer coordinates, but
the image lattice has the extra parity condition
\texttt{a=\allowbreak{}b mod 2} and index two in the full coordinate
lattice of the subspace. The integer vector
\texttt{(\allowbreak{}1,\allowbreak{}-\allowbreak{}1,\allowbreak{}0,\allowbreak{}0)}
has inverse
\texttt{(\allowbreak{}1/\allowbreak{}2,\allowbreak{}1/\allowbreak{}2)},
which is not integral. In fact \texttt{F(\allowbreak{}P)} equals the
real square container \texttt{|a|<=\allowbreak{}1},
\texttt{|b|<=\allowbreak{}1}; the missing assertion concerns its
lattice, not the real inclusion or hull.

This independent example explains exactly why the signatures of
\texttt{CircuitB\allowbreak{}asis.\allowbreak{}reflexiv\allowbreak{}e\_\allowbreak{}containe\allowbreak{}r},
\texttt{algorith\allowbreak{}m\_\allowbreak{}enumerat\allowbreak{}es\_\allowbreak{}cws\_\allowbreak{}containe\allowbreak{}r},
and the assembly coverage theorems cannot be read as unimodular
equivalence. They exhibit a real linear bijection, forward preservation
of lattice points, and \texttt{F '' P} contained in the CWS lattice
hull. They do not assert that every integer point of the target subspace
has an integral preimage. Requiring that extra assertion
indiscriminately would exclude valid examples; the correct completeness
route records the image sublattice and later explores lattice
realizations.

Duality reverses the lattice inclusions as well as polar inclusions.
Refining \texttt{N\_\allowbreak{}0} to a larger lattice \texttt{N}
corresponds to replacing its dual \texttt{M\_\allowbreak{}f} by a
smaller sublattice \texttt{M}. Depending on which side is described, the
associated finite quotient may be presented as a quotient/orbifold
condition or a sublattice restriction. An implementation must state the
side, the actual integer embedding, and the pairing compatibility; a
real change of coordinates alone is insufficient.

\section{Effective finite lattice choices: the exact scope of KS Lemma
7}\label{cws-and-lattice-realizations--effective-finite-lattice-choices-the-exact-scope-of-ks-lemma-7}

{[}KS00{]} §2.3 begins by stating that reflexive polytopes are recovered
as subpolytopes of suitable CWS containers, possibly on sublattices of
\texttt{M\_\allowbreak{}finest}. For a \textbf{fixed full-dimensional IP
real polytope} \texttt{Delta} and its polar, form the vertex-pairing
matrix \texttt{X}. If \texttt{X} is integral, both vertex-generated
coarsest lattices pair integrally, so

\texttt{M\_\allowbreak{}coarsest\allowbreak{} <=\allowbreak{} M <=\allowbreak{} M\_\allowbreak{}finest}

for any reflexive realization. This is now a bounded lattice interval:
its endpoints have the same rank and their quotient is finite. Full span
and integrality of the VPM matter. Arbitrarily choosing all sublattices
of \texttt{M\_\allowbreak{}finest} without the lower endpoint would give
infinitely many choices.

Lemma 7, page 13, diagonalizes \texttt{X} by unimodular integer row and
column operations: \texttt{X=\allowbreak{}W D U}, with \texttt{D} an
\texttt{n}-by-\texttt{n} positive diagonal integer matrix after
discarding zero rows/columns. It describes compatible lattices by
integer upper-triangular factorizations \texttt{D=\allowbreak{}T S}. The
columns of \texttt{S U} give primal vertices in lattice coordinates; the
rows of \texttt{W T} give dual vertices. The finest/coarsest
realizations correspond to \texttt{D U}/\texttt{U} on the primal side
and \texttt{W D}/\texttt{W} on the dual side. This is an actual
lattice-coordinate reconstruction, stronger than a real forward-integral
embedding.

The proof on page 14 selects diagonal divisors of \texttt{D}, bounded
off-diagonal representatives, and solves the triangular equations
\texttt{T S=\allowbreak{}D} for integral entries. There are finitely
many divisor choices and finitely many residue choices, giving an
effective finite enumeration for the fixed integral VPM. Equivalently,
one may enumerate subgroups of the finite quotient between the two
endpoint lattices and recover intermediate lattices. The concrete
\texttt{KS/\allowbreak{}Realizat\allowbreak{}ions} implementation uses
this finite-quotient route.

An indexing caution accompanies direct transcription: the extracted v1
statement/final proof inequality reads
\texttt{0<=\allowbreak{}T\_\allowbreak{}\{ji\}<T\_\allowbreak{}\{ii\}},
while the worked discussion specifies
\texttt{0<=\allowbreak{}T\_\allowbreak{}\{12\}<T\_\allowbreak{}\{11\}}.
The finite-factorization principle does not resolve those row/column
conventions. The Lean reference implementation uses finite-quotient
subgroup enumeration rather than treating the extracted inequalities as
a canonical matrix algorithm.

The next paragraph, page 15, says to construct all subpolytopes with
integral VPM of the CWS containers and apply Lemma 7. It supplies no
uniform finite-index enumeration for unrestricted sublattices
independently of a chosen polytope. The program needs a finite container
lattice-point search, exact hull/facet/VPM computations, integral-VPM
tests, and then the finite interval reconstruction. Source existence
claims and an implementable fixed-matrix procedure are different parts
of that chain.

The need is visible in the paper's own data: Proposition 2, page 17,
identifies a three-dimensional reflexive polytope requiring a
\texttt{Z\_\allowbreak{}2} sublattice realization beyond the
finest-lattice subpolytope list. This is historical evidence for the
missing lattice stage, not a new census result or a benchmark.

\section{Abstract reachability and finite
assembly}\label{cws-and-lattice-realizations--abstract-reachability-and-finite-assembly}

\begin{longtable}[]{@{}
  >{\raggedright\arraybackslash}p{(\linewidth - 4\tabcolsep) * \real{0.1766}}
  >{\raggedright\arraybackslash}p{(\linewidth - 4\tabcolsep) * \real{0.4649}}
  >{\raggedright\arraybackslash}p{(\linewidth - 4\tabcolsep) * \real{0.3584}}@{}}
\toprule\noalign{}
\begin{minipage}[b]{\linewidth}\raggedright
Declaration
\end{minipage} & \begin{minipage}[b]{\linewidth}\raggedright
Inspected conclusion
\end{minipage} & \begin{minipage}[b]{\linewidth}\raggedright
What it does not promise
\end{minipage} \\
\midrule\noalign{}
\endhead
\bottomrule\noalign{}
\endlastfoot
\texttt{Geometri\allowbreak{}cCWS.\allowbreak{}lean}:
\texttt{algorith\allowbreak{}m\_\allowbreak{}reaches\_\allowbreak{}simplex\_\allowbreak{}block}
& Each geometric dual-lattice circuit has a reachable WS node with
selected weights equal to its normalized row; node size is bounded by
support size plus one. & Actual native row ordering, pruning, integer
representation, or execution trace. \\
\texttt{Algorith\allowbreak{}m.\allowbreak{}lean}:
\texttt{Selectio\allowbreak{}nRule}, \texttt{Reachabl\allowbreak{}e} & A
rule picks from nonempty feasible regions; an inductive tree starts at
the ones vector and inserts admissible nonnegative lattice points under
a strict inequality. & A supplied computable rule, exhaustive executable
traversal, or equivalence with native stopping conditions. \\
\texttt{CWSConta\allowbreak{}iner.\allowbreak{}lean}:
\texttt{ReachedC\allowbreak{}ircuit},
\texttt{CircuitB\allowbreak{}asis.\allowbreak{}blocks\_\allowbreak{}reached}
& Each basis row has a genuine simplex support and a WS reachability
witness; length is at most \texttt{n+1}, or at most \texttt{n} for
nonsimplicial minimal seeds. & A chosen optimized overlap descriptor or
complete catalogue of geometric types. \\
\texttt{CWSConta\allowbreak{}iner.\allowbreak{}lean}:
\texttt{CWS.\allowbreak{}RowsReac\allowbreak{}hed},
\texttt{CircuitB\allowbreak{}asis.\allowbreak{}rows\_\allowbreak{}reached}
& Every row is reached and injectively zero-padded with exact
support/\allowbreak{}off-\allowbreak{}support equations. & A native
assembler's completeness or an independently checked relation-rank
computation. \\
\texttt{CWSAssem\allowbreak{}bly.\allowbreak{}lean}:
\texttt{candidat\allowbreak{}eRows\_\allowbreak{}finite},
\texttt{candidat\allowbreak{}eCWS\_\allowbreak{}finite} & Bounded row
length and coordina\allowbreak{}te/\allowbreak{}row counts give finite
mathematical candidate sets, using normalized IP weight finiteness. & An
explicit degree cutoff or executable finite list. \\
\texttt{cwsCatal\allowbreak{}ogue},
\texttt{dimensio\allowbreak{}nCWSCata\allowbreak{}logue} & Noncomputable
finite-set choices from those finiteness proofs. & A computed output,
canonical representatives, or proven equality with the 46
descriptors. \\
\texttt{ws\_\allowbreak{}cws\_\allowbreak{}catalogu\allowbreak{}e\_\allowbreak{}coverage},
\texttt{finite\_\allowbreak{}ws\_\allowbreak{}cws\_\allowbreak{}containe\allowbreak{}r\_\allowbreak{}coverage}
& Every reflexive polytope has a forward-integral real bijective
embedding into an IP container in the abstract catalogue, with an
existential dimension bound in the latter theorem. & Inverse
integrality, all image sublattices, all lattice realizations,
maximality, subpolytope enumeration, or terminating end-to-end code. \\
\end{longtable}

The abstract finite row catalogues avoid an empirical numerical degree
bound. \texttt{KS/\allowbreak{}WeightSe\allowbreak{}arch} and
\texttt{KS/\allowbreak{}Assembly} supply effective generation
separately. Native input-pool and descriptor coverage on the PALP
reference lattice is not implied by these generic results.

\section{Concrete reconstruction and target
boundary}\label{cws-and-lattice-realizations--concrete-reconstruction-and-target-boundary}

\begin{longtable}[]{@{}
  >{\raggedright\arraybackslash}p{(\linewidth - 2\tabcolsep) * \real{0.2108}}
  >{\raggedright\arraybackslash}p{(\linewidth - 2\tabcolsep) * \real{0.7892}}@{}}
\toprule\noalign{}
\begin{minipage}[b]{\linewidth}\raggedright
Connection
\end{minipage} & \begin{minipage}[b]{\linewidth}\raggedright
Formal guarantee or boundary
\end{minipage} \\
\midrule\noalign{}
\endhead
\bottomrule\noalign{}
\endlastfoot
Complete equations and dimension & \texttt{KS/\allowbreak{}Circuits}
retains full relation-kernel span; \texttt{KS/\allowbreak{}Assembly}
checks exact independence, dimension and coverage. \\
IP component projection & The geometric simplex theorem supplies real
surjectivity; bare covering rows still do not discharge
\texttt{hlift}. \\
Saturated kernel coordinates &
\texttt{KS/\allowbreak{}IntegerK\allowbreak{}ernel} proves integer
preimages of every integer kernel point, not just real spanning. \\
Original image lattice & \texttt{KS/\allowbreak{}Recovery} places the
target image lattice between the actual-vertex coarse and polar-defined
fine endpoints. \\
All compatible realizations & \texttt{KS/\allowbreak{}Lattices} proves
interval finiteness; \texttt{KS/\allowbreak{}Realizat\allowbreak{}ions}
exhausts the finite quotient and proves exact
hull/\allowbreak{}image-\allowbreak{}lattice reconstruction. \\
Finite subpolytope recovery &
\texttt{KS/\allowbreak{}IntegerP\allowbreak{}oints} enumerates all
container lattice points; \texttt{KS/\allowbreak{}Recovery} proves an
enumerated sublist recovers the target image. \\
Maximality & IP containers need not be reflexive or maximal; Lemma 8's
automatic reflexivity remains limited to \texttt{n<=\allowbreak{}4}. \\
Native enumeration and distinct classes & PALP refinement, optimized
pruning, native descriptors and exact normal-form deduplication are
separate guarantees; equal hashes are insufficient. \\
\end{longtable}

For reflexive targets, the integral polar supplies a minimal seed and
the compatible-lattice stage recovers the target's original lattice.
These stages compose in
\texttt{Refpoly.\allowbreak{}KS.\allowbreak{}classifi\allowbreak{}es\_\allowbreak{}reflexiv\allowbreak{}e}.
For general nonreflexive one-point polytopes, polar integrality is
unavailable, so the same argument does not extend automatically. The
projection and parity examples above are independent deductions, not
additional cited or formal theorems.

\chapter{Finiteness, volume bounds, and reduction to a finite lattice
box}\label{finiteness-and-box-reduction--finiteness-volume-bounds-and-reduction-to-a-finite-lattice-box}

\noindent\textit{Markdown source: }\texttt{docs/\allowbreak{}learning\allowbreak{}/\allowbreak{}finitene\allowbreak{}ss-\allowbreak{}and-\allowbreak{}box-\allowbreak{}reductio\allowbreak{}n.\allowbreak{}md}

This note concerns full-dimensional lattice polytopes in a fixed
lattice, first with a fixed positive number of interior lattice points,
then with the additional reflexivity hypothesis. The supporting Lean
declarations provide abstract volume and box results. The
\href{../../proofs/STATEMENTS.md}{concrete KS endpoint} uses a
separately proved executable construction; neither generic theorem is a
native descriptor census.

\section{Definitions and equivalence
conventions}\label{finiteness-and-box-reduction--definitions-and-equivalence-conventions}

Let the dimension be \texttt{n ≥ 1}, let
\texttt{M ≅ Z\textasciicircum{}n}, and let \texttt{P} be the convex hull
of finitely many points of \texttt{M}, with affine span
\texttt{M\_\allowbreak{}R}. Put

\[
k(P)=\#(\operatorname{int}P\cap M),\qquad
J(P)=\#(P\cap M).
\]

Here interior is ambient interior in \texttt{M\_\allowbreak{}R}. For a
lower-dimensional polytope, one must instead specify its intrinsic
lattice and relative interior before applying a theorem in that smaller
dimension. Ambient interior is empty for a proper affine subspace.

The \textbf{IP property} means \texttt{0 ∈ int P}. It specifies an
interior origin and implies \texttt{k(\allowbreak{}P)\allowbreak{} ≥ 1};
it does not imply
\texttt{k(\allowbreak{}P)\allowbreak{} =\allowbreak{} 1}. For example,
\texttt{[\allowbreak{}-\allowbreak{}r,\allowbreak{}r]\allowbreak{}\textasciicircum{}n},
with positive integral \texttt{r}, has the IP property and unbounded
volume as \texttt{r} grows. A \textbf{one-point lattice polytope} has
exactly one interior lattice point. Translating that point to zero makes
it IP. A \textbf{reflexive polytope} is an IP lattice polytope whose
polar is a lattice polytope in the dual lattice. It is a one-point
lattice polytope, but the converse need not hold.

Unmarked polytopes are compared by affine lattice isomorphisms
\texttt{x ↦ Ux + t}, with \texttt{U ∈ GL(\allowbreak{}n,\allowbreak{}Z)}
and \texttt{t ∈ Z\textasciicircum{}n}. Origin-marked polytopes are
compared by \texttt{x ↦ Ux}. For \texttt{k =\allowbreak{} 1}, every
affine equivalence sends the unique interior lattice point to the unique
such point; after centering both polytopes it has zero translation. For
\texttt{k > 1}, there is no distinguished interior point. Marking one
gives at most \texttt{k} possibilities per affine class, so marked GL
finiteness follows, but the marked and unmarked classifications are
different. Arbitrarily translating polytopes without requiring an
interior point at zero does not produce a finite GL classification.

Volume below is ordinary Euclidean volume in lattice coordinates: a
fundamental parallelepiped has volume one, and the standard lattice
simplex has volume \texttt{1/\allowbreak{}n!}. For a lattice of covolume
\texttt{c}, first divide Euclidean volume by \texttt{c}. The
integer-valued normalized volume convention is
\texttt{Vol\_\allowbreak{}Z(\allowbreak{}P)\allowbreak{} =\allowbreak{} n! vol(\allowbreak{}P)}.
These conventions differ by \texttt{n!}, which must be carried through
both the volume and box bounds.

\section{Verified primary
results}\label{finiteness-and-box-reduction--verified-primary-results}

\subsection{Hensley's positive-interior-point
bound}\label{finiteness-and-box-reduction--hensleys-positive-interior-point-bound}

Douglas Hensley, \emph{Lattice vertex polytopes with interior lattice
points}, \textbf{Pacific Journal of Mathematics 105} (1) (1983),
183--191, \href{https://msp.org/pjm/1983/105-1/p10.xhtml}{publisher
record},
\href{https://msp.org/pjm/1983/105-1/pjm-v105-n1-p10-s.pdf}{journal
PDF}, \href{https://doi.org/10.2140/pjm.1983.105.183}{DOI
10.\allowbreak{}2140/\allowbreak{}pjm.\allowbreak{}1983.\allowbreak{}105.\allowbreak{}183}.
The article byline is Douglas Hensley; the cover/index and publisher
record give Douglas Austin Hensley. The scanned issue says September
1983, whereas the current publisher record says 1 March 1983. The year,
volume, issue, pagination and DOI are verified; no resolution of this
month discrepancy is assumed.

Hensley constructs positive constants
\texttt{ε(\allowbreak{}k,\allowbreak{}n)} in §2, pp.~184--186: Lemma 2.3
treats \texttt{k =\allowbreak{} 1}, and Lemma 2.4 treats \texttt{k ≥ 2}.
Theorem 3.1, p.~186, bounds the barycentric coordinates of interior
lattice points of a lattice simplex away from zero using these
constants. Corollary 3.2, pp.~186--187, transfers the relevant
coefficient bound to a lattice pyramid. Lemma 3.3, p.~187, explains why
a simplex's relative interior lies in a containing convex body's ambient
interior when their relative interiors meet there; relative and ambient
interiors are not being identified unconditionally.

Theorem 3.4, p.~187, gives for a lattice simplex \texttt{S} with exactly
\texttt{k ≥ 1} interior lattice points

\[
\operatorname{vol}(S)\le \frac{k}{n!}\varepsilon(k,n)^{-n}.
\]

Theorem 3.6, p.~188, applies to \textbf{any convex lattice polytope}
\texttt{P} in \texttt{R\textasciicircum{}n} with
\texttt{k =\allowbreak{} \#(\allowbreak{}int P ∩ Z\textasciicircum{}n)\allowbreak{} ≥ 1}
and gives

\[
\operatorname{vol}(P)\le k\varepsilon(k,n)^{-n}.
\]

Corollary 3.7 on that same page gives

\[
J(P)\le n+k\,n!\,\varepsilon(k,n)^{-n}.
\]

The general-polytope argument controls asymmetry about an interior
lattice point and invokes a generalized Mahler bound; it is not just the
simplex formula with the same factorial. The scanned statements include
\texttt{k =\allowbreak{} 1}. OCR repeatedly renders \texttt{≥} as
\texttt{>} or \texttt{<}; the hypotheses, inequalities and factorials
above were checked against the scanned page. No sharp estimate or later
resolution of Hensley's conjectures is needed here.

\subsection{Lagarias--Ziegler: explicit volume and box
bounds}\label{finiteness-and-box-reduction--lagariasziegler-explicit-volume-and-box-bounds}

Jeffrey C. Lagarias and Günter M. Ziegler, \emph{Bounds for Lattice
Polytopes Containing a Fixed Number of Interior Points in a Sublattice},
\textbf{Canadian Journal of Mathematics 43} (5) (1991), 1022--1035,
\href{https://www.cambridge.org/core/journals/canadian-journal-of-mathematics/article/bounds-for-lattice-polytopes-containing-a-fixed-number-of-interior-points-in-a-sublattice/0D4E5B45AB64DC6069A9BD7DD45B450C}{publisher
record}, \href{https://doi.org/10.4153/CJM-1991-058-4}{DOI
10.\allowbreak{}4153/\allowbreak{}CJM-\allowbreak{}1991-\allowbreak{}058-\allowbreak{}4}.
Print issue date is 1 October 1991; Cambridge's online publication date,
20 November 2018, is not the paper's research publication year.

Use \texttt{s ≥ 1} for their integer sublattice multiplier, called
\texttt{d} in the paper, to avoid confusing it with dimension. Theorem
1, p.~1023, assumes lattice vertices in \texttt{Z\textasciicircum{}n}
and exactly \texttt{k ≥ 1} interior points belonging to
\texttt{s Z\textasciicircum{}n}. Its formulas (1.1)--(1.2) are

\[
V(n,k,s)\le k s^n\,[7(ks+1)]^{\,n2^{n+1}},
\qquad
J(n,k,s)\le n+n! k s^n\,[7(ks+1)]^{\,n2^{n+1}}.
\]

\texttt{J} counts \textbf{all} \texttt{Z\textasciicircum{}n} points, not
only sublattice points. The exponent is
\texttt{n · 2\textasciicircum{}(\allowbreak{}n+1)}, and the base starts
with \texttt{7}, not an OCR letter \texttt{l}. Both were visually
verified on p.~1023. With \texttt{s =\allowbreak{} 1}, define the
explicit integer bound

\[
C(n,k)=k[7(k+1)]^{\,n2^{n+1}}.
\]

Thus every full-dimensional lattice polytope with exactly \texttt{k ≥ 1}
interior lattice points has Euclidean volume at most
\texttt{C(\allowbreak{}n,\allowbreak{}k)}, normalized volume at most
\texttt{n! C(\allowbreak{}n,\allowbreak{}k)}, and total lattice-point
count at most \texttt{n+n! C(\allowbreak{}n,\allowbreak{}k)}. The proof
is in §2, pp.~1025--1031: one-sided Diophantine approximation,
barycentric/asymmetry bounds, then a convex-body volume argument. This
is an effective bound, without a practicality claim.

Theorem 2, p.~1024, states that a lattice polytope of volume at most
\texttt{V} can be sent by a \textbf{linear} integral unimodular
transformation into a lattice cube of side at most \texttt{n · n! · V}.
The cube need not be based at the origin. The paragraph immediately
after the theorem uses an \textbf{integer translation} to put the cube
in \texttt{0 ≤ x\_\allowbreak{}i ≤ n · n! · V} and obtains finite affine
integral equivalence classes. If preserving
\texttt{s Z\textasciicircum{}n} is required, it uses translations in
\texttt{s Z\textasciicircum{}n} and the enlarged bound
\texttt{n · n! · V + s}. The affine group and its sublattice-preserving
subgroup are defined earlier on p.~1024.

The theorem's proof in §3, p.~1032, reduces a lattice simplex by Hermite
normal form and then covers a general polytope using a maximum-volume
simplex. Theorem 3(a), p.~1025, for a closed bounded convex body
\texttt{K} with nonempty interior and any maximum-volume contained
simplex \texttt{S}, centroid \texttt{b}, gives

\[
K\subseteq -nS+(n+1)b,
\qquad K\subseteq (n+2)S-(n+1)b.
\]

Theorem 3(b) allows \texttt{S} to use vertices of a polytope \texttt{K};
when \texttt{K} is lattice, the two enlarged simplices are lattice
because \texttt{(\allowbreak{}n+1)\allowbreak{}b} is a sum of lattice
vertices. Proof: §4, pp.~1032--1034. This maximum-volume simplex is not
an inclusion-maximal member of the fixed-\texttt{k} target class.

The sublattice extension (1.3)--(1.4), pp.~1023--1024, is for a fixed
full-rank sublattice \texttt{Λ ⊆ Z\textasciicircum{}n}. Its reduction
uses an integral map with determinant
\texttt{s\textasciicircum{}n/\allowbreak{}det Λ}; that map need not be
unimodular. A change of lattice coordinates for deriving a volume bound
is therefore not itself an equivalence witness for a fixed-lattice
classification.

\section{Consequences for termination and maximal
containers}\label{finiteness-and-box-reduction--consequences-for-termination-and-maximal-containers}

These paragraphs are deductions from the verified results, not
additional numbered theorems in the papers.

For fixed \texttt{(\allowbreak{}n,\allowbreak{}k)} with \texttt{k ≥ 1},
put \texttt{L=\allowbreak{}n · n! · C(\allowbreak{}n,\allowbreak{}k)}.
There are finitely many lattice points in
\texttt{[\allowbreak{}0,\allowbreak{}L]\allowbreak{}\textasciicircum{}n}.
Enumerating their subsets, taking convex hulls, retaining full dimension
and testing the exact interior count produces a terminating finite list
covering every affine class. The list has duplicates until an exact
equivalence procedure is supplied. Exact rational polytope geometry can
implement these predicates, but neither paper provides a proof that this
particular repository implements that enumeration.

If \texttt{0 ∈ int P}, the \emph{linear} map in Theorem 2 preserves
zero. A cube containing its image contains zero, so each coordinate of
that cube lies within
\texttt{[\allowbreak{}-\allowbreak{}L,\allowbreak{}L]}. Therefore the
same explicit \texttt{L} gives a centered box without adding a
translation to the origin-marked object. This applies to every marked
fixed-\texttt{k} polytope, not just reflexive ones. With integer
normalized volume bound \texttt{B}, the corresponding cube side is
\texttt{n B}, not \texttt{n · n! · B}.

The positive-count restriction is essential. For \texttt{n ≥ 2}, boxes
\texttt{[\allowbreak{}0,\allowbreak{}m]\allowbreak{} × [\allowbreak{}0,\allowbreak{}1]\allowbreak{}\textasciicircum{}(\allowbreak{}n-\allowbreak{}1)}
have no interior lattice points and volume \texttt{m}. Their distinct
volumes preclude affine equivalence. Hence there is no general finite
classification at \texttt{k =\allowbreak{} 0} in these dimensions by
this route. In dimension one, an empty-interior lattice interval has
length one up to affine equivalence; this exception does not salvage a
uniform all-dimension claim. Nor is the union over all positive
\texttt{k} a finite list for fixed dimension.

Fix a class \texttt{C\_\allowbreak{}\{n,\allowbreak{}k\}} of
full-dimensional lattice polytopes on the same lattice, with exactly
\texttt{k} interior points; optionally require a marked origin. The
uniform bound on \texttt{J} proves that every \texttt{P} in this class
has an \textbf{inclusion-maximal superset within this class}: maximize
the integer \texttt{J(\allowbreak{}Q)} among members containing
\texttt{P}. The attainable values form a nonempty bounded set of natural
numbers. If lattice polytopes \texttt{Q ⊊ R}, then
\texttt{J(\allowbreak{}Q)\allowbreak{} < J(\allowbreak{}R)}, because
each is the convex hull of its lattice points and equal point sets would
give equal polytopes. A maximizer therefore admits no proper
class-preserving extension. This argument uses a numerical invariant on
actual embeddings; it does not order arbitrarily chosen class
representatives by literal inclusion.

There are only finitely many equivalence classes of these maximal
objects. Choose representatives
\texttt{M\_\allowbreak{}1,\allowbreak{}…,\allowbreak{}M\_\allowbreak{}t}.
The correct covering statement is
\texttt{g(\allowbreak{}P)\allowbreak{} ⊆ M\_\allowbreak{}i} for some
allowed equivalence \texttt{g} and some \texttt{i}. It is generally
false that every original \texttt{P} is literally contained in one of
the fixed representatives without a change of coordinates. For unmarked
objects \texttt{g} is affine; for centered reflexive or marked objects
it is origin-fixed GL.

Every lattice subpolytope of a bounded lattice container
\texttt{M\_\allowbreak{}i} is generated by a subset of the finite set
\texttt{M\_\allowbreak{}i ∩ M}. Thus finite subpolytope enumeration can
recover all targets, once the covering containers have actually been
constructed. It must allow \textbf{all lattice points} of the container
as generators. Subsets of only the container's vertices need not recover
a subpolytope with vertices at other lattice points. Container existence
does not identify the KS WS/CWS construction or supply an executable
list of containers.

\section{Supporting formal volume and box
statements}\label{finiteness-and-box-reduction--supporting-formal-volume-and-box-statements}

Paths are relative to \texttt{proofs/\allowbreak{}Refpoly/};
\href{../../proofs/SOURCES.md}{SOURCES.md} identifies the original
editions and hashes. Names below are in namespace \texttt{Refpoly},
except explicitly qualified subnamespaces. Line locations identify the
original source edition, not the current checkout.

\begin{longtable}[]{@{}
  >{\raggedright\arraybackslash}p{(\linewidth - 2\tabcolsep) * \real{0.2679}}
  >{\raggedright\arraybackslash}p{(\linewidth - 2\tabcolsep) * \real{0.7321}}@{}}
\toprule\noalign{}
\begin{minipage}[b]{\linewidth}\raggedright
File/\allowbreak{}declarat\allowbreak{}ion
\end{minipage} & \begin{minipage}[b]{\linewidth}\raggedright
Actual supported statement and limit
\end{minipage} \\
\midrule\noalign{}
\endhead
\bottomrule\noalign{}
\endlastfoot
\texttt{Basic.\allowbreak{}lean}, \texttt{IsIP},
\texttt{IsReflex\allowbreak{}ive}, \texttt{IsUnimod\allowbreak{}ular},
\texttt{GLEquiv} & \texttt{IsIP} uses ambient \texttt{interior};
reflexivity requires lattice primal and polar. \texttt{GLEquiv} is a
real linear equivalence with \textbf{both} forward and inverse lattice
preservation. It has no translation parameter. \\
\texttt{Basic.\allowbreak{}lean},
\texttt{interior\allowbreak{}\_\allowbreak{}latticeP\allowbreak{}oint\_\allowbreak{}unique}
&
\texttt{IsReflex\allowbreak{}ive P → interior\allowbreak{} P ∩ latticeP\allowbreak{}oints n =\allowbreak{} \{0\}}.
This bridges reflexivity to \texttt{k=\allowbreak{}1}; it does not
classify all one-point polytopes as reflexive. \\
\texttt{BoxReduc\allowbreak{}tion.\allowbreak{}lean:\allowbreak{}295},
\texttt{BoxReduc\allowbreak{}tion.\allowbreak{}box\_\allowbreak{}of\_\allowbreak{}volume\_\allowbreak{}bound}
& For any \texttt{n} and real \texttt{C}, there exists a natural radius
\texttt{R}; every lattice polytope with origin in its interior and
\texttt{(\allowbreak{}volume P)\allowbreak{}.\allowbreak{}toReal ≤ C}
has a \texttt{GLEquiv} image in
\texttt{[\allowbreak{}-\allowbreak{}R,\allowbreak{}R]\allowbreak{}\textasciicircum{}n}.
This is general conditional box reduction, including dimension zero, but
its radius is existential. \\
\texttt{Reflexiv\allowbreak{}eVolume.\allowbreak{}lean:\allowbreak{}112},
\texttt{Reflexiv\allowbreak{}eVolume.\allowbreak{}volume\_\allowbreak{}bound};
\texttt{Finitene\allowbreak{}ss.\allowbreak{}lean},
\texttt{hensley\_\allowbreak{}volume\_\allowbreak{}bound} &
\texttt{∃ C :\allowbreak{} R,\allowbreak{} ∀ P,\allowbreak{} IsReflex\allowbreak{}ive P → (\allowbreak{}volume P)\allowbreak{}.\allowbreak{}toReal ≤ C}.
The name does \textbf{not} mean the general fixed-\texttt{k} Hensley
theorem is formalized, or that its numerical estimate is supplied. \\
\texttt{Finitene\allowbreak{}ss.\allowbreak{}lean},
\texttt{lagarias\allowbreak{}\_\allowbreak{}ziegler\_\allowbreak{}box},
\texttt{latticeP\allowbreak{}oint\_\allowbreak{}count\_\allowbreak{}bound}
& Existential dimension-dependent centered radius and total
lattice-point bound, specialized to reflexive polytopes. \\
\texttt{Finitene\allowbreak{}ss.\allowbreak{}lean},
\texttt{finitely\allowbreak{}\_\allowbreak{}many\_\allowbreak{}reflexiv\allowbreak{}e}
&
\texttt{∃ T,\allowbreak{} T.\allowbreak{}Finite ∧ ∀ P,\allowbreak{} IsReflex\allowbreak{}ive P → ∃ Q ∈ T,\allowbreak{} GLEquiv P Q}.
The proof takes \textbf{all lattice polytopes in a box} as \texttt{T}.
The type neither requires every \texttt{Q ∈ T} to be reflexive nor
promises one representative per class. \\
\texttt{Finitene\allowbreak{}ss.\allowbreak{}lean},
\texttt{ncard\_\allowbreak{}lt\_\allowbreak{}of\_\allowbreak{}ssubset} &
Strict inclusion of two lattice polytopes strictly increases their total
lattice-point \texttt{ncard}; no reflexivity needed for this lemma. \\
\texttt{MaxMin.\allowbreak{}lean},
\texttt{IsMaxima\allowbreak{}lReflexi\allowbreak{}ve},
\texttt{exists\_\allowbreak{}maximal\_\allowbreak{}superset} &
Maximality is under literal inclusion on the same ambient lattice. For
reflexive \texttt{P}, obtains
\texttt{∃ Δ,\allowbreak{} P ⊆ Δ ∧ IsMaxima\allowbreak{}lReflexi\allowbreak{}ve Δ}.
Its proof maximizes the bounded point count. This is not an executable
enumeration of maximal representatives. \\
\texttt{MaxMin.\allowbreak{}lean},
\texttt{IsMinima\allowbreak{}lReflexi\allowbreak{}ve},
\texttt{maximal\_\allowbreak{}iff\_\allowbreak{}dual\_\allowbreak{}minimal},
\texttt{exists\_\allowbreak{}minimal\_\allowbreak{}subset} &
Inclusion-minimal \textbf{reflexive} objects and their polarity
correspondence. These types do not extend this correspondence to general
nonreflexive one-point objects. \\
\texttt{MaxMin.\allowbreak{}lean}, \texttt{IsMinima\allowbreak{}lIPGen},
\texttt{exists\_\allowbreak{}minimalI\allowbreak{}PGen\_\allowbreak{}subset}
& Minimality of a finite \textbf{generating set} under deletion while
preserving origin interior. This differs from inclusion-minimal
reflexivity. \\
\end{longtable}

The conditional box proof uses a maximum-determinant generator tuple,
bounds its determinant by volume via an anchored scaled parallelepiped,
proves that only finitely many row lattices have bounded nonzero
determinant, and chooses one representative matrix per row lattice.
\texttt{bounded\_\allowbreak{}matrix\_\allowbreak{}reductio\allowbreak{}n}
uses classical finite choices. Thus supplying an explicit input
\texttt{C} does not turn its conclusion into the LZ numerical radius or
into code computing \texttt{R}. Compactness ensures the volume is finite
before the conversions from extended nonnegative reals to
\texttt{toReal} are used.

The frozen reflexive volume proof uses a different route from Hensley's
general theorem:

\begin{enumerate}
\def\labelenumi{\arabic{enumi}.}
\tightlist
\item
  \texttt{WeightSy\allowbreak{}stemFini\allowbreak{}te.\allowbreak{}normaliz\allowbreak{}ed\_\allowbreak{}ip\_\allowbreak{}weight\_\allowbreak{}systems\_\allowbreak{}finite}
  proves finiteness of normalized IP weights from finite positive
  branching and \texttt{algorith\allowbreak{}m\_\allowbreak{}complete}
  in \texttt{Algorith\allowbreak{}m.\allowbreak{}lean}.
  \texttt{positive\allowbreak{}Selectio\allowbreak{}n} is noncomputable.
  Its import chain starts with
  \texttt{Algorith\allowbreak{}m → WeightSy\allowbreak{}stem → Basic},
  rather than with \texttt{Finitene\allowbreak{}ss} or a volume bound.
\item
  \texttt{Reflexiv\allowbreak{}eVolume.\allowbreak{}uniform\_\allowbreak{}weight\_\allowbreak{}lower\_\allowbreak{}bound}
  takes a common positive lower bound for all coefficients of those
  finite weight sets with supports of size at most \texttt{n+1}.
\item
  \texttt{dual\_\allowbreak{}point\_\allowbreak{}ip\_\allowbreak{}circuit}
  obtains a positive simplex circuit through each dual generator.
  \texttt{reflexiv\allowbreak{}e\_\allowbreak{}simplex\_\allowbreak{}block\_\allowbreak{}ip}
  in \texttt{Geometri\allowbreak{}cCWS.\allowbreak{}lean} uses the
  integral primal/dual pairings and a surjective open real-linear map to
  prove \textbf{integer Newton-hull IP} in
  \texttt{wsHyperp\allowbreak{}lane q}. Positivity of simplex
  coefficients alone would only establish simplex relative interior. The
  map's forward integral pairing property is not a proof that it is a
  lattice isomorphism onto every hyperplane lattice.
\item
  \texttt{dual\_\allowbreak{}generato\allowbreak{}r\_\allowbreak{}pairing\_\allowbreak{}bound}
  bounds pairings above by \texttt{ε⁻¹}; polar inequalities bound them
  below by \texttt{-\allowbreak{}1}. Independent integral dual normals
  send \texttt{P} into a fixed box. Their determinant has absolute value
  at least one, which suffices to bound the \textbf{volume} of
  \texttt{P}; this map need not be unimodular.
\item
  The separate conditional box theorem then supplies actual GL
  representatives; box counting gives finiteness and terminates
  ascending reflexive chains.
\end{enumerate}

The inspected dependency route avoids using reflexive-polytope
finiteness to prove the finite weights from which its volume bound is
derived. One must preserve that direction: replacing weight finiteness
by an argument that assumes the same polytope finiteness would create a
circular justification. The compiler-derived dependency and axiom checks
preserve that direction.

\section{Role in the concrete
classification}\label{finiteness-and-box-reduction--role-in-the-concrete-classification}

\begin{longtable}[]{@{}
  >{\raggedright\arraybackslash}p{(\linewidth - 2\tabcolsep) * \real{0.1696}}
  >{\raggedright\arraybackslash}p{(\linewidth - 2\tabcolsep) * \real{0.8304}}@{}}
\toprule\noalign{}
\begin{minipage}[b]{\linewidth}\raggedright
Result
\end{minipage} & \begin{minipage}[b]{\linewidth}\raggedright
Scope and connection
\end{minipage} \\
\midrule\noalign{}
\endhead
\bottomrule\noalign{}
\endlastfoot
Fixed \texttt{(\allowbreak{}n,\allowbreak{}k≥1)}
volume/\allowbreak{}count bounds & Hensley and LZ treat the general
positive-interior-count problem. The supporting Lean volume theorem is
reflexive-specific; the broader theorem is not needed for the
classifier's target. \\
Equivalence and marking & LZ uses affine lattice maps. The formal
classifier fixes the origin and proves two-sided integral
\texttt{GLEquiv}; these conventions must not be conflated for general
\texttt{k>1}. \\
Effective boxes & The supporting box radii are existential. Executable
\texttt{KS/\allowbreak{}IntegerP\allowbreak{}oints} derives its own
bounds from CWS weights and a certified kernel basis, not an unevaluated
LZ radius. \\
Maximal containers & Abstract maximality existence does not construct a
KS container. \texttt{KS/\allowbreak{}Containe\allowbreak{}rs} supplies
explicit seed-pairing witnesses; its containers need not be maximal. \\
Subpolytope and lattice recovery & \texttt{KS/\allowbreak{}Recovery}
proves exact hull recovery from enumerated container points;
\texttt{KS/\allowbreak{}Realizat\allowbreak{}ions} exhausts compatible
lattices. \\
Polarity for broader one-point targets & A nonreflexive one-point
polytope may have a rational polar. Such final targets are outside the
theorem. \\
WS/\allowbreak{}CWS coverage &
\texttt{KS/\allowbreak{}WeightSe\allowbreak{}arch},
\texttt{KS/\allowbreak{}Assembly} and the composition theorem supply the
concrete terminating reflexive enumeration. Native pruning and
descriptor refinement are separate guarantees. \\
Unique representatives & Output duplicates are permitted. Exact normal
forms and collision-safe class merging are not implied by volume or
point-count bounds. \\
\end{longtable}

\section{Source evidence and reusable
provenance}\label{finiteness-and-box-reduction--source-evidence-and-reusable-provenance}

Both original publisher PDFs were downloaded, read through their
relevant proof sections, and retained without byte changes in the
ignored cache. Mathematical formulas were checked on embedded scanned
pages after observing OCR errors. Each directory contains
\texttt{provenan\allowbreak{}ce.\allowbreak{}json} with primary URL,
edition, retrieval date \textbf{2026-10-14}, copyright observations and
SHA-256. Derived \texttt{paper.\allowbreak{}txt} files are reading aids,
not original source editions. No author TeX source was linked by the
consulted publisher records; none is claimed to have been cached.

\begin{longtable}[]{@{}
  >{\raggedright\arraybackslash}p{(\linewidth - 4\tabcolsep) * \real{0.3429}}
  >{\raggedright\arraybackslash}p{(\linewidth - 4\tabcolsep) * \real{0.2800}}
  >{\raggedright\arraybackslash}p{(\linewidth - 4\tabcolsep) * \real{0.3771}}@{}}
\toprule\noalign{}
\begin{minipage}[b]{\linewidth}\raggedright
Source/\allowbreak{}version
\end{minipage} & \begin{minipage}[b]{\linewidth}\raggedright
Cache path relative to \texttt{docs/\allowbreak{}learning}
\end{minipage} & \begin{minipage}[b]{\linewidth}\raggedright
Original PDF SHA-256
\end{minipage} \\
\midrule\noalign{}
\endhead
\bottomrule\noalign{}
\endlastfoot
Hensley 1983, published journal scan &
\texttt{sources/\allowbreak{}hensley-\allowbreak{}1983/\allowbreak{}journal/\allowbreak{}paper.\allowbreak{}pdf}
&
\texttt{db6d8851\allowbreak{}05160df9\allowbreak{}c53be90f\allowbreak{}4e5f468d\allowbreak{}25a36319\allowbreak{}dccef69d\allowbreak{}c5dfbfe5\allowbreak{}813e271c} \\
Lagarias--Ziegler 1991, published journal scan, publisher PDF &
\texttt{sources/\allowbreak{}lagarias\allowbreak{}-\allowbreak{}ziegler-\allowbreak{}1991/\allowbreak{}journal/\allowbreak{}paper.\allowbreak{}pdf}
&
\texttt{5d3b5a7c\allowbreak{}0bc2616e\allowbreak{}b08b176c\allowbreak{}569c46cb\allowbreak{}b868f6c7\allowbreak{}e9e29f91\allowbreak{}85038fe7\allowbreak{}5698ac63} \\
\end{longtable}

The LZ PDF download URL is the publisher's
\href{https://www.cambridge.org/core/services/aop-cambridge-core/content/view/0D4E5B45AB64DC6069A9BD7DD45B450C/S0008414X00007483a.pdf/bounds-for-lattice-polytopes-containing-a-fixed-number-of-interior-points-in-a-sublattice.pdf}{journal
PDF}. The cached Hensley publisher HTML record additionally verifies the
DOI; its original-byte hash is recorded in the local provenance file.
Cached literature is not an input to the Lean build.

The main inspected Lean file hashes identify the source evidence
independently of declaration names:

\begin{longtable}[]{@{}
  >{\raggedright\arraybackslash}p{(\linewidth - 2\tabcolsep) * \real{0.2903}}
  >{\raggedright\arraybackslash}p{(\linewidth - 2\tabcolsep) * \real{0.7097}}@{}}
\toprule\noalign{}
\begin{minipage}[b]{\linewidth}\raggedright
Snapshot file in \texttt{Refpoly/}
\end{minipage} & \begin{minipage}[b]{\linewidth}\raggedright
SHA-256
\end{minipage} \\
\midrule\noalign{}
\endhead
\bottomrule\noalign{}
\endlastfoot
\texttt{Basic.\allowbreak{}lean} &
\texttt{a716a150\allowbreak{}8842b4f4\allowbreak{}ed1b24ea\allowbreak{}f978d437\allowbreak{}58efbdb8\allowbreak{}7990b4a6\allowbreak{}4e80c879\allowbreak{}5089dfc4} \\
\texttt{BoxReduc\allowbreak{}tion.\allowbreak{}lean} &
\texttt{85e9dbae\allowbreak{}f06e9702\allowbreak{}267cd8cf\allowbreak{}5f88b1b7\allowbreak{}339f0261\allowbreak{}77f1cfde\allowbreak{}359e9329\allowbreak{}1dacebfd} \\
\texttt{Reflexiv\allowbreak{}eVolume.\allowbreak{}lean} &
\texttt{23d54cf6\allowbreak{}206a41b1\allowbreak{}0bfaf206\allowbreak{}fde3e67b\allowbreak{}f532fb2c\allowbreak{}06207f78\allowbreak{}5e3da312\allowbreak{}4781c589} \\
\texttt{Finitene\allowbreak{}ss.\allowbreak{}lean} &
\texttt{0152ac35\allowbreak{}97310ff9\allowbreak{}783b3066\allowbreak{}5faf7fc4\allowbreak{}9c4ddf3e\allowbreak{}3ad8af5b\allowbreak{}c83b9f6c\allowbreak{}02a222af} \\
\texttt{MaxMin.\allowbreak{}lean} &
\texttt{cf4bf17f\allowbreak{}61a68356\allowbreak{}43710932\allowbreak{}26ea904d\allowbreak{}a42a5be6\allowbreak{}f9d0cf07\allowbreak{}183e0e82\allowbreak{}6a74b6ef} \\
\end{longtable}

\chapter{Normal forms, exact equality, and classification
output}\label{normal-forms-and-classification--normal-forms-exact-equality-and-classification-output}

\noindent\textit{Markdown source: }\texttt{docs/\allowbreak{}learning\allowbreak{}/\allowbreak{}normal-\allowbreak{}forms-\allowbreak{}and-\allowbreak{}classifi\allowbreak{}cation.\allowbreak{}md}

A normal form can decide lattice equivalence only after the target
lattice and equivalence relation are fixed. Its canonical integer
matrix, serialized bytes and a hash of those bytes are three different
objects. This note reviews the primary normal-form algorithm and
pinned/archived computational contracts. The
\href{../../proofs/STATEMENTS.md}{Lean classifier} proves coverage with
possible duplicates; native normal forms, exact deduplication and census
counts are separate guarantees.

\section{Primary references and reusable
evidence}\label{normal-forms-and-classification--primary-references-and-reusable-evidence}

\textbf{{[}GK13{]}} Roland Grinis and Alexander M. Kasprzyk,
\emph{Normal forms of convex lattice polytopes},
\href{https://arxiv.org/abs/1301.6641v1}{arXiv:\allowbreak{}1301.\allowbreak{}6641v1},
submitted 28 January 2013, 28 pages. The primary version inspected is
the \href{https://arxiv.org/pdf/1301.6641v1}{PDF} and author
\href{https://arxiv.org/src/1301.6641v1}{TeX archive}. No journal
publication metadata was present in the checked arXiv record, and none
is inferred. Relevant statements are §3, pages 8--12; Definitions
3.1--3.2, pages 8--9; affine normal form in §3.3, page 11; PALP's extra
column ordering in §3.4, pages 11--12; Appendix A, pages 16--25,
especially A.8 on pages 24--25. Section 2 treats a different, more
general face-graph approach, including rational and lower-dimensional
polytopes.

\textbf{{[}PALP04{]}} Maximilian Kreuzer and Harald Skarke, \emph{PALP:
A Package for Analyzing Lattice Polytopes with Applications to Toric
Geometry}, Computer Physics Communications \textbf{157} (2004), 87--106,
\href{https://doi.org/10.1016/S0010-4655\%2803\%2900491-0}{DOI
10.\allowbreak{}1016/\allowbreak{}S0010-\allowbreak{}4655(\allowbreak{}03)\allowbreak{}00491-\allowbreak{}0}.
Title/authors/publication metadata and DOI were verified against
\href{https://arxiv.org/abs/math/0204356v1}{arXiv:\allowbreak{}math/\allowbreak{}0204356v\allowbreak{}1},
submitted 30 April 2002. The actual paper read is the primary
\href{https://arxiv.org/pdf/math/0204356v1}{22-page v1 PDF} and
\href{https://arxiv.org/src/math/0204356v1}{author source}. Section 1,
pages 1--2, states full-dimension and fixed-precision limitations; §2's
\texttt{-\allowbreak{}N}/\texttt{-\allowbreak{}t} documentation, page 6,
describes the isomorphism test and stages. The preprint year differs
from the journal year.

The consulted editions have the following original-file hashes;
extracted TeX/text are reading aids. Retrieval date is 2026-10-14 UTC.
Paths are relative to
\texttt{docs/\allowbreak{}learning\allowbreak{}/\allowbreak{}sources/}:

\begin{longtable}[]{@{}
  >{\raggedright\arraybackslash}p{(\linewidth - 4\tabcolsep) * \real{0.1902}}
  >{\raggedright\arraybackslash}p{(\linewidth - 4\tabcolsep) * \real{0.4049}}
  >{\raggedright\arraybackslash}p{(\linewidth - 4\tabcolsep) * \real{0.4049}}@{}}
\toprule\noalign{}
\begin{minipage}[b]{\linewidth}\raggedright
Source/\allowbreak{}version
\end{minipage} & \begin{minipage}[b]{\linewidth}\raggedright
Original PDF SHA-256
\end{minipage} & \begin{minipage}[b]{\linewidth}\raggedright
Original source SHA-256
\end{minipage} \\
\midrule\noalign{}
\endhead
\bottomrule\noalign{}
\endlastfoot
\texttt{grinis-\allowbreak{}kasprzyk\allowbreak{}-\allowbreak{}1301-\allowbreak{}6641/\allowbreak{}v1/}
&
\texttt{0e2fb84a\allowbreak{}6bd81009\allowbreak{}924079ca\allowbreak{}02a39dec\allowbreak{}a2f87ff5\allowbreak{}75d942cb\allowbreak{}b05def37\allowbreak{}b045273d}
&
\texttt{fde1cb48\allowbreak{}7f13d1d1\allowbreak{}2d0ce992\allowbreak{}f71c9001\allowbreak{}2e9ceafc\allowbreak{}b2e61bed\allowbreak{}b292c123\allowbreak{}85e4526a} \\
\texttt{palp-\allowbreak{}math-\allowbreak{}0204356/\allowbreak{}v1/} &
\texttt{5134b66b\allowbreak{}aac3b000\allowbreak{}9b4f4cdd\allowbreak{}02aa642d\allowbreak{}958d7cef\allowbreak{}03f66f7b\allowbreak{}cccf2c12\allowbreak{}7c30e247}
&
\texttt{627ca351\allowbreak{}5b570be5\allowbreak{}560bd939\allowbreak{}dc4db3ff\allowbreak{}69466995\allowbreak{}fd65beff\allowbreak{}fb6a619f\allowbreak{}61c4bf51} \\
\end{longtable}

Each directory contains \texttt{paper.\allowbreak{}pdf},
\texttt{source.\allowbreak{}tar.\allowbreak{}gz},
\texttt{provenan\allowbreak{}ce.\allowbreak{}json},
\texttt{paper.\allowbreak{}txt}, and supplementary \texttt{tex/} files.
The PALP source download is a gzip-compressed \textbf{single TeX file},
despite the local archive filename; provenance records that format.
Source-license status is unverified; these are local reading caches.
Printed page numbers, rather than guessed journal pages, are used above.

Implementation evidence is fixed to PALP commit
\texttt{968bb4b3\allowbreak{}c6c42a08\allowbreak{}93c95340\allowbreak{}e8171854\allowbreak{}8e13f12a},
root archived \texttt{main} commit
\texttt{7c4dc2e6\allowbreak{}2da1ff8d\allowbreak{}539684d4\allowbreak{}977d4de9\allowbreak{}656d7a4f},
and \texttt{gpu-\allowbreak{}run-\allowbreak{}43types} commit
\texttt{452fbe8e\allowbreak{}1b296bf7\allowbreak{}5fa9e8be\allowbreak{}cf6c34d8\allowbreak{}d48bf00a}.
Current regression contracts are documented in
\texttt{tests/\allowbreak{}README.\allowbreak{}md},
\texttt{tests/\allowbreak{}cws\_\allowbreak{}support.\allowbreak{}py},
and \texttt{tests/\allowbreak{}test\_\allowbreak{}cws.\allowbreak{}py}.
Formal equivalence is defined in
\texttt{proofs/\allowbreak{}Refpoly/\allowbreak{}Basic.\allowbreak{}lean};
\href{../../proofs/SOURCES.md}{SOURCES.md} identifies its original
edition.

\section{Linear and affine equivalence are different
contracts}\label{normal-forms-and-classification--linear-and-affine-equivalence-are-different-contracts}

For a fixed rank-\texttt{n} lattice, linear lattice isomorphism means
\texttt{Q=\allowbreak{}g(\allowbreak{}P)} for \texttt{g} in
\texttt{GL(\allowbreak{}n,\allowbreak{}Z)}. The origin is distinguished
and fixed. Affine lattice equivalence allows
\texttt{Q=\allowbreak{}g(\allowbreak{}P)\allowbreak{}+t} for an integral
translation \texttt{t}. Neither is merely combinatorial equivalence or
an arbitrary invertible real transformation.

A lattice polytope with exactly one interior lattice point has an
intrinsic centering rule: translate that unique point to zero. Every
affine lattice equivalence sends that point to the unique point of the
image; after centering both polytopes, it becomes linear. This applies
to reflexive and general one-point lattice polytopes. If several
interior lattice points exist, choosing one arbitrarily can distinguish
affine-equivalent but differently marked objects. For instance,
intervals \texttt{[\allowbreak{}-\allowbreak{}1,\allowbreak{}3]} and
\texttt{[\allowbreak{}-\allowbreak{}2,\allowbreak{}2]} are integral
translates, both have origin interior, and are not related by
\texttt{GL(\allowbreak{}1,\allowbreak{}Z)\allowbreak{}=\allowbreak{}\{+1,\allowbreak{}-\allowbreak{}1\}}.
A marked origin or a translation-invariant normal form must be
specified.

{[}PALP04{]} describes ordinary normal form in the linear setting.
{[}GK13{]} §3.3 separately defines an affine normal form by minimizing
over vertex translations as well as the relevant vertex permutations.
Its finite vertex-translation formula covers affine lattice equivalence
of full-dimensional lattice polytopes. Neither origin interior nor
reflexivity is required by the normal-form construction in §3. The
full-dimensional \textbf{lattice} hypothesis is essential to its integer
Hermite-normal-form step. The paper separately discusses
lower-dimensional reductions in §3.2 and general rational isomorphism
testing in §2; those extensions should not be silently assumed to be the
pinned \texttt{poly -\allowbreak{}N} command's general input contract.

\section{The primary normal-form
algorithm}\label{normal-forms-and-classification--the-primary-normal-form-algorithm}

With vertices arranged as columns of an integer matrix \texttt{V},
changing the lattice basis acts on the left by a unimodular matrix;
vertex relabeling acts on the columns. Hermite normal form removes the
first freedom for a fixed column order, but does not remove column-order
dependence. Therefore choosing a single arbitrary vertex order and
taking HNF is not sufficient.

{[}GK13{]} Definition 3.1 forms the vertex--facet \textbf{height} matrix

\texttt{PM\_\allowbreak{}ij =\allowbreak{} <w\_\allowbreak{}i,\allowbreak{}v\_\allowbreak{}j> + c\_\allowbreak{}i},

where \texttt{w\_\allowbreak{}i} is a primitive inward integral facet
normal and its supporting facet has equation
\texttt{<w\_\allowbreak{}i,\allowbreak{}x>=\allowbreak{}-\allowbreak{}c\_\allowbreak{}i}.
This matrix is invariant under lattice automorphisms and translations,
up to row and column permutations. For a nonreflexive lattice polytope
the \texttt{c\_\allowbreak{}i} need not equal one; no reflexive polar is
needed. Its purpose is to organize candidate vertex orders, not to serve
by itself as a complete lattice-equivalence invariant.

The construction takes the lexicographically \textbf{maximal} height
matrix over facet/vertex permutations, retains the permutations
attaining it (equivalently one maximizing permutation composed with all
its matrix automorphisms), and then minimizes the HNF of the
corresponding reordered vertex matrices. Definition 3.2 gives this
ordinary normal form. PALP adds a deterministic column order based on
each maximal height-matrix column's maximum entry, sum, and relative
position, as explained in §3.4 and Appendix A.8. The final choice is the
lexicographically \textbf{minimal} vertex HNF over the remaining
candidate permutations. Matrix symmetries can exceed actual lattice
automorphisms; retaining all the relevant candidates before HNF is part
of the correctness argument.

Appendix A, page 16, explicitly states the canonical-output contract:
for full-dimensional lattice polytopes in the specified lattice, equal
PALP normal forms are equivalent to linear lattice isomorphism. This is
an algorithmic description and correctness statement in the primary
source, not a numbered Lean theorem or a proof that every future
reimplementation matches it. Finite permutation sets and integer HNF
give a terminating mathematical procedure; arithmetic range, capacities,
and complete permutation handling in a C build remain implementation
obligations. Definitions 3.1 and 3.2 are the relevant numbered items;
this note does not invent a separate theorem number for that contract.

Determinism is strictly weaker. A routine could always return the same
wrong matrix for the same input and be deterministic without being
canonical. Canonicity needs both invariance under every allowed
equivalence and separation of inequivalent inputs. Even a canonical
matrix is not a collision-free finite-width digest.

\section{Pinned PALP and the actual regression
target}\label{normal-forms-and-classification--pinned-palp-and-the-actual-regression-target}

At the pinned PALP commit, \texttt{Polynf.\allowbreak{}c} implements the
expected stages: \texttt{Init\_\allowbreak{}rVM\_\allowbreak{}VPM}
computes integer vertex coordinates and facet evaluations;
\texttt{Make\_\allowbreak{}VPM\_\allowbreak{}NF} organizes the maximal
pairing matrix and its admissible permutations;
\texttt{New\_\allowbreak{}pNF\_\allowbreak{}Order} supplies the
additional column ordering;
\texttt{Aux\_\allowbreak{}pNF\_\allowbreak{}from\_\allowbreak{}vNF}/\texttt{Aux\_\allowbreak{}Make\_\allowbreak{}Triang}
compare triangularized vertex matrices;
\texttt{GLZ\_\allowbreak{}Make\_\allowbreak{}Trian\_\allowbreak{}NF}
supplies the unimodular triangular/HNF step.
\texttt{Aux\_\allowbreak{}XltY\_\allowbreak{}Poly\_\allowbreak{}NF}
compares entries row by row.
\texttt{Init\_\allowbreak{}rVM\_\allowbreak{}VPM} records whether facet
constants all equal one, but
\texttt{Make\_\allowbreak{}Poly\_\allowbreak{}NF} computes the normal
form regardless and returns that reflexivity flag separately. Therefore
the normal-form routine is not restricted mathematically to reflexive
objects.

\texttt{poly.\allowbreak{}c} invokes
\texttt{Make\_\allowbreak{}Poly\_\allowbreak{}Sym\_\allowbreak{}NF} on
\texttt{\_\allowbreak{}P} for option \texttt{N}, and
\texttt{Polynf.\allowbreak{}c} labels the result \textbf{Normal form of
vertices of P}. Option \texttt{A} calls the separate
\texttt{Make\_\allowbreak{}ANF} path; it temporarily adds an auxiliary
coordinate and origin, obtains a linear normal form, and
reduces/restores the affine representation. The root regression command
is \texttt{poly -\allowbreak{}fN}, with no dual-selection flag. Its
goldens are vertex normal forms of \textbf{P}, not of its polar.

\texttt{tests/\allowbreak{}cws\_\allowbreak{}support.\allowbreak{}py}
parses each normal form as
\texttt{(\allowbreak{}dimensio\allowbreak{}n,\allowbreak{} vertex\_\allowbreak{}count,\allowbreak{} tuple(\allowbreak{}matrix\_\allowbreak{}rows)\allowbreak{})}.
It requires the expected P header, exactly
\texttt{dimensio\allowbreak{}n} rows, and exactly
\texttt{vertex\_\allowbreak{}count} integers per row.
\texttt{serializ\allowbreak{}e\_\allowbreak{}forms} writes the shape and
integer rows as normalized text. Only whitespace and the informational
\texttt{perm=} annotation are removed.
\texttt{tests/\allowbreak{}test\_\allowbreak{}cws.\allowbreak{}py}
compares every resulting matrix to its corresponding golden in sequence,
requires one matrix per accepted CWS row, and separately checks
\texttt{len(\allowbreak{}set(\allowbreak{}forms)\allowbreak{})} for the
distinct count. Repeated NFs remain in the accepted-row list: their
multiplicities are part of the regression contract.

These are exact PALP consistency regressions for the checked-in pools.
They do not independently prove mathematical canonicity, all-dimensional
completeness, equivalence with the abstract catalogue, or the
correctness of another lattice realization. A changed golden requires
the repository's explicit baseline-review procedure; regenerating it
merely to pass a test would erase evidence.

\section{What the archived classifiers actually
compare}\label{normal-forms-and-classification--what-the-archived-classifiers-actually-compare}

Both inspected archived classifiers,
\texttt{main:\allowbreak{}src/\allowbreak{}classify\allowbreak{}/\allowbreak{}classifi\allowbreak{}er.\allowbreak{}cpp}
and
\texttt{gpu-\allowbreak{}run-\allowbreak{}43types:\allowbreak{}src/\allowbreak{}classify\allowbreak{}/\allowbreak{}classifi\allowbreak{}er.\allowbreak{}cpp},
compute a 128-bit xxHash of the used NF entries, stored contiguously as
native \texttt{Long} values. Their maps use \texttt{Hash128} as the key.
On a key hit, counts are merged without comparing the full integer NF
matrices. The payload stores counts and representative metadata, not the
full matrix used for an exact equality check. The same digest-key logic
appears in local-to-global merges. This is hash-based aggregation; no
collision-resolution equality path was present in those inspected
sections.

The archived
\texttt{gpu-\allowbreak{}run-\allowbreak{}43types:\allowbreak{}src/\allowbreak{}classify\allowbreak{}/\allowbreak{}compare\_\allowbreak{}polytope\allowbreak{}\_\allowbreak{}outputs.\allowbreak{}cpp}
compares sorted
\texttt{(\allowbreak{}hash\_\allowbreak{}lo,\allowbreak{}hash\_\allowbreak{}hi,\allowbreak{}count)}
records. Passing it shows equality of those records, not exact equality
of the underlying NF matrices.
\texttt{src/\allowbreak{}verify/\allowbreak{}proofs/\allowbreak{}nf\_\allowbreak{}determin\allowbreak{}ism.\allowbreak{}md}
discusses same-input deterministic sorting and hashing, with assertions
about other stages; it is an archived argument, not a formal proof of NF
canonicity or hash injectivity. A total-order comparator addresses one
sorting concern, while leaving the geometric canonicalization and
arithmetic contracts distinct.

The archived \texttt{palp\_\allowbreak{}api.\allowbreak{}h} wrappers
construct P from CWS data, find its vertices/equations, sort its vertex
indices, and call
\texttt{Make\_\allowbreak{}Poly\_\allowbreak{}Sym\_\allowbreak{}NF} on
those points. \texttt{result.\allowbreak{}ok} is set after the IP
pipeline succeeds; it is not, by that assignment alone, an independently
asserted reflexivity flag. Again the returned matrix is primal P's NF.

The newer archived \texttt{Polytope\allowbreak{}Info} preserves
first-CWS replay fields: schema version, structure/profile IDs, row
count, coordinate count, source index, degrees, all padded weights, and
a legacy six-weight projection, alongside geometric/Hodge metadata.
These fields can help reconstruct a representative after aggregation,
provided their input/lattice/build interpretation is pinned. They do not
retain all members of a digest bucket or automatically detect an
incorrect earlier merge.

\texttt{gpu-\allowbreak{}run-\allowbreak{}43types:\allowbreak{}src/\allowbreak{}classify\allowbreak{}/\allowbreak{}add\_\allowbreak{}nf.\allowbreak{}cpp}
replays those first-CWS inputs and appends
\texttt{nf\_\allowbreak{}vertices} as row-major native \texttt{int32}
bytes. Its optional hash verification checks agreement with the stored
digest, not equality against a separately preserved full canonical
matrix. The conversion's range check is guarded by
\texttt{first\_\allowbreak{}block}, so the inspected check does not
cover every later successful block per worker; the source's general
expectation that values fit is not a proved universal bound. A lossless
output contract must justify all converted entries or check them all.
This review records the historical behavior without modifying it or
inferring that any actual retained entry overflowed.

\section{Exact reduction and serialization
obligations}\label{normal-forms-and-classification--exact-reduction-and-serialization-obligations}

For a mathematically exact class merge, the decisive object should
include the declared equivalence/lattice convention, dimension, vertex
count, and every integer NF entry. A digest may select a partition or
candidate bucket. A hit must then compare complete canonical values, or
replay enough immutable data to recover and compare them exactly.
Distinct matrices in the same digest bucket must remain distinct
classes. The counts of equivalent rows should be accumulated only after
that equality is established.

A persistent binary encoding needs a defined signed-integer width, byte
order, matrix layout, shape, and algorithm/schema version. The native
\texttt{Long} hash buffer and native \texttt{int32} output are different
encodings; equality of hashes computed from one does not prove lossless
serialization into the other. Fixed platform/build assumptions may
support replay within a pinned run, but should be recorded instead of
being mistaken for a portable mathematical normal form. Shape and
integer comparisons in the current text goldens avoid relying on native
byte width or hash collision assumptions.

Exact replay also needs immutable CWS input, the relevant
descriptor/support ordering, PALP commit/build capacities/arithmetic
choices, and the lattice realization. The same real polytope on a
different lattice can represent a different equivalence class, so a
normal form only classifies the lattice coordinates it is actually
given. Recording weights alone may be insufficient for finite-index
realizations beyond the reference lattice.

\section{Formal support and remaining
gaps}\label{normal-forms-and-classification--formal-support-and-remaining-gaps}

\texttt{Basic.\allowbreak{}lean} defines
\texttt{IsUnimod\allowbreak{}ular g} by \textbf{both} forward and
inverse preservation of integer lattice points. \texttt{GLEquiv P Q} is
the existence of a real linear equivalence with that unimodularity
property carrying P onto Q. Its declarations prove reflexivity,
symmetry, transitivity, lattice-image/cardinality preservation, and
preservation of lattice-polytope structure. They give an abstract
equivalence relation; they do not define NF, its serializer, a hash, or
a native reduction theorem.

\begin{longtable}[]{@{}
  >{\raggedright\arraybackslash}p{(\linewidth - 4\tabcolsep) * \real{0.1831}}
  >{\raggedright\arraybackslash}p{(\linewidth - 4\tabcolsep) * \real{0.3005}}
  >{\raggedright\arraybackslash}p{(\linewidth - 4\tabcolsep) * \real{0.5164}}@{}}
\toprule\noalign{}
\begin{minipage}[b]{\linewidth}\raggedright
Contract
\end{minipage} & \begin{minipage}[b]{\linewidth}\raggedright
Evidence available
\end{minipage} & \begin{minipage}[b]{\linewidth}\raggedright
Missing or bounded scope
\end{minipage} \\
\midrule\noalign{}
\endhead
\bottomrule\noalign{}
\endlastfoot
Mathematical linear lattice equivalence &
\texttt{IsUnimod\allowbreak{}ular}, \texttt{GLEquiv}; primary NF
contract & No formal theorem
\texttt{NF(\allowbreak{}P)\allowbreak{}=\allowbreak{}NF(\allowbreak{}Q)\allowbreak{} iff GLEquiv P Q}. \\
Affine equivalence & GK §3.3 algorithm; unique-interior centering
deduction & No inspected formal
affine-\allowbreak{}equivale\allowbreak{}nce/\allowbreak{}NF bridge;
multiple origins need a declared convention. \\
P versus polar target & Pinned PALP call/\allowbreak{}output and exact
root parser & Switching targets requires an explicit contract and
fixture change. \\
General nonreflexive NF & GK full-dimensional lattice scope; PALP height
matrix &
Rational\allowbreak{}/\allowbreak{}lower-\allowbreak{}dimensio\allowbreak{}nal
reductions and native capacity bounds need separate verification. \\
Same-input determinism & Archived argument and finite algorithm
structure & Does not prove equivalent-input canonicity or separation. \\
Collision-safe class merging & Required exact-NF equality policy &
Inspected archived maps compare digest keys only; no full payload
collision resolution. \\
Lossless serialization & Exact integer text goldens; archived row-major
binary path & Portable binary convention and complete
width/\allowbreak{}range guarantee are not established by those
wrappers. \\
Replay & First-CWS
schema/\allowbreak{}degree/\allowbreak{}weight/\allowbreak{}source
metadata & Representative replay alone cannot recover collided classes;
lattice/\allowbreak{}build/\allowbreak{}input provenance must be
immutable. \\
Scientific output counts & Regression distinguishes accepted rows and
distinct exact NFs & Hash-count equality is insufficient to certify an
exact geometric census. \\
End-to-end completeness & Concrete Lean KS coverage with lattices and
subpolytope recovery & Native
NF/\allowbreak{}refineme\allowbreak{}nt/\allowbreak{}dedup correctness
is not supplied by that theorem. \\
\end{longtable}

A complete enumeration theorem can permit duplicates: every output must
be valid and every target class represented, with termination.
Duplicate-free output is an additional theorem only if one
representative per equivalence class is promised. Normal forms are
useful for that stronger promise and for exact distinct-class counts;
they do not replace seed, lattice-realization or subpolytope coverage.
Native refinement and census guarantees therefore keep coverage,
correctness and exact equivalence reduction separate.

\appendix
\chapter{Markdown source snapshot}
\label{markdown-source-snapshot}
The review chapters are a presentation of the complete Markdown notes,
not replacements for them. Their SHA-256 fingerprints below identify
the exact source bytes used for this compilation. Source bibliographies
and primary-document provenance remain in their respective chapters.
\begin{description}
\item[\texttt{ip-\allowbreak{}reflexiv\allowbreak{}ity-\allowbreak{}and-\allowbreak{}pyramids\allowbreak{}.\allowbreak{}md}]
\texttt{42ac78da\allowbreak{}8865d9eb\allowbreak{}f1de4ae9\allowbreak{}3cdd444e\allowbreak{}ca9fa549\allowbreak{}4285b07d\allowbreak{}e0e71ca6\allowbreak{}b1bbb8a0}
\item[\texttt{weight-\allowbreak{}systems-\allowbreak{}and-\allowbreak{}wps.\allowbreak{}md}]
\texttt{f1fe431c\allowbreak{}1444b19d\allowbreak{}c79fc5d9\allowbreak{}63f9388a\allowbreak{}334ac7a2\allowbreak{}5028078e\allowbreak{}7be5775e\allowbreak{}1aec72a9}
\item[\texttt{minimal-\allowbreak{}ip-\allowbreak{}circuits\allowbreak{}-\allowbreak{}and-\allowbreak{}types.\allowbreak{}md}]
\texttt{9c1a076d\allowbreak{}ff8cdd0f\allowbreak{}906ba8a4\allowbreak{}94440f87\allowbreak{}bfe7a255\allowbreak{}11246db4\allowbreak{}f0b6590d\allowbreak{}8b723638}
\item[\texttt{weight-\allowbreak{}enumerat\allowbreak{}ion-\allowbreak{}and-\allowbreak{}degree-\allowbreak{}bounds.\allowbreak{}md}]
\texttt{2645d215\allowbreak{}ed905af8\allowbreak{}179b486b\allowbreak{}d543aeb8\allowbreak{}96400429\allowbreak{}a17b5f49\allowbreak{}cad6779a\allowbreak{}ff879f17}
\item[\texttt{cws-\allowbreak{}and-\allowbreak{}lattice-\allowbreak{}realizat\allowbreak{}ions.\allowbreak{}md}]
\texttt{9ee15150\allowbreak{}f165f6f4\allowbreak{}3a187b99\allowbreak{}be76cd50\allowbreak{}3eaaebd1\allowbreak{}616c4fd1\allowbreak{}9285a627\allowbreak{}3eb31bc6}
\item[\texttt{finitene\allowbreak{}ss-\allowbreak{}and-\allowbreak{}box-\allowbreak{}reductio\allowbreak{}n.\allowbreak{}md}]
\texttt{332ab831\allowbreak{}137db6d7\allowbreak{}52c508bb\allowbreak{}b6649f2a\allowbreak{}80f8e134\allowbreak{}f37b4840\allowbreak{}41ba477b\allowbreak{}5070b12d}
\item[\texttt{normal-\allowbreak{}forms-\allowbreak{}and-\allowbreak{}classifi\allowbreak{}cation.\allowbreak{}md}]
\texttt{8c41f928\allowbreak{}5e206f1a\allowbreak{}bcf0a984\allowbreak{}b3e6fdaf\allowbreak{}5ff6d90d\allowbreak{}ef629e5c\allowbreak{}f441b9b3\allowbreak{}9cc6460a}
\end{description}
% END GENERATED LITERATURE REVIEW

\end{document}
